

\documentclass[11pt]{book}

\usepackage{amsmath,amssymb,amsfonts,amsthm}
\usepackage{stmaryrd}
\SetSymbolFont{stmry}{bold}{U}{stmry}{m}{n}
\usepackage{mathtools}
\usepackage{geometry}
\geometry{a4paper,margin=1in}
\usepackage{enumitem}
\usepackage{bm}
\usepackage{graphicx}
\usepackage{hyperref} \usepackage{tikz-cd}

\usepackage{fancyhdr}
\pagestyle{fancy}
\fancyhf{} 

\fancyhead[LE,RO]{\footnotesize
  \thepage\quad CHAPTER~\thechapter
}



\theoremstyle{plain}
\newtheorem{theorem}{Theorem}[chapter]
\newtheorem{proposition}[theorem]{Proposition}
\newtheorem{lemma}[theorem]{Lemma}
\newtheorem{corollary}[theorem]{Corollary}
\newtheorem{assumption}[theorem]{Assumption}

\theoremstyle{definition}
\newtheorem{definition}[theorem]{Definition}
\newtheorem{example}[theorem]{Example}

\theoremstyle{remark}
\newtheorem{remark}[theorem]{Remark}



\newcommand{\PLPIMMTCN}{\PLPIMMTEN} \newcommand{\PLPIMMTEN}{\text{Meta-Mathematical Theory based on Pan-Logic Analysis and Pan-Iterative Analysis}}
\newcommand{\PLPIMMT}{\text{PL-PI MMT}}  

\newcommand{\GFramework}{\text{G-Framework}} \newcommand{\GAlgebra}{\text{G-Algebra}}     

\newcommand{\GZLStxt}{\textsf{GZ-LS}}   \newcommand{\GZLTtxt}{\textsf{GZ-LT}}   \newcommand{\GZLOCtxt}{\textsf{GZ-LOC}} \newcommand{\GZLCurvtxt}{\textsf{GZ-LCurv}} 

\newcommand{\GZLS}{\ensuremath{\mathfrak{L}_{\mathrm{GZ}}}}       \newcommand{\GZLT}{\ensuremath{M_{!w}}}                           \newcommand{\GZLOC}{\ensuremath{A_M}}                             \newcommand{\GZLC}{\ensuremath{\omega_{\mathrm{law}}}}            \newcommand{\GZLCurv}{\ensuremath{\mathcal{F}_{\mathrm{law}}}}    

\newcommand{\GZTLCtxt}{\textsf{GZ-TLC}}                           \newcommand{\GZTLC}{\ensuremath{\mathcal{A}}}                     \newcommand{\GZGMS}{\textsf{GZ-GMS}}                              

\newcommand{\LawSpace}{\ensuremath{\mathcal{L}}}                  \newcommand{\Base}{\ensuremath{M}}                                \newcommand{\Bundle}{\ensuremath{P}}                              \newcommand{\pr}{\operatorname{pr}}                               \newcommand{\Group}{\ensuremath{G}}                               \newcommand{\Conn}{\ensuremath{\omega}}                           \newcommand{\Curv}{\ensuremath{\Omega}}                           \newcommand{\Homotopy}{\ensuremath{\mathcal{H}}}                  \newcommand{\Jac}{\ensuremath{\mathsf{Jac}}}                      \newcommand{\Layer}{\ensuremath{\ell}}                            

\newcommand{\R}{\mathbb{R}}
\newcommand{\N}{\mathbb{N}}
\newcommand{\Z}{\mathbb{Z}}
\newcommand{\C}{\mathbb{C}}

\setlist[itemize]{leftmargin=2em}
\setlist[enumerate]{leftmargin=2em}



\title{\textbf{GaoZheng G-Framework and GaoZheng G-Algebra:\\
Applied Mathematics Volume III -- HACA, PACER, and Certificate-Based AI}}

\author{GaoZheng}
\date{Version 1.0, 2025}

\begin{document}

\raggedbottom 

\maketitle

\frontmatter



\chapter*{License and Permissions}


Copyright \textcopyright{} 2025 GaoZheng.

This arXiv version of the monograph is made available under the
Creative Commons Attribution 4.0 International license (CC BY 4.0).
In particular, any reuse or adaptation of this text must provide
appropriate credit to the author, include a reference to this
monograph, and indicate if changes were made, but is not otherwise
restricted by non-commercial or no-derivative clauses.

Earlier project-level source materials in the public repository: \\
\url{https://github.com/CTaiDeng/open_meta_mathematical_theory} \\
may be released under more restrictive licenses (such as
CC BY-NC-ND 4.0 for markdown source documents or GPL-3.0-only for
code). Those repository licenses govern public use of the
corresponding source files. The present arXiv monograph is an
author-prepared, curated derivative work of that project and is
separately licensed under CC BY 4.0 as stated above.


\chapter*{Abstract}
This Applications Volume~III continues the integrated construction of the
GaoZheng \\G-Framework and GaoZheng G-Algebra by developing a law-space–aware
semantics for hierarchical cognitive architectures (HACA) and pack-aligned
compressive–expansion reasoners (PACER). It is the third applied volume in the
series, building on the law-space geometry and GaoZheng Reinforcement Learning
(GRL) constructions of the pure mathematics volume and Applications Volumes~I–II.
The central idea is to treat semantics itself as a law-space object:
multi-scale operator universes, semantic functor fields, and GRL in law-space
path integrals are all embedded into the principal-bundle–based generalized
noncommutative Lie algebra (PFB-GNLA) background constructed in the previous
volumes.

Heuristically, the law-space layering perspective is aligned with the
HACA/PACER construction as follows. The $L_1$-layer is the raw operator
universe
\[
  \Omega(L)
  =
  \Omega^{(0)}(L)\sqcup\Omega^{(1)}(L)\sqcup\Omega^{(2)}(L)\sqcup\Omega^{(3)}(L),
\]
consisting of characters, word-packs, sentence structures, and paragraph
paths on KAT action monoids~\cite{KozenSmithKAT1996}. The $L_2$-layer resolves $\Omega(L)$ into a
semantic property layer $\mathrm{Prop}(L)$, built from basic semantic
functors and an intentional power-set construction of observable, homotopy-stable,
law-space–realizable subsets; in this sense $L_2$ addresses the \emph{problem
of meaning}. The $L_3$-layer adds homotopy and Jacobiator-gap data on top of
$L_2$, providing a framework in which robustness and safety can be formulated
as homotopy-level constraints rather than merely pointwise predicates; this
layer addresses the \emph{problem of safety}. A further $L_4$-layer is
proposed to encode metaphorical and cross-domain connections via
structure-preserving functors between law-objects and law-sectors, thus
targeting the \emph{problem of metaphor and analogy}. In principle, higher
layers $L_n$ can be introduced to capture additional forms of structure and
constraint (value geometry, ethics, multi-agent games, and so on), each
layer resolving classes of problems that cannot be fully treated on the
layers below.

Within this layered picture, the present volume develops three main
constructions. First, semantic functor fields and their noncommutative
semantic functor algebras are introduced as degenerate, semantically
constrained sectors of the underlying PFB-GNLA; each semantic functor can
be interpreted as a differential dynamical quantum of meaning (semantic MDQ)
for HACA/PACER. Second, GRL path integrals are developed at two coupled
levels: lexical GRL path integrals on lexical KAT action monoids govern
minimal-action principles for word-pack organization, while semantic GRL
path integrals on paragraph-level KAT paths yield long semantic expressions
whose admissibility is certified by semantic predicates and law-space
sectors. In particular, a paragraph operator is fully admissible only if it
satisfies both lexical (minimal-action) and semantic (certificate-based)
constraints. Third, the CH-layering viewpoint is used to organize meaning,
safety, metaphor, and higher-order constraints into a single law-space–first
architecture.

On top of these constructions, the volume develops a certificate-based AI
framework for HACA/PACER systems. Property-layer Jacobiator-gaps, LA/LB/LC
homotopy completeness levels, and certificate bundles over law-space sectors
are combined with GRL/MDQ-based structural coordinates to define reading-path
certificates for PACER runs. Homotopy networks of packs and law-space
certificates are then used to analyze behavior on large documents and
collections of documents (standards, technical specifications, regulatory
corpora), supporting certificate-based evaluation pipelines for tasks such as
version-to-version comparison, audit evidence generation, and rollback
decisions. In this way, semantic behaviour is not merely learned from data,
but can, at least in principle, be stated and verified in a mathematically
controlled way.


\chapter*{Preface}

The GaoZheng G-Framework and GaoZheng G-Algebra project~\cite{GaoZheng2025GFramework} is conceived as
an integrated, multi-volume construction. The aim is to place a wide
range of mathematical, physical, biological, and cognitive systems
under a single law-space–first architecture, in which ``laws'' and
``properties'' are treated as geometric objects equipped with
connections, curvatures, and higher homotopy data.

The \emph{pure mathematics volume}~\cite{GaoZheng2025GFramework} establishes the foundational layer of
this architecture. It introduces the GaoZheng law-space~$\GZLS$, the
principal-bundle–based generalized noncommutative Lie algebra
(PFB-GNLA), and their associated law-connections and curvatures. Within
this setting, Jacobiator-gaps are formalized as obstructions in
generalized Lie-algebraic structures, and CH-layering is used to
organize different logical and structural regimes into a stratified
law-space geometry. This volume provides the core algebraic–geometric
framework on which all subsequent applications are built.

\emph{Applications Volume~I} (GRL, quantum computing, and physics)~\cite{GaoZhengAppVol1}
instantiates GaoZheng Reinforcement Learning (GRL) in law-space and
introduces differential dynamical quanta \\(MDQ) as local quanta of
law-aware dynamics along law-space paths. GRL path integrals are used to
express minimal-action principles for physical and quantum systems,
including applications to quantum circuit models and superconductivity.
From the standpoint of the present volume, Volume~I supplies the GRL and
MDQ machinery that will later be adapted from physical state–action
spaces to lexical and semantic operator monoids.

\emph{Applications Volume~II} (LBOPB, life-science monoids, and
generative precision medicine)~\cite{GaoZhengAppVol2} exports the same law-space and GRL/MDQ
framework to life-science domains. It introduces the Life Bio-Operator
Principal Bundle (LBOPB) and related monoids, together with
pharmacokinetic/pharmacodynamic (PK/PD) structures and pharmacogenomics
monoids (PGOM). Law-space trajectories are used to model pipelines from
pathology to reverse drug design; homotopy and MDQ-based minimal-action
principles provide a structured language for generative precision
medicine. For the present volume, LBOPB and its monoidal structures
serve as prototypical examples of domain-specific law-objects~\cite{GaoZhengAppVol2} that can
be coupled to cognitive and linguistic law-objects in a unified
law-space.

\emph{Applications Volume~III}, the present volume, turns to
\emph{cognitive, linguistic, and reasoning systems}. Its focus is on a
law-space description of hierarchical cognitive architectures and
language-based agents, with special attention to character-level
reinforcement learning, PACER-style context iteration, and certificate-
based AI.

At the core of this volume are \emph{hierarchical algebraic cognitive
architectures} (HACA) and \emph{PACER} (Pack-Aligned
Compressive–Expansion Reasoner) engines:
\begin{itemize}
  \item HACA law-objects are defined as multi-layer law-space objects
        with explicit symbolic, semantic, and decision layers
        $L^{(\mathrm{sym})}\to L^{(\mathrm{sem})}\to L^{(\mathrm{dec})}$
        supported on sectors $B_L\subseteq\GZLS$. They carry multi-scale
        operator universes $\Omega(L)$ (from character-level operators
        to paragraph-level KAT paths), semantic structure spaces
        $\mathcal{S}(L)$, property layers $\mathrm{Prop}(L)$, and
        triple-layer connections~$\GZTLC_L$.
  \item PACER engines implement compressive–alignment–expansion–
        recompression cycles on lexical KAT action monoids
        $\mathcal{M}\subseteq\mathrm{End}(\Sigma^*)$, viewed as the
        symbolic backbone of HACA. In particular, character-level and
        token-level operators, together with higher-level packs and
        sentence operators, are treated as law-space–aware objects rather
        than mere syntactic tokens. This perspective provides a natural
        landing zone for \emph{char-level RL} schemes: character- and
        token-level GRL path integrals can be formulated directly on
        KAT action monoids and coupled to higher semantic layers via
        HACA law-objects.
\end{itemize}

On top of this structural layer, the volume develops \emph{two coupled
GRL path integrals}: a lexical GRL path integral on $\mathcal{M}$
governing minimal-action organization of KAT word-pack and sentence
structures, and a semantic GRL path integral on paragraph-level KAT
paths $\Omega^{(3)}(L)$ governing the admissibility of long semantic
expressions with respect to semantic certificates and law-space sectors.
Both are formulated within the PFB-GNLA background and share the same
law-space connections and curvatures, but they enforce distinct
constraints on lexical and semantic dynamics.

A central theme of the volume is the use of \emph{Jacobiator-gap
invariants}, \emph{LA/LB/LC homotopy completeness levels}, and
\emph{certificate bundles} as the backbone of a \emph{certificate-based
AI} framework. Property-layer Jacobiator-gaps and their higher homotopy
fillings are used to construct operator bundles over property
intersections and certificate bundles over law-space sectors. These, in
turn, provide reading-path certificates for PACER runs, with explicit
constraints on semantic functors, MDQ budgets, and law-space safe
sectors. Homotopy networks of packs and reading-path certificates enable
a geometric treatment of interpretability, audit, and governance for
language-based systems.

From an engineering standpoint, the intended use of this volume is two-
fold:
\begin{itemize}
  \item It provides a law-space and operator-theoretic envelope within
        which \emph{char-level RL} and \emph{PACER-style context
        iteration} can be embedded as instances of GRL on KAT action
        monoids and HACA law-objects, rather than as ad hoc heuristics
        on raw token sequences.
  \item It proposes a certificate-based AI architecture in which
        cognitive and language systems are observed and constrained via
        law-space certificates. Reading-path certificates, occupancy
        measures, value geometry, and Jacobiator-gap invariants serve as
        machine-checkable counterparts to concepts such as ``safety
        envelope'', ``test coverage'', and ``audit trail'', making it
        possible—at least in principle—to connect engineering practice
        to governance and policy interfaces in a mathematically
        controlled way.
\end{itemize}

The exposition assumes familiarity with the law-space and GRL/MDQ
foundations developed in the earlier volumes, but it is written so that
the core HACA/PACER constructions and certificate frameworks can be
followed independently once the basic notation has been introduced. The
hope is that this volume will be useful both as a continuation of the
mathematical program initiated in the pure mathematics volume and as a
conceptual reference for engineers and governance practitioners seeking
to build law-space–aware, certificate-based cognitive and language
systems.


\tableofcontents

\mainmatter



\part[HACA Law-Objects and Semantic Gauge Fields]{HACA Law-Objects and Semantic Gauge Fields\\
      as Applications of the GaoZheng G-Framework}

\chapter[From Law-Space Geometry to HACA Law-Objects and Semantic Layers]{From Law-Space Geometry to HACA Law-Objects and Semantic Layers:\\
         Hierarchical Algebraic Cognitive Architectures in the GaoZheng G-Framework}


\section[HACA Law-Objects and Semantic Layers in GZ-LS]{HACA Law-Objects and Semantic Layers in GZ-LS}


In this chapter we view hierarchical algebraic cognitive architectures
(HACA) as law-level objects internal to the GaoZheng law-space~\GZLS.
The goal of this section is to fix notation for the category
$\mathcal{L}_{\mathrm{HACA}}$ of HACA law-objects and to record how the
symbolic, semantic, and decision layers of a HACA instance sit over
~\GZLS. The subsequent sections will then build property layers,
property-changing operators, and Jacobiator-gap invariants on top of
this basic structure.

\begin{definition}[HACA Law-Objects and the Category $\mathcal{L}_{\mathrm{HACA}}$]
\label{def:haca-law-objects-category}
A \emph{HACA law-object} $L$ in the GaoZheng G-Framework consists of the
following data:
\begin{enumerate}
  \item A distinguished law-space sector $B_L\subseteq\GZLS$, called the
        \emph{law-support} of $L$.
  \item A triple of law-level components
        \[
          L^{(\mathrm{sym})},\quad
          L^{(\mathrm{sem})},\quad
          L^{(\mathrm{dec})},
        \]
        called the \emph{symbolic}, \emph{semantic}, and
        \emph{decision} layers of $L$, each supported on $B_L$ and
        connected by law-space maps
        \[
          L^{(\mathrm{sym})}\longrightarrow L^{(\mathrm{sem})}
          \longrightarrow L^{(\mathrm{dec})}
        \]
        encoding, respectively, symbol-to-meaning and
        meaning-to-decision transitions.
  \item A multi-scale operator universe $\Omega(L)$ together with a
        semantic structure space $\mathcal{S}(L)$ and maps
        \[
          \operatorname{Sem}_L:\Omega(L)\to\mathcal{S}(L),
          \qquad
          \operatorname{Law}_L:\mathcal{S}(L)\to B_L,
        \]
        whose composition
        $\Phi_L=\operatorname{Law}_L\circ\operatorname{Sem}_L:
         \Omega(L)\to B_L$
        records the law-space footprint of operator configurations. The
        precise shape of $\Omega(L)$ is spelled out in
        Definition~\ref{def:haca-operator-universe} later in this
        chapter.
  \item A semantic degeneration of the PFB-GNLA over $B_L$ in the sense
        of Definition~\ref{def:semantic-degeneration-pfb-gnla} below,
        so that $L$ inherits a semantic G-Algebra sector
        $\GAlgebra^{\mathrm{sem}}_L$ and a constrained principal
        subbundle $P^{\mathrm{sem}}_L$ compatible with the global
        law-space connection.
\end{enumerate}
A \emph{HACA morphism} $f:L\to L'$ is a law-space map compatible with
all of these structures: it maps $B_L$ into $B_{L'}$, preserves the
three-layer decomposition (up to admissible refinement),
commutes with the semantic and law projections
$\operatorname{Sem}_L,\operatorname{Law}_L$ up to controlled homotopy,
and respects the semantic degenerations of the PFB-GNLA.

The resulting category is denoted by
\[
  \mathcal{L}_{\mathrm{HACA}}.
\]
\end{definition}

\begin{remark}[Three-Layer Architecture and CH-Layering]
\label{rem:haca-three-layer-vs-ch-layer}
The triple
$L^{(\mathrm{sym})}\to L^{(\mathrm{sem})}\to L^{(\mathrm{dec})}$ is the
categorical counterpart of the three-layer HACA architecture used in the
engineering manuscripts: the symbolic layer governs characters,
tokens, and word-packs; the semantic layer governs paragraph graphs,
KAT models and other structured meaning representations; the decision
layer governs action-selection and policy updates. In the CH-layering
picture adopted later in this section, these three layers are
consistently embedded into the $l_1/l_2/l_3$ stratification of
law-space: the symbolic layer is anchored at the raw operator level,
the semantic layer at the property level, and the decision layer at the
homotopy-constrained level where safety and robustness enter as
additional structure.
\end{remark}



\section[Property Layers, Property-Changing Operators, and Jacobiator-Gaps in HACA]{Property Layers, Property-Changing Operators, and \\Jacobiator-Gaps in HACA}


\begin{definition}[Property Layer for HACA Law-Objects]
\label{def:haca-property-layer}
Let $\mathcal{L}_{\mathrm{HACA}}$ denote the category of HACA law-objects
(hierarchical algebraic cognitive architectures viewed as law-level
structures in the GaoZheng G-Framework). A \emph{property layer} on
$\mathcal{L}_{\mathrm{HACA}}$ is given by a functor
\[
  \mathrm{Prop} : \mathcal{L}_{\mathrm{HACA}} \longrightarrow \mathbf{Set},
\]
which assigns to each law-object $L$ a set of properties
$\mathrm{Prop}(L)$ encoding, for example, semantic invariants, structural
constraints, value baselines, or safety/compatibility conditions.  An
element $p \in \mathrm{Prop}(L)$ is called a \emph{property instance} of
the law-object $L$.
\end{definition}



\begin{definition}[Property-Changing Operators and Heterogeneous Property-Changing Morphisms]
\label{def:property-changing-operators}
Let $L_1,L_2 \in \mathcal{L}_{\mathrm{HACA}}$ be two law-objects. A
\emph{property-changing operator} from $L_1$ to $L_2$ is a map on the
property layer
\[
  S_{12} : \mathrm{Prop}(L_1) \longrightarrow \mathrm{Prop}(L_2),
\]
with well-defined composition and preimage on property sets.

A \emph{heterogeneous property-changing morphism} from $L_1$ to $L_2$ is a
pair
\[
  (f_{12}, S_{12}),
\]
where $f_{12}: L_1 \to L_2$ is a morphism in
$\mathcal{L}_{\mathrm{HACA}}$ (a law-space map) and $S_{12}$ is a
property-changing operator as above. Composition is defined componentwise:
given $(f_{12},S_{12}) : L_1 \to L_2$ and $(f_{23},S_{23}) : L_2 \to L_3$,
the composite is
\[
  (f_{23},S_{23}) \circ (f_{12},S_{12})
  \;:=\;
  (f_{23}\circ f_{12},\, S_{23}\circ S_{12}),
\]
where $S_{23}\circ S_{12}$ acts on $\mathrm{Prop}(L_1)$ by composition of
maps. In this way, law-space morphisms and property-level transformations
are coupled into a single ``heterogeneous'' structure.
\end{definition}



\begin{definition}[Jacobiator-Gap on the Property Layer for HACA]
\label{def:haca-jacobiator-gap-property-layer}
Let
\[
  L_1 \xrightarrow{(f_{12},S_{12})} L_2
      \xrightarrow{(f_{23},S_{23})} L_3,
  \qquad
  L_1 \xrightarrow{(f_{13},S_{13})} L_3
\]
be three heterogeneous property-changing morphisms in
$\mathcal{L}_{\mathrm{HACA}}$ as in
Definition~\ref{def:property-changing-operators}. The
\emph{Jacobiator-gap on the property layer} for the triple
$(L_1,L_2,L_3)$ is the defect
\[
  J_{\mathrm{prop}}
  \;:=\;
  S_{23} \circ S_{12} \;-\; S_{13},
\]
measuring to what extent the three-step property evolution
$L_1 \to L_2 \to L_3$ fails to coincide with the ``direct'' property
transformation $L_1 \to L_3$.

We say that the Jacobiator-gap at $(L_1,L_2,L_3)$ is
\emph{homotopy-filled on the property layer} if there exists homotopy data
$H$ (property-level higher data) such that the triple intersection
\[
  \mathrm{Prop}(L_1)
  \;\cap\;
  S_{12}^{-1}\bigl(\mathrm{Prop}(L_2)\bigr)
  \;\cap\;
  (S_{23}\circ S_{12})^{-1}\bigl(\mathrm{Prop}(L_3)\bigr)
  \;\neq\; \varnothing,
\]
i.e.\ there exists at least one configuration of properties that is
simultaneously compatible with the two-step path
$L_1 \to L_2 \to L_3$ and with the direct path $L_1 \to L_3$ after
incorporating $H$.
\end{definition}

\begin{remark}[Relation to the Jacobiator-Gap in the Pure Mathematics Volume]
\label{rem:haca-jacobiator-gap-puremath}
In the pure mathematics volume~\cite{GaoZheng2025GFramework} of the GaoZheng G-Framework, the
Jacobiator-gap is defined at the level of generalized Lie-algebraic and
$L_\infty$ structures (PFB-GNLA and related objects), capturing a
three-fold failure of strict associativity in law-level brackets. The
property-layer Jacobiator-gap in
Definition~\ref{def:haca-jacobiator-gap-property-layer} should be viewed as
a HACA-specific reinterpretation of the same phenomenon: instead of
measuring a failure of algebraic closure on brackets, it measures a failure
of closure of heterogeneous property-changing data under three-step
composition.

In this volume we will use the property-layer Jacobiator-gap, together
with its homotopy-filling, as the basic invariant behind homotopy-complete
HACA/PACER reading paths, where PACER (Pack-Aligned Compressive–Expansion
Reasoner) is the pack-aligned compressive–expansion reasoning engine
built on hierarchical algebraic cognitive architectures (HACA), and behind
the associated logical-chain metrics and certificate-based AI. A
more detailed algebraic theory of property-changing operators and higher
law-space morphisms is developed throughout this volume and is not required
in the life-science-oriented Volume II.
\end{remark}



\begin{proposition}[Heterogeneous Homotopy Smoothing and Gaugeable Operator Bundles]
\label{prop:gaugeable-operator-bundle-from-property-intersections}
Let $\mathcal{L}_{\mathrm{HACA}}$ be the category of HACA law-objects with
heterogeneous property-changing morphisms $(f_{ij},S_{ij})$ as in
Definition~\ref{def:property-changing-operators}. For any finite chain
\[
  L_0 \xrightarrow{(f_{01},S_{01})} L_1
      \xrightarrow{(f_{12},S_{12})} \cdots
      \xrightarrow{(f_{n-1,n},S_{n-1,n})} L_n
\]
we define its property-layer intersection set
\[
\begin{aligned}
I(L_0,\dots,L_n)
:={}& \mathrm{Prop}(L_0)
      \;\cap\;
      S_{01}^{-1}\bigl(\mathrm{Prop}(L_1)\bigr)
      \;\cap\;
      (S_{12}\circ S_{01})^{-1}\bigl(\mathrm{Prop}(L_2)\bigr) \\
 &\;\cap\;\cdots\;\cap\;
      (S_{n-1,n}\circ\cdots\circ S_{01})^{-1}
      \bigl(\mathrm{Prop}(L_n)\bigr).
\end{aligned}
\]
Intuitively, $I(L_0,\dots,L_n)$ consists of those property configurations
that remain compatible along the entire heterogeneous chain
$L_0\to\cdots\to L_n$.

Let ``good chains'' be those chains $(L_0,\dots,L_n)$ that are
homotopy-admissible in the sense of this volume: they admit $n$-th order
homotopy smoothing data that fill the relevant Jacobiator-gaps on the
property layer, and they satisfy the local admissibility and quantitative
constraints (e.g.\ non-empty intersections, extendability, and bounded
gap-energy). Define the union of property intersections over all good
chains by
\[
  \mathcal{U}
  :=
  \bigcup_{\text{good chains }(L_0,\dots,L_n)}
  I(L_0,\dots,L_n).
\]

Suppose in addition that the following conditions hold:
\begin{enumerate}
  \item \textbf{Local sufficiency and stability.}
        For every good chain $(L_0,\dots,L_n)$, the intersection
        $I(L_0,\dots,L_n)$ is non-empty, and as the chain is perturbed
        within the class of good chains, the corresponding intersections
        vary in a locally stable manner so that $\mathcal{U}$ inherits a
        reasonable topology from the ambient law-space and property layer.
  \item \textbf{Operator fibers and structural group.}
        For each $p\in\mathcal{U}$, let
        \[
        \begin{aligned}
          \mathcal{F}_p
          := \bigl\{\, T \in \mathcal{M} \;\big|\;&
          T \text{ is an admissible operator at }p,\\
          & T \text{ preserves the good-chain admissibility constraints}
          \,\bigr\},
        \end{aligned}
        \]
        where $\mathcal{M}$ is the relevant operator monoid (e.g.\ a
        Lex-KAT action monoid or an LBOPB operator monoid). Assume there
        exists a subgroup $G_{\mathrm{op}}\subset\mathcal{M}$ that acts
        freely and transitively on each non-empty fiber $\mathcal{F}_p$
        (at least locally), so that the fibers can be identified (locally)
        with $G_{\mathrm{op}}$.
  \item \textbf{$n$-th order homotopy smoothing of Jacobiator-gaps.}
        The higher-homotopy data assigned to good chains trivialize the
        relevant Jacobiator-gap cohomology classes over $\mathcal{U}$:
        in other words, the property-layer Jacobiator-gap
        $J_{\mathrm{prop}}$ (as in
        Definition~\ref{def:haca-jacobiator-gap-property-layer})
        becomes a coboundary in the appropriate cohomology over
        $\mathcal{U}$.
\end{enumerate}
Then there exists a principal $G_{\mathrm{op}}$-bundle
\[
  \pi_{\mathrm{op}} : P_{\mathrm{op}} \longrightarrow \mathcal{U}
\]
whose fibers coincide (locally) with the admissible operator fibers
$\mathcal{F}_p$. In particular, $\mathcal{U}$ can be interpreted as the
base of a gaugeable operator bundle built from heterogeneous
property-changing data that have been homotopy-smoothed up to order $n$.

Moreover, the homotopy completeness levels LA/LB/LC introduced in
Definition~\ref{def:homotopy-completeness-levels} admit the following
geometric interpretation in this setting:
\begin{itemize}
  \item Along a single homotopy-complete path (Level A), the restriction
        of $P_{\mathrm{op}}$ to the image of that path is automatically
        trivial (as the path is contractible), so a local gauge can always
        be chosen there.
  \item In the presence of a homotopy common core
        $\mathcal{C}\subset\mathcal{U}$ (Level B), all good paths pass
        through a region where $P_{\mathrm{op}}$ admits a common local
        trivialization, so a shared gauge can be used as a semantic and
        structural ``core chart'' for auditing and comparison.
  \item When a normative subspace $\mathcal{N}\subset\mathcal{U}$ (Level C)
        is such that all good paths in a given regime can be homotopically
        projected into $\mathcal{N}$, the restriction
        $P_{\mathrm{op}}\vert_{\mathcal{N}}$ is (by design) globally
        gaugeable: one can choose a global section or, at least, a canonical
        baseline path $\gamma_\ast$ to which all good paths are semantically
        homotopy-equivalent. This realizes the strongest form of
        gauge-fixability under the chosen baseline.
\end{itemize}
\end{proposition}

A detailed proof of
Proposition~\ref{prop:gaugeable-operator-bundle-from-property-intersections}
is given in Section~\ref{app:sec-operator-bundles} of the Technical Appendix.

In informal terms, the union of property intersections generated by
heterogeneous property-changing morphisms, together with their $n$-th
order homotopy smoothing of Jacobiator-gaps, can be regarded, under the
above conditions, as the base (or a subspace of the base) of a gaugeable
operator bundle. The levels LA/LB/LC, introduced later in
Definition~\ref{def:homotopy-completeness-levels}, then quantify at which
scales such an operator bundle admits canonical gauges (local
trivializations or global sections under the chosen baseline): along single
paths, around homotopy common cores, or globally on normative subspaces.





\begin{definition}[Multi-Scale Operator Universes in HACA and PACER]
\label{def:haca-operator-universe}
Let $L$ be a HACA law-object. We write
\[
  \Omega^{(0)}(L),\ \Omega^{(1)}(L),\ \Omega^{(2)}(L),\ \Omega^{(3)}(L)
\]
for the collections of operator configurations at the levels of
characters/tokens, word- or pack-level operators, sentence-level
structures, and paragraph- or reading-path operators (KAT path integrals)
respectively. The \emph{multi-scale operator universe} of $L$ is the
disjoint union
\[
  \Omega(L)
  \;:=\;
  \Omega^{(0)}(L)\ \sqcup\ \Omega^{(1)}(L)\ \sqcup\ \Omega^{(2)}(L)\ \sqcup\ \Omega^{(3)}(L).
\]
An element $\omega\in\Omega(L)$ is called an \emph{operator configuration}.



\end{definition}

\begin{definition}[Semantic and Law-Space Projections]
\label{def:haca-sem-law-projections}
For each HACA law-object $L$, we fix:
\begin{itemize}
  \item a \emph{semantic structure space} $\mathcal{S}(L)$ (for instance a
        KAT model, a paragraph-graph of PACER nodes and edges, or a
        higher-level semantic state space), and
  \item maps
        \[
          \operatorname{Sem}_L : \Omega(L) \longrightarrow \mathcal{S}(L),
          \qquad
          \operatorname{Law}_L : \mathcal{S}(L) \longrightarrow B_L
          \subseteq \GZLS,
        \]
        where $B_L$ is the portion of law-space $\GZLS$ supporting $L$.
\end{itemize}
The composition
\[
  \Phi_L
  \;:=\;
  \operatorname{Law}_L \circ \operatorname{Sem}_L
  : \Omega(L) \longrightarrow B_L
\]
associates to each operator configuration $\omega\in\Omega(L)$ a
law-space point (or sector) $\Phi_L(\omega)\in B_L$ that records its
semantic law-level footprint.



\end{definition}

\begin{definition}[Semantic Predicates and Truth-Valued Functors]
\label{def:haca-semantic-predicates}
With notation as in Definition~\ref{def:haca-operator-universe}
and Definition~\ref{def:haca-sem-law-projections}, a
\emph{semantic predicate} for a HACA law-object $L$ is a map
\[
  \chi:\mathcal{S}(L)\longrightarrow \{\mathrm{true},\mathrm{false}\}
\]
satisfying the following conditions:
\begin{enumerate}
  \item \textbf{Functoriality with respect to semantic operations.}
        For each primitive semantic operation
        \[
          \mathrm{op}_S : \mathcal{S}(L)^n \longrightarrow \mathcal{S}(L)
        \]
        used in the HACA/PACER semantics (for example, KAT composition,
        control-flow join, or paragraph-graph edge composition), there
        exists a Boolean operation
        \[
          F_{\mathrm{op}} : \{\mathrm{true},\mathrm{false}\}^n
          \longrightarrow \{\mathrm{true},\mathrm{false}\}
        \]
        such that
        \[
          \chi\bigl(\mathrm{op}_S(s_1,\dots,s_n)\bigr)
          =
          F_{\mathrm{op}}\bigl(\chi(s_1),\dots,\chi(s_n)\bigr)
        \]
        for all $s_1,\dots,s_n\in\mathcal{S}(L)$.
  \item \textbf{Semantic closure of the ``true'' sector.}
        For each primitive semantic operation $\mathrm{op}_S$ as above,
        whenever all inputs satisfy $\chi(s_i)=\mathrm{true}$, we require
        that
        \[
          \chi\bigl(\mathrm{op}_S(s_1,\dots,s_n)\bigr)
          = \mathrm{true}.
        \]
        In other words, the ``true'' fiber
        \[
          \mathcal{S}_\chi^{\mathrm{true}}
          :=
          \chi^{-1}(\{\mathrm{true}\})
          \subseteq \mathcal{S}(L)
        \]
        is closed under the semantic operations of HACA/PACER.
\end{enumerate}
The corresponding \emph{truth-valued observable} on operator
configurations is the composite
\[
  \tilde{\chi}
  :=
  \chi \circ \operatorname{Sem}_L :
  \Omega(L) \longrightarrow \{\mathrm{true},\mathrm{false}\},
\]
and its ``true'' fiber
\[
  P_\chi
  :=
  \tilde{\chi}^{-1}(\{\mathrm{true}\})
  =
  \bigl\{\, \omega\in\Omega(L)\ \big|\ \chi(\operatorname{Sem}_L(\omega))
  = \mathrm{true} \,\bigr\}
\]
is called the \emph{semantic property class} associated with~$\chi$.

\end{definition}

\begin{definition}[Observable Property Families and Property Partitions]
\label{def:haca-observable-properties}
Let $L$ be a HACA law-object with operator universe $\Omega(L)$ and
semantic--law map $\Phi_L$ as above. We fix a family of observable
functionals
\[
  \mathsf{Obs}(L)\subseteq \bigl\{\, f:\Omega(L)\to V_f \,\bigr\},
\]
where each $f\in\mathsf{Obs}(L)$ takes values in some set $V_f$ (for
example $\{0,1\}$, a totally ordered set, or a metric space). We allow in
particular the truth-valued observables
\[
  \tilde{\chi}
  =
  \chi\circ\operatorname{Sem}_L :
  \Omega(L)\to\{\mathrm{true},\mathrm{false}\}
\]
arising from semantic predicates $\chi$ as in
Definition~\ref{def:haca-semantic-predicates}.

\begin{itemize}
  \item A subset $P\subseteq \Omega(L)$ is called \emph{observable} if
        there exist finitely many observable functionals
        $f_i\in\mathsf{Obs}(L)$ with value sets $V_{f_i}$ and subsets
        $B_i\subseteq V_{f_i}$ such that
        \[
          P
          =
          \bigl\{\, \omega\in\Omega(L)\ \big|\ f_i(\omega)\in B_i
          \text{ for all }i \,\bigr\},
        \]
        and $P$ is stable under all admissible HACA/law-space homotopies
        (up to Jacobiator-gap boundaries) in the sense of
        Definition~\ref{def:haca-jacobiator-gap-property-layer} and
        Definition~\ref{def:homotopy-completeness-levels}.
        In particular, for each semantic predicate $\chi$ the semantic
        property class $P_\chi$ from
        Definition~\ref{def:haca-semantic-predicates} is an observable
        subset, closed under the induced semantic operations.
  \item We write $\mathcal{L}_{\mathrm{obs}}(L)$ for the collection of all
        observable subsets of $\Omega(L)$.
  \item Two observable subsets $P,Q\in\mathcal{L}_{\mathrm{obs}}(L)$ are
        declared equivalent, written $P\sim Q$, if they induce the same
        law-space data, for example if $\Phi_L(P)$ and $\Phi_L(Q)$ agree
        up to refinement inside $B_L$ and the same homotopy-completeness
        profile (LA/LB/LC) holds on both.
\end{itemize}
A property layer functor $\mathrm{Prop}$ as in
Definition~\ref{def:haca-property-layer} is said to be \emph{realized by
observable partitions} if for each $L$ there is a bijection
\[
  \mathrm{Prop}(L)
  \;\cong\;
  \mathcal{L}_{\mathrm{obs}}(L) / \sim,
\]
compatible with pullback along HACA morphisms. In this case we may view
each property instance $p\in\mathrm{Prop}(L)$ as the label of an
equivalence class of observable subsets
$[P]\in\mathcal{L}_{\mathrm{obs}}(L)/\sim$.

Equivalently, we can encode the same data by a ``property evaluation
map''
\[
  \pi_L : \Omega(L) \longrightarrow \mathrm{Prop}(L),
\]
whose fibers $\Omega_p(L):=\pi_L^{-1}(\{p\})$ give an admissible
partition of $\Omega(L)$ into observable law-space sectors.



\end{definition}

\begin{remark}[Trivial and Degenerate Property Partitions]
\label{rem:trivial-property-partitions}
Nothing in the raw set-theoretic picture prevents us from choosing
extreme partitions, such as the ``discrete'' partition where each
$\omega\in\Omega(L)$ forms its own property class, or the ``coarse''
partition where all of $\Omega(L)$ is assigned a single property label.
In the present framework such choices are excluded as degenerate: they
fail to provide semantic compression, are unstable under admissible
homotopies, and cannot be realized as meaningful sectors of law-space
geometry. The realizability and homotopy-stability requirements build
these exclusions into the definition of observable property families.



\end{remark}

\begin{remark}[Structural Properties as Law-Space Sectors in GZ-LS]
\label{rem:structural-properties-as-laws}
Given the semantic--law map $\Phi_L$ from
Definition~\ref{def:haca-sem-law-projections}, any admissible law-space
sector $\Sigma\subseteq B_L$ determines an observable subset
\[
  P_\Sigma
  :=
  \Phi_L^{-1}(\Sigma)
  \subseteq\Omega(L),
\]
and conversely each observable subset
$P\in\mathcal{L}_{\mathrm{obs}}(L)$ is intended to arise (up to negligible
differences) as such a pullback. In this sense a structural property of
$L$ can be regarded either as a law-space sector $\Sigma\subseteq B_L$
together with its connection and curvature data, or as the corresponding
class of operator configurations $P_\Sigma$; the two descriptions carry
the same information once we work inside the GZ-LS formalism. This is the
precise sense in which ``structural properties'' and ``structural laws''
are two equivalent descriptions of the same GZ-LS data.



\end{remark}





\noindent\textbf{Informal Summary.}
In this section we adopt a layered view of law-space, aligned with the
HACA/PACER construction. The $L_1$-layer is the raw operator universe
of characters, word-packs, sentences, and paragraph paths. The $L_2$-layer
resolves this universe into a semantic property layer, addressing the
problem of meaning. The $L_3$-layer adds homotopy and Jacobiator-gap
data, addressing robustness and safety as homotopy-level constraints
rather than merely pointwise predicates. The $L_4$-layer is proposed as
a layer for metaphor and cross-domain connections, and higher layers
$L_n$ can, in principle, be used to encode additional forms of structure
and constraint (value geometry, ethics, multi-agent games, and so on).



Heuristically, one may align the law-space layering perspective with the
HACA/PACER construction as follows. The $L_1$-layer is the raw operator
universe
\[
  \Omega(L)
  =
  \Omega^{(0)}(L)\sqcup\Omega^{(1)}(L)\sqcup\Omega^{(2)}(L)\sqcup\Omega^{(3)}(L),
\]
consisting of characters, word-packs, sentence structures, and paragraph
paths. The $L_2$-layer resolves $\Omega(L)$ into a semantic property
layer $\mathrm{Prop}(L)$, built from basic semantic functors and an
intentional power-set construction; in this sense, the $L_2$-layer
addresses the \emph{problem of meaning}. The $L_3$-layer adds homotopy
and Jacobiator-gap data on top of $L_2$, providing a framework in which
robustness and safety can be formulated as homotopy-level constraints
rather than merely pointwise predicates, thus addressing the
\emph{problem of safety}. A further $L_4$-layer may then be used to
encode metaphorical and cross-domain connections by considering
structure-preserving functors between law-objects and law-sectors,
i.e.\ addressing the \emph{problem of metaphor and analogy}. In
principle, higher layers $L_n$ can be introduced to capture additional
forms of structure and constraint (value geometry, ethics, multi-agent
games, and so on), each layer resolving problems that cannot be fully
treated on the layers below.



\begin{remark}[Layered View: Operator States, Properties, and Homotopy]
\label{rem:ch-layering-l1-l3}
In the layered picture inherited from the CH-layering perspective of the
earlier volumes, one may summarize the role of semantic functors and
observable subsets as follows. At the $l_1$-level, the HACA law-object $L$
carries the multi-scale operator universe
\[
  \Omega(L)
  =
  \Omega^{(0)}(L)\sqcup\Omega^{(1)}(L)\sqcup\Omega^{(2)}(L)\sqcup\Omega^{(3)}(L),
\]
which is the raw state space of characters, word-packs, sentence
structures, and paragraph paths. This level is too rich to be handled
directly for semantic or safety decisions.

At the $l_2$-level, one selects a family of basic semantic functors
$\{F_\alpha\}_\alpha$ (including truth-valued predicates as in
Definition~\ref{def:haca-semantic-predicates}) and uses them, together
with observable functionals, to carve out an ``intentional power set''
\[
  \mathcal{L}_{\mathrm{obs}}(L)\subseteq\mathcal{P}(\Omega(L))
\]
of semantically meaningful, homotopy-stable, and law-space–realizable
subsets. The resulting property layer
\[
  \mathrm{Prop}(L)\cong\mathcal{L}_{\mathrm{obs}}(L)/\sim
\]
then serves as an $l_2$-resolution of the $l_1$ state space: many
questions about operator configurations are transferred to questions
about semantic property assignments induced by the basic functors.

When this $l_2$-resolution is still insufficient—for instance, in
questions of robustness and safety—one may pass to an $l_3$-layer based
on homotopy and Jacobiator-gap data. At this level, properties are no
longer evaluated only pointwise on individual paths, but classified up
to controlled homotopy equivalence (LA/LB/LC) within the relevant
law-space sectors. Safety can then be formulated as a homotopy-level
constraint: acceptable semantic evolutions must lie in the same
controlled homotopy class as a certified safe baseline, rather than
merely satisfying pointwise $l_2$-properties along a single
representative path.

\end{remark}

\begin{definition}[Semantic Gauge Fields of Functors]
\label{def:semantic-gauge-field}
Let $L$ be a HACA law-object with operator universe $\Omega(L)$, semantic
structure space $\mathcal{S}(L)$, and law-space sector $B_L\subseteq\GZLS$
as in Definitions~\ref{def:haca-operator-universe}
and~\ref{def:haca-sem-law-projections}. A \emph{semantic functor field}
(or \emph{semantic gauge field}) on $L$ consists of the following data:
\begin{enumerate}
  \item A family of value spaces $\{V_\alpha\}_\alpha$ (for example
        $\{0,1\}$, intervals in $\mathbb{R}$, or finite label sets) and
        corresponding semantic functors
        \[
          F_\alpha : \mathcal{S}(L)\longrightarrow V_\alpha,
        \]
        each of which is functorially compatible with the semantic
        operations on $\mathcal{S}(L)$ in the sense that for every
        primitive semantic operation
        $\mathrm{op}_S : \mathcal{S}(L)^n\to\mathcal{S}(L)$ there is an
        induced operation
        \[
          G_{\mathrm{op}}^\alpha : V_\alpha^n\to V_\alpha
        \]
        such that
        \[
          F_\alpha\bigl(\mathrm{op}_S(s_1,\dots,s_n)\bigr)
          =
          G_{\mathrm{op}}^\alpha\bigl(F_\alpha(s_1),\dots,F_\alpha(s_n)\bigr)
        \]
        holds for all $s_1,\dots,s_n\in\mathcal{S}(L)$.
        Truth-valued semantic predicates
        $\chi:\mathcal{S}(L)\to\{\mathrm{true},\mathrm{false}\}$ from
        Definition~\ref{def:haca-semantic-predicates} are the special
        case $V_\alpha=\{\mathrm{true},\mathrm{false}\}$.

  \item A compatible law-space realization of these functors, in the form
        of a map
        \[
          \mathcal{A}_L :
          B_L \longrightarrow \mathsf{Conn}(\mathcal{F}_L),
        \]
        where $\mathcal{F}_L$ is the bundle over $B_L$ whose fiber at
        $b\in B_L$ consists of the restrictions of the semantic functors
        $F_\alpha$ to the semantic configurations mapping into $b$, and
        $\mathsf{Conn}(\mathcal{F}_L)$ denotes the space of connections
        on this bundle. Intuitively, $\mathcal{A}_L$ specifies how the
        family $\{F_\alpha\}_\alpha$ is parallel transported along
        law-space paths in $B_L$, so that functor values vary
        covariantly with respect to the law-space connection on $L$.
\end{enumerate}
For each semantic functor $F_\alpha$ and each operator configuration
$\omega\in\Omega(L)$, the composite
\[
  F_\alpha(\operatorname{Sem}_L(\omega))
\]
is called the \emph{semantic evaluation} of $\omega$ with respect to
$F_\alpha$. For a given law-space path $\gamma:[0,1]\to B_L$ and an
operator path $\omega_t$ projecting to $\gamma$, the covariant derivative
of $F_\alpha$ along $\gamma$,
\[
  \nabla_{\dot{\gamma}} F_\alpha,
\]
as determined by the connection $\mathcal{A}_L$, is called the
\emph{semantic gradient} of $F_\alpha$ along $\gamma$. When
$V_\alpha=\{\mathrm{true},\mathrm{false}\}$, the semantic gradient reduces
to the controlled appearance or disappearance of the corresponding
semantic property class along the path; when $V_\alpha\subseteq\mathbb{R}$
or $V_\alpha\subseteq[0,1]$, it provides a quantitative, law-space–aware
measure of how the semantics of $\omega_t$ changes along $\gamma$.

\end{definition}

\begin{definition}[Noncommutative Semantic Functor Algebra]
\label{def:noncommutative-semantic-functor-algebra}
Let $L$ be a HACA law-object with semantic structure space
$\mathcal{S}(L)$ and law-space sector $B_L\subseteq\GZLS$ as before.
Consider a family of semantic functors
\[
  \{F_\alpha:\mathcal{S}(L)\to V_\alpha\}_\alpha
\]
forming a semantic gauge field in the sense of
Definition~\ref{def:semantic-gauge-field}. We define the
\emph{semantic functor space}
\[
  \mathfrak{F}_L
  :=
  \bigl\{\, F:\mathcal{S}(L)\to V \ \big|\ F
  \text{ is a semantic functor on }L \,\bigr\}.
\]

A \emph{noncommutative semantic functor algebra} on $L$ is the datum of
\begin{enumerate}
  \item a bilinear operation
        \[
          [\cdot,\cdot]_{\mathrm{sem}} :
          \mathfrak{F}_L\times\mathfrak{F}_L\longrightarrow\mathfrak{F}_L,
        \]
        called the \emph{semantic bracket}, and
  \item for each triple $(F,G,H)\in\mathfrak{F}_L^3$, a
        \emph{Jacobiator-gap}
        \[
          J(F,G,H)
          :=
          [F,[G,H]_{\mathrm{sem}}]_{\mathrm{sem}}
          + [G,[H,F]_{\mathrm{sem}}]_{\mathrm{sem}}
          + [H,[F,G]_{\mathrm{sem}}]_{\mathrm{sem}},
        \]
\end{enumerate}
such that:
\begin{itemize}
  \item (Noncommutativity) For generic pairs $(F,G)$ one has
        \(
          [F,G]_{\mathrm{sem}}\neq 0
        \)
        and \\$[F,G]_{\mathrm{sem}}\neq [G,F]_{\mathrm{sem}}$; in
        particular, the functors do not commute in general.
  \item (Closure) For all $F,G\in\mathfrak{F}_L$ the bracket
        $[F,G]_{\mathrm{sem}}$ is again a semantic functor, so that
        $\mathfrak{F}_L$ is closed under $[\cdot,\cdot]_{\mathrm{sem}}$.
  \item (Controlled Jacobiator-gap) The Jacobiator-gap $J(F,G,H)$ is
        either zero or lies in a distinguished subspace
        $\mathfrak{Z}_L\subseteq\mathfrak{F}_L$ which is compatible with
        the law-space connection on $B_L$ (e.g.\ a central or controlled
        sector in the sense of the GaoZheng generalized noncommutative
        Lie algebra). In this sense $\mathfrak{F}_L$ carries the
        structure of a generalized, possibly $H$-twisted, noncommutative
        Lie algebra with Jacobiator-gaps.
\end{itemize}

\end{definition}

\begin{remark}[Semantic Functors as Gauge Quanta and Their Mutual Constraints]
\label{rem:semantic-functors-as-quanta}
In this picture, each individual semantic functor $F_\alpha$ may be
regarded as an elementary ``mode'' or ``quantum'' of the semantic gauge
field: it specifies one direction of semantic evaluation on the
multi-scale operator universe $\Omega(L)$ via
\[
  \omega \longmapsto F_\alpha(\operatorname{Sem}_L(\omega)).
\]
The semantic bracket
$[F_\alpha,F_\beta]_{\mathrm{sem}}$ measures the noncommutativity of
applying the semantic effects represented by $F_\alpha$ and $F_\beta$ in
different orders on the same semantic configurations. The fact that this
bracket is generically nonzero reflects a strictly noncommutative
structure on the semantic side: semantic evaluations in different
``directions'' do not, in general, commute.

The closure and Jacobiator-gap conditions ensure that the family
$\{F_\alpha\}_\alpha$ is not just a collection of independent labels,
but a noncommutative algebra of semantic modes constrained by the
underlying law-space geometry. In this sense, semantic functors play a
role analogous to fundamental quanta in gauge field theory, while the
noncommutative semantic functor algebra encodes their mutual
interactions and constraints within the GZ-LS framework.

\end{remark}

\begin{definition}[PFB-GNLA Background and HACA Law-Sectors]
\label{def:pfb-gnla-background-haca-sector}
We recall that the \\GaoZheng G-Framework is built on a principal-bundle–
based generalized noncommutative Lie algebra (PFB-GNLA) consisting of:
\begin{itemize}
  \item a law-space manifold $\GZLS$;
  \item a principal bundle
        \[
          \pi_{\mathrm{law}} : P_{\mathrm{law}}\longrightarrow \GZLS
        \]
        with structure group $G$;
  \item a bundle of generalized noncommutative Lie algebras
        \[
          \GAlgebra \longrightarrow \GZLS,
        \]
        whose fiber $\GAlgebra_x$ over $x\in\GZLS$ carries a (possibly
        $H$-twisted) noncommutative Lie bracket with controlled
        Jacobiator-gaps;
  \item a law-space connection $\GZLC$ and curvature $\GZLCurv$ on
        $P_{\mathrm{law}}$ compatible with $\GAlgebra$ in the sense of
        the PFB-GNLA construction of the previous volumes.
\end{itemize}

For a given HACA law-object $L$, we assume that its law-level behavior
is supported on a distinguished law-sector
\[
  B_L\subseteq \GZLS,
\]
and we write
\[
  P_{\mathrm{law}}|_{B_L}\to B_L,
  \qquad
  \GAlgebra|_{B_L}\to B_L
\]
for the restrictions of the principal bundle and the G-Algebra bundle to
$B_L$.

\end{definition}

\begin{definition}[Semantic Degeneration of PFB-GNLA]
\label{def:semantic-degeneration-pfb-gnla}
Let $L$ be a HACA law-object with associated law-sector $B_L\subseteq
\GZLS$ as in Definition~\ref{def:pfb-gnla-background-haca-sector}.
A \emph{semantic degeneration} of the PFB-GNLA over $B_L$ consists of:
\begin{enumerate}
  \item A multi-scale operator universe
        \[
          \Omega(L)
          =
          \Omega^{(0)}(L)\ \sqcup\ \Omega^{(1)}(L)\ \sqcup\ \Omega^{(2)}(L)\ \sqcup\ \Omega^{(3)}(L)
        \]
        as in Definition~\ref{def:haca-operator-universe}, together with
        a semantic structure space $\mathcal{S}(L)$ and a semantic map
        \[
          \operatorname{Sem}_L:\Omega(L)\longrightarrow \mathcal{S}(L).
        \]
  \item A law-space projection
        \[
          \operatorname{Law}_L:\mathcal{S}(L)\longrightarrow B_L
        \]
        such that the composite
        \(
          \Phi_L:=\operatorname{Law}_L\circ\operatorname{Sem}_L
          :\Omega(L)\to B_L
        \)
        is compatible with the PFB-GNLA connection $\GZLC$ on
        $P_{\mathrm{law}}|_{B_L}$ (in particular, semantic operator paths
        in $\Omega(L)$ project to law-space paths in $B_L$ along which
        $\GZLC$ defines a meaningful parallel transport).
  \item A subbundle
        \[
          \GAlgebra^{\mathrm{sem}}_L
          \longrightarrow B_L
        \]
        of the G-Algebra bundle $\GAlgebra|_{B_L}$ such that each fiber
        $\GAlgebra^{\mathrm{sem}}_{L,x}\subseteq \GAlgebra_x$ is closed
        under the generalized noncommutative Lie bracket up to a
        controlled Jacobiator-gap, i.e.
        \[
          [\GAlgebra^{\mathrm{sem}}_{L,x},
           \GAlgebra^{\mathrm{sem}}_{L,x}]
          \subseteq \GAlgebra^{\mathrm{sem}}_{L,x}
        \]
        and the Jacobiator-gap lies in a distinguished subspace
        $\mathfrak{Z}^{\mathrm{sem}}_{L,x}\subseteq
         \GAlgebra^{\mathrm{sem}}_{L,x}$ compatible with the law-space
        connection.
        We call $\GAlgebra^{\mathrm{sem}}_L$ the \emph{semantic
        G-Algebra sector} of $L$.
  \item A constrained principal subbundle
        \[
          P^{\mathrm{sem}}_L \subseteq P_{\mathrm{law}}|_{B_L}
        \]
        with structure group $G^{\mathrm{sem}}\subseteq G$, such that the
        restriction of $\GZLC$ to $P^{\mathrm{sem}}_L$ has curvature
        valued in $\GAlgebra^{\mathrm{sem}}_L$. Informally, this means
        that all non-semantic G-Algebra directions have been frozen by
        structural constraints, and only the semantic sector remains
        dynamically active over $B_L$.
\end{enumerate}
The data $(\mathcal{S}(L),\operatorname{Sem}_L,\operatorname{Law}_L,
 \GAlgebra^{\mathrm{sem}}_L,P^{\mathrm{sem}}_L)$ will be called a
\emph{semantic degeneration} of the PFB-GNLA for HACA/PACER.
\end{definition}

\begin{definition}[Semantic Structure as a Degenerate PFB-GNLA Slice]
\label{def:semantic-structure-from-degeneration}
Given a semantic degeneration
\[
  (\mathcal{S}(L),\operatorname{Sem}_L,\operatorname{Law}_L,
   \GAlgebra^{\mathrm{sem}}_L,P^{\mathrm{sem}}_L)
\]
of the PFB-GNLA over $B_L$, the \emph{semantic structure} of the HACA
law-object $L$ is defined as the following category (or higher
category, if desired):
\begin{itemize}
  \item Objects are points $s\in\mathcal{S}(L)$ together with chosen
        lifts to $P^{\mathrm{sem}}_L$ over $\operatorname{Law}_L(s)\in
        B_L$, i.e.\ pairs $(s,\tilde{s})$ with
        $\tilde{s}\in P^{\mathrm{sem}}_L$ projecting to
        $\operatorname{Law}_L(s)$.
  \item Morphisms between $(s,\tilde{s})$ and $(s',\tilde{s}')$ are
        homotopy classes of $P^{\mathrm{sem}}_L$-horizontal paths
        connecting $\tilde{s}$ to $\tilde{s}'$ whose projected base
        paths in $B_L$ are compatible with the semantic operator
        paths in $\Omega(L)$ mapping to $s$ and $s'$ via
        $\operatorname{Sem}_L$.
  \item Composition of morphisms is induced by concatenation of
        horizontal paths and the generalized noncommutative Lie
        bracket on $\GAlgebra^{\mathrm{sem}}_L$ along the path, with
        Jacobiator-gaps controlled by the semantic sector
        $\mathfrak{Z}^{\mathrm{sem}}_L$.
\end{itemize}
We denote this semantic structure category by
\[
  \mathbf{SemStr}(L).
\]
By construction, $\mathbf{SemStr}(L)$ is a degenerate slice of the
global PFB-GNLA: it inherits its noncommutative composition laws and
Jacobiator-gap structure from $\GAlgebra^{\mathrm{sem}}_L$, but all
non-semantic directions have been constrained away.

\end{definition}

\begin{proposition}[Semantic Functor Algebra as a Degenerate PFB-GNLA Representation]
\label{prop:semantic-functor-as-degenerate-rep}
Let $L$ be a HACA law-object equipped with a semantic degeneration of the
PFB-GNLA as in
Definition~\ref{def:semantic-degeneration-pfb-gnla}, and let
$\mathfrak{F}_L$ denote the noncommutative semantic functor algebra
from Definition~\ref{def:noncommutative-semantic-functor-algebra}.
Then:
\begin{enumerate}
  \item The semantic bracket
        $[\cdot,\cdot]_{\mathrm{sem}}$ on $\mathfrak{F}_L$ is induced,
        via pullback along $\operatorname{Sem}_L$ and $\operatorname{Law}_L$,
        from the generalized noncommutative Lie bracket on
        $\GAlgebra^{\mathrm{sem}}_L$; in particular, the Jacobiator-gap
        $J(F,G,H)$ in $\mathfrak{F}_L$ is the image of the Jacobiator-gap
        in $\GAlgebra^{\mathrm{sem}}_L$ under the corresponding
        representation.
  \item Each semantic functor $F_\alpha\in\mathfrak{F}_L$ can be viewed
        as a generalized representation of the semantic G-Algebra sector
        $\GAlgebra^{\mathrm{sem}}_L$ on its value space $V_\alpha$, and
        the collection $\{F_\alpha\}_\alpha$ forms a constrained,
        noncommutative family analogous to the elementary modes of a
        gauge field.
  \item In this precise sense, the semantic functor algebra
        $\mathfrak{F}_L$ is a degenerate structural-constraint
        realization of the underlying PFB-GNLA: it retains the
        noncommutative and Jacobiator-gap features of the full theory,
        but only along those directions singled out by the semantic
        sector and subject to the structural constraints defining the
        HACA/PACER law-object $L$.
\end{enumerate}

A detailed proof of
Proposition~\ref{prop:semantic-functor-as-degenerate-rep}
is given in Section~\ref{app:sec-semantic-functor-proof} of the
Technical Appendix.

\end{proposition}

\begin{remark}[Relation to Differential Dynamical Quanta in Applications Volume~I]
\label{rem:semantic-functors-vs-mdq}
In the Applications Volume~I, differential dynamical quanta (MDQ) were
introduced as local quanta of law-aware dynamics along GRL path
integrals: at each point of a law-space path, the MDQ encode the
infinitesimal contribution of the underlying generalized
noncommutative Lie algebra and its connection to the evolution of
observables. They provide a microscopic description of how PFB-GNLA
data drive the dynamics of GRL systems at the level of physical or
application-specific observables.

In the present volume, semantic functors play an analogous role at the
HACA/PACER level. Each semantic functor $F_\alpha$ can be regarded as a
differential dynamical quantum of \emph{semantic} evaluation for the
law-object $L$: given a multi-scale operator configuration
$\omega\in\Omega(L)$ and its semantic image $\operatorname{Sem}_L(\omega)$,
the value
\[
  F_\alpha(\operatorname{Sem}_L(\omega))
\]
is the infinitesimal semantic contribution attached to $\omega$ along a
law-space path in $B_L$, while the semantic gradient
$\nabla_{\dot{\gamma}}F_\alpha$ encodes how this contribution varies
under the law-space connection driven by $\GZLC$ and the semantic
sector $\GAlgebra^{\mathrm{sem}}_L$.

In this precise sense, the noncommutative semantic functor algebra
$\mathfrak{F}_L$ is the semantic-sector counterpart of the MDQ algebra
constructed in the Applications Volume~I: both are realizations of the
same PFB-GNLA background, specialized by structural constraints
(MDQ for GRL dynamics, semantic functors for HACA/PACER semantics) and
equipped with Jacobiator-gap–controlled homotopy behaviour. Semantic
functors may thus be interpreted as ``differential dynamical quanta of
meaning'' for HACA/PACER, standing in direct correspondence with the
differential dynamical quanta of GRL dynamics in the earlier volume.

\end{remark}

\section{Structural and Semantic Homotopy Across Volumes I--III}


The role of homotopy in this monograph family naturally splits into three
levels.

First, in the pure mathematics volume~\cite{GaoZheng2025GFramework}, homotopy is a classical structural
notion: it classifies maps and paths in generalized law-spaces up to
continuous deformation, and underlies the construction of homotopy groups,
cohomology classes, and bundle classifications. At this level, homotopy is
rigorous but intentionally semantic-neutral: it says that two objects are
equivalent with respect to a given geometric or algebraic structure, but
it does not express any opinion about whether the corresponding
trajectories or configurations are ``reasonable'' interpretations of a
physical, biological, or cognitive process.

Second, in Applied Mathematics Volumes I and II~\cite{GaoZhengAppVol1,GaoZhengAppVol2}, this structural homotopy
framework is instantiated in specific law-spaces. For GRL-based physics,
it supports statements such as the existence of law-space paths realizing a
transition between two regimes. For LBOPB-based life-science models, it
supports qualitative claims that there exist law-space trajectories
connecting a pathological regime to a reverse drug-design law-object. In
both cases, structural homotopy is used as a qualitative solvability
criterion: a homotopy class witnesses that the underlying system admits, in
principle, a continuous law-level evolution between two regimes. The choice
and detailed evaluation of particular trajectories is carried out using the
domain-specific functionals introduced in those volumes (actions and
performance objectives, spectral and fidelity measures, PK/PD and toxicity
constraints, PGOM, etc.), but these functionals are not yet organized into
a single, explicit semantic homotopy framework.

Third, in the present HACA volume, homotopy is upgraded to a semantic and
quantitative notion. By adding a property layer, property-changing
operators, and quantitative functionals (MDQ, value geometry, logical-chain
scores), we restrict attention to deformations that:
\begin{itemize}
  \item remain admissible on the property layer (property intersections
        do not collapse, Jacobiator-gaps can be homotopy-filled);
  \item keep quantitative measures within acceptable bounds.
\end{itemize}
Two HACA/PACER trajectories are considered homotopic only when they can be
connected by such a constrained deformation. This semantic homotopy thus
refines structural homotopy: it preserves all the rigor of the classical
notion, while adding the extra structure needed to support auditing,
filtering, and baseline selection in cognitive and language-generation
systems.

From a structural point of view, this refinement does not go beyond the
abstract framework of the pure mathematics volume: the law-spaces, bundles
and homotopy notions used here are all instances of that general shell. What
changes in this volume is the semantic layer on top of that shell. The
property functor, property-changing operators and homotopy common cores are
introduced as constrained specializations, so that domain-specific quantitative
schemes from Volumes I and II can be uniformly recast as instances of a single
semantic homotopy and logical completeness framework. In short, Volume III
does not raise the structural ceiling established by the pure mathematics
volume; it pushes the semantic floor up to meet that ceiling in a unified way.



\chapter[Semantic Gauge Fields and Triple-Layer Connections in HACA]{Semantic Gauge Fields and Triple-Layer Connections in HACA:\\
         From Law-Connections to Cognitive Dynamics}


In the pure mathematics volume, the GaoZheng triple-layer connection
\GZTLC was introduced as a law-space level structure packaging three
coupled connections and curvature data on law-space bundles. In this
chapter we specialize \GZTLC to the hierarchical algebraic cognitive
architectures (HACA) of Definition~\ref{def:haca-law-objects-category}.
The goal is twofold. First, we explain how the symbolic, semantic, and
decision layers of a HACA law-object inherit restricted pieces of
\GZTLC, giving rise to a triple-layer connection that governs how local
updates propagate across layers. Second, we use this structure to define
semantic gauge fields and logical pressure fields, which constrain
GaoZheng Reinforcement Learning (GRL) in law-space and PACER dynamics
to follow ``geodesic-like'' trajectories in a suitable sense.

\section[Triple-Layer Connections for HACA Law-Objects]{Triple-Layer Connections for HACA Law-Objects}


In the pure mathematics volume, the GaoZheng triple-layer connection
\GZTLC on the law-space~\GZLS is given, schematically, by a triple
\[
  \GZTLC
  =
  \bigl(\GZLC^{(1)},\GZLC^{(2)},\GZLC^{(3)}\bigr),
\]
where each component is a connection on a law-space bundle, and the
three connections are coupled by compatibility relations and curvature
constraints. The precise construction is recalled there; for the
present chapter we only need the fact that \GZTLC yields three
distinguished connection components and associated curvature components
on any law-sector of~\GZLS.

Let $L$ be a HACA law-object in the sense of
Definition~\ref{def:haca-law-objects-category}, with law-support
$B_L\subseteq\GZLS$ and three-layer structure
\[
  L^{(\mathrm{sym})}
  \longrightarrow
  L^{(\mathrm{sem})}
  \longrightarrow
  L^{(\mathrm{dec})}
\]
as in Remark~\ref{rem:haca-three-layer-vs-ch-layer}. Restricting
\GZTLC to the sector $B_L$ and projecting to the three components
$L^{(\mathrm{sym})}$, $L^{(\mathrm{sem})}$, $L^{(\mathrm{dec})}$, we
obtain three induced connections which we denote by
\[
  {\GZLC}_L^{(\mathrm{sym})},
  \qquad
  {\GZLC}_L^{(\mathrm{sem})},
  \qquad
  {\GZLC}_L^{(\mathrm{dec})}.
\]
Their curvature tensors will be written
\[
  {\GZLCurv}_L^{(\mathrm{sym})},
  \qquad
  {\GZLCurv}_L^{(\mathrm{sem})},
  \qquad
  {\GZLCurv}_L^{(\mathrm{dec})},
\]
respectively.

\begin{definition}[Triple-layer connection on a HACA law-object]
\label{def:haca-triple-layer-connection}
Let $L$ be a HACA law-object with law-support $B_L\subseteq\GZLS$.
A \emph{triple-layer connection} on $L$ is the restriction of the
global triple-layer connection \GZTLC to $B_L$, viewed through the
three-layer structure of $L$:
\[
  \GZTLC_L
  :=
  \bigl(
    {\GZLC}_L^{(\mathrm{sym})},
    {\GZLC}_L^{(\mathrm{sem})},
    {\GZLC}_L^{(\mathrm{dec})}
  \bigr),
\]
together with their curvature tensors
${\GZLCurv}_L^{(\mathrm{sym})}$,
${\GZLCurv}_L^{(\mathrm{sem})}$,
${\GZLCurv}_L^{(\mathrm{dec})}$.
\end{definition}

Informally, ${\GZLC}_L^{(\mathrm{sym})}$ governs how symbolic
configurations (characters, tokens, and word-packs) can be transported
along law-space paths in $B_L$; ${\GZLC}_L^{(\mathrm{sem})}$ governs the
transport of semantic structures (paragraph graphs, KAT models, and
other elements of $\mathcal{S}(L)$); and ${\GZLC}_L^{(\mathrm{dec})}$
governs the transport of decision-level data (policies, value
baselines, and decision rules). The coupling of these three connections
enforces a compatibility between symbol-level, semantic-level, and
decision-level dynamics.

One may summarize the situation by the following schematic diagram:
\[
\begin{array}{ccccc}
  L^{(\mathrm{sym})} & \xrightarrow{\;\Phi^{(\mathrm{sym}\to\mathrm{sem})}\;}
                     & L^{(\mathrm{sem})} & \xrightarrow{\;\Phi^{(\mathrm{sem}\to\mathrm{dec})}\;}
                     & L^{(\mathrm{dec})} \\
  \Big\downarrow &   & \Big\downarrow &   & \Big\downarrow \\
  B_L            & = & B_L            & = & B_L
\end{array}
\]
where the vertical arrows denote the projections to the law-space
sector $B_L$ equipped with the restricted triple-layer connection
$\GZTLC_L$, and the horizontal arrows encode the layer-to-layer maps
within $L$.

\begin{remark}[Triple-layer connections and local cross-layer propagation]
\label{rem:haca-triple-layer-propagation}
Consider a law-space path $\gamma:[0,1]\to B_L$ together with three
lifts
\[
  \tilde{\gamma}^{(\mathrm{sym})},
  \quad
  \tilde{\gamma}^{(\mathrm{sem})},
  \quad
  \tilde{\gamma}^{(\mathrm{dec})}
\]
to $L^{(\mathrm{sym})}$, $L^{(\mathrm{sem})}$, and $L^{(\mathrm{dec})}$
which are horizontal with respect to
${\GZLC}_L^{(\mathrm{sym})}$,
${\GZLC}_L^{(\mathrm{sem})}$,
${\GZLC}_L^{(\mathrm{dec})}$.
A local update at the symbolic layer—e.g.\ replacing a phrase by an
equivalent one—can be viewed as a perturbation of
$\tilde{\gamma}^{(\mathrm{sym})}$. The coupling conditions in \GZTLC
then constrain how this perturbation must be propagated to the semantic
and decision layers, through corresponding perturbations of
$\tilde{\gamma}^{(\mathrm{sem})}$ and $\tilde{\gamma}^{(\mathrm{dec})}$,
so that the combined triple remains compatible with the law-space
geometry. In this sense, the triple-layer connection formalizes the
idea that ``small changes in wording induce controlled changes in
meaning and in decisions'' when viewed at the law-space level.
\end{remark}

\section[Semantic Gauge Fields and Logical Pressure Fields]{Semantic Gauge Fields and Logical Pressure Fields}


In Definition~\ref{def:semantic-gauge-field} of the previous chapter we
introduced semantic gauge fields on a HACA law-object $L$ in terms of
semantic functors $F_\alpha:\mathcal{S}(L)\to V_\alpha$ and a connection
$\mathcal{A}_L$ on the associated functor bundle over $B_L$. We now
relate this construction to the triple-layer connection
$\GZTLC_L$ of Definition~\ref{def:haca-triple-layer-connection}, and
introduce logical pressure fields that quantify how strongly semantic
trajectories are compressed, stretched, or twisted by the underlying
law-space curvature.

Let ${\GZLC}_L^{(\mathrm{sem})}$ be the semantic-layer component of the
triple-layer connection on $L$, with curvature
${\GZLCurv}_L^{(\mathrm{sem})}$. The semantic gauge field
$\mathcal{A}_L$ of Definition~\ref{def:semantic-gauge-field} can be
regarded as a refinement of ${\GZLC}_L^{(\mathrm{sem})}$ obtained by
pulling the law-space connection back along the semantic--law map
$\operatorname{Law}_L:\mathcal{S}(L)\to B_L$ and then pushing it
forward to the functor bundle of semantic evaluations:
\[
  \mathcal{A}_L
  \quad\text{encodes the parallel transport of the values }
  F_\alpha(\operatorname{Sem}_L(\omega))
  \text{ along law-space paths in }B_L.
\]
The curvature of this semantic gauge field will be denoted by
$\mathcal{F}_L$; it measures the failure of semantic evaluations to be
path-independent under parallel transport.

\begin{definition}[Logical pressure field]
\label{def:logical-pressure-field}
Let $L$ be a HACA law-object with semantic-layer connection
${\GZLC}_L^{(\mathrm{sem})}$, curvature ${\GZLCurv}_L^{(\mathrm{sem})}$,
and semantic gauge field $\mathcal{A}_L$ with curvature $\mathcal{F}_L$.
A \emph{logical pressure field} for $L$ is a function
\[
  \mathcal{F}_{\mathrm{log}} : B_L \longrightarrow \mathbb{R}_{\ge 0}
\]
obtained by evaluating the semantic curvature along a fixed family of
reference directions, i.e.\ there exists a family of law-space vectors
$\{v_x\in T_x B_L\}_{x\in B_L}$ and an evaluation functional
$\Psi$ such that
\[
  \mathcal{F}_{\mathrm{log}}(x)
  =
  \Psi\bigl({\GZLCurv}_L^{(\mathrm{sem})}(v_x,\cdot)\bigr)
  =
  \Psi\bigl(\mathcal{F}_L(v_x,\cdot)\bigr),
  \qquad x\in B_L.
\]
Intuitively, $\mathcal{F}_{\mathrm{log}}(x)$ measures the amount of
semantic ``pressure'' exerted by the law-space curvature at $x$, in the
sense of how sensitive semantic evaluations are to small loops based at
$x$.
\end{definition}

The logical pressure field interacts with differential dynamical quanta
(MDQ) in the following way. Let $\gamma:[0,1]\to B_L$ be a law-space
path and let $\omega_t\in\Omega(L)$ be an operator path lifting
$\gamma$, so that $\Phi_L(\omega_t)=\gamma(t)$ for all $t$. In
Applications Volume~I, the MDQ density $q_{\mathrm{MDQ}}(t)$ along
$\gamma$ is a local law-space quantity encoding how the PFB-GNLA
geometry drives GRL dynamics at the physical or application level. In
the present volume, semantic functors $F_\alpha$ define an analogous
semantic MDQ density $q_{\mathrm{MDQ}}^{\mathrm{sem}}(t)$ by measuring
the infinitesimal change of semantic evaluations along $\gamma$:
\[
  q_{\mathrm{MDQ}}^{\mathrm{sem}}(t)
  \;\sim\;
  \bigl\|
    \nabla_{\dot{\gamma}(t)} F_\alpha
  \bigr\|
  \quad\text{for a suitable family of }F_\alpha.
\]

\begin{remark}[Semantic MDQ as a law-space pushforward]
\label{rem:semantic-mdq-pushforward}
The construction of \(q_{\mathrm{MDQ}}^{\mathrm{sem}}\) is by design a
law-space pushforward of the MDQ density developed in Applications
Volume~I~\cite{GaoZhengAppVol1}: the semantic MDQ is obtained by
restricting the MDQ functional to HACA law-objects and admissible
semantic functors and then evaluating it along semantic trajectories.
Consequently, whenever the hypotheses of Applications Volume~I ensure
existence, integrability, or stability properties for the MDQ along GRL
trajectories, the same analytical properties hold for
\(q_{\mathrm{MDQ}}^{\mathrm{sem}}\) in the semantic sector under the
corresponding semantic regularity assumptions.
\end{remark}

\begin{proposition}[Curvature control of semantic MDQ]
\label{prop:logical-pressure-mdq-control}
Under suitable regularity and boundedness assumptions on the semantic
functors and on the law-space connection, there exist constants
$C_1,C_2\ge 0$ such that, for every HACA law-object $L$, every
law-space path $\gamma:[0,1]\to B_L$, and every admissible operator
lift $\omega_t$, the semantic MDQ density satisfies the pointwise
estimate
\[
  \bigl|q_{\mathrm{MDQ}}^{\mathrm{sem}}(t)\bigr|
  \;\le\;
  C_1 + C_2\,\mathcal{F}_{\mathrm{log}}\bigl(\gamma(t)\bigr)
\]
for almost every $t\in[0,1]$.
\end{proposition}

A detailed proof of
Proposition~\ref{prop:logical-pressure-mdq-control}
is given in Section~\ref{app:sec-logical-pressure-mdq-proof} of the
Technical Appendix. Intuitively, this estimate can be viewed as an
adaptation of standard curvature-control arguments in gauge theory and
GRL: the covariant derivative of a functor along $\gamma$ is bounded
by a combination of its value and the curvature of the connection; the
logical pressure field extracts the relevant curvature component and
thus provides an a priori bound on the magnitude of semantic MDQ. In
particular, regions where $\mathcal{F}_{\mathrm{log}}$ is small behave
like ``low-pressure'' semantic regimes in which semantic trajectories
can evolve with small local cost, whereas regions with large
$\mathcal{F}_{\mathrm{log}}$ correspond to high-pressure regimes in
which semantic changes are expensive and strongly constrained.

\begin{remark}[Semantic geodesics and law-space guidance]
\label{rem:semantic-geodesics}
The logical pressure field suggests a natural notion of
``semantic geodesics'' for PACER and GRL in law-space. Given a
Lagrangian combining semantic MDQ and logical pressure,
\[
  \mathcal{L}(\gamma,\dot{\gamma})
  =
  \bigl|q_{\mathrm{MDQ}}^{\mathrm{sem}}(t)\bigr|
  + \lambda\,\mathcal{F}_{\mathrm{log}}(\gamma(t)),
\]
with a coupling parameter $\lambda\ge 0$, one may define semantic
geodesics as critical points of the corresponding action functional.
Trajectories that minimize this action balance the desire to keep
semantic MDQ small with the requirement to avoid or carefully traverse
regions of high logical pressure. In later chapters, PACER and
SAC-PACER-HACA engines will be interpreted as learning mechanisms that
approximate such geodesic-like trajectories in the HACA law-space.
\end{remark}

\section[CH-Layering, Safety Layers, and Higher-Order Constraints]{CH-Layering, Safety Layers, and Higher-Order Constraints}


The CH-layering viewpoint developed in the pure mathematics volume
assigns to the Continuum Hypothesis and its relatives a layered status
within law-space: the layers
\[
  \Layer_2,\ \Layer_3,\ \Layer_4,\ \dots
\]
represent increasingly rich structures and constraints built on top of
the same underlying law-space geometry. In the HACA/PACER context, this
layered perspective provides a natural way to organize semantic,
safety, and higher-order constraints.

At the $\Layer_2$-level, the focus is on meaning and truth: the
property layer $\mathrm{Prop}(L)$ of
Definition~\ref{def:haca-property-layer} and the observable partitions
of Definition~\ref{def:haca-observable-properties} resolve the raw
operator universe $\Omega(L)$ into semantically meaningful sectors. The
basic question at this level is whether an operator configuration or
trajectory is semantically correct with respect to a family of
predicates and semantic functors. The logical pressure field
$\mathcal{F}_{\mathrm{log}}$ enters here as a quantitative measure of
how tightly semantic correctness is constrained by the law-space
curvature.

At the $\Layer_3$-level, homotopy and Jacobiator-gaps are added on
top of the $\Layer_2$ structure. Safety and robustness are no longer
formulated purely as pointwise semantic conditions along a single path,
but as conditions stable under controlled homotopies and subject to
Jacobiator-gap constraints. In particular, acceptable HACA/PACER
trajectories must remain in the same controlled homotopy class (LA/LB/LC)
as a certified safe baseline, and property intersections must remain
non-empty along homotopy-smoothing operations as in
Proposition~\ref{prop:gaugeable-operator-bundle-from-property-intersections}.
The triple-layer connection and the semantic gauge field ensure that
safety constraints are respected simultaneously at the symbolic,
semantic, and decision layers.

The $\Layer_4$-level and higher layers are designed to capture
metaphor, analogy, and additional forms of higher-order structure. At
these levels, structure-preserving functors between different HACA
law-objects and law-sectors encode cross-domain connections: safe
patterns of reasoning in one domain can be transported, under suitable
constraints, to other domains. Higher layers can also incorporate
value geometry, ethical constraints, and multi-agent game-theoretic
conditions, all expressed as additional structure compatible with the
underlying triple-layer connection and semantic gauge fields.

\begin{definition}[Safe and unsafe law-space sectors for HACA/PACER]
\label{def:safe-unsafe-sectors}
Let $L$ be a HACA law-object with triple-layer connection
$\GZTLC_L$, semantic gauge field $\mathcal{A}_L$, logical pressure
field $\mathcal{F}_{\mathrm{log}}$, and homotopy-completeness profile
(LA/LB/LC) as in Definition~\ref{def:homotopy-completeness-levels}.
A subset $S\subseteq B_L$ is called a \emph{safe law-space sector} if:
\begin{enumerate}
  \item For every law-space path $\gamma:[0,1]\to S$ and every
        admissible HACA/PACER trajectory lifting $\gamma$, there exists
        a certified safe baseline trajectory in the same homotopy class
        (with respect to the $\Layer_3$ structure) and with MDQ and
        logical pressure bounded by prescribed thresholds.
  \item The restrictions of $\GZTLC_L$ and $\mathcal{A}_L$ to $S$ admit
        a choice of gauges in which the logical pressure field
        $\mathcal{F}_{\mathrm{log}}$ remains uniformly bounded and
        Jacobiator-gap data remain homotopy-fillable in the sense of
        Proposition~\ref{prop:gaugeable-operator-bundle-from-property-intersections}.
\end{enumerate}
The complement $B_L\setminus S$ is called an \emph{unsafe sector}
(relative to the chosen thresholds and baseline). In subsequent
chapters, PACER and SAC-PACER-HACA engines will be designed to operate
predominantly inside safe sectors, with explicit mechanisms to detect
and avoid or quarantine excursions into unsafe regions.
\end{definition}

\begin{remark}[Pre-compiled safety vs.\ ad hoc patching]
\label{rem:precompiled-safety}
From the CH-layering point of view, safety in HACA/PACER is not an
afterthought implemented by ad hoc patching or post hoc filtering.
Instead, safety is pre-compiled into the law-space structure: the
$\Layer_2$ property layer ensures that only semantically meaningful
sectors are considered; the $\Layer_3$ homotopy and Jacobiator-gap
structure ensures that safe behavior is stable under controlled
deformations; and higher layers encode cross-domain and value-related
constraints. The triple-layer connection and semantic gauge fields
provide the geometric infrastructure for these layers, and the notion
of safe/unsafe sectors in Definition~\ref{def:safe-unsafe-sectors}
formalizes the idea that PACER and SAC-PACER-HACA engines should run
along law-space trajectories that are geometrically compatible with
pre-established safety envelopes.
\end{remark}

\chapter[Law-Space Baselines for Cognitive Agents and Pack-Level Value Geometry]{Law-Space Baselines for Cognitive Agents and Pack-Level Value Geometry:\\
         MDQ and Value Geometry in HACA}
In the Applications Volume~I, baselines and occupancy measures were used
to describe how GaoZheng Reinforcement Learning (GRL) in law-space
distributes its attention and resources across different sectors of a
law-space, and how small perturbations of these distributions change the
value of a given objective. In the present volume we develop the
analogous notions for HACA-based cognitive agents: law-space baselines
for ``neutral'' cognitive behavior, pack-level occupancy measures for
semantic attention and usage, and value geometry vectors that encode how
the objective reacts to redistributions of attention across packs. These
constructions provide the geometric background for PACER and
SAC-PACER-HACA engines in later chapters.

\section[Law-Space Baselines for HACA Agents]{Law-Space Baselines for HACA Agents}


In GRL for physical and engineering systems, a baseline configuration
$(x^{\mathrm{base}},\mu^{\mathrm{base}})$ is often used to represent a
neutral or default regime: an equilibrium state, a reference control
policy, or a default operating mode. In the HACA setting, the analogous
notion is a baseline HACA law-object representing a neutral reading or
answering behavior in law-space.

\begin{definition}[Baseline law-object for a HACA agent]
\label{def:haca-baseline-law-object}
A \emph{baseline law-object} for a HACA agent is a HACA law-object
$L^{\mathrm{base}}$ with law-support
$B_{L^{\mathrm{base}}}\subseteq\GZLS$ together with a choice of
triple-layer connection
\[
  \GZTLC_{L^{\mathrm{base}}}
  =
  \bigl(
    {\GZLC}_{L^{\mathrm{base}}}^{(\mathrm{sym})},
    {\GZLC}_{L^{\mathrm{base}}}^{(\mathrm{sem})},
    {\GZLC}_{L^{\mathrm{base}}}^{(\mathrm{dec})}
  \bigr)
\]
such that:
\begin{enumerate}
  \item The semantic layer $L^{(\mathrm{sem})}_{\mathrm{base}}$ encodes
        neutral or default reading/answering behavior: for instance, the
        behavior of reading a document and restating it in a faithful,
        non-optimizing way without additional goals such as persuasion,
        compression, or adversarial obfuscation.
  \item The decision layer
        $L^{(\mathrm{dec})}_{\mathrm{base}}$ represents a neutral,
        non-optimizing policy, for example a policy that selects
        responses according to a fixed, law-space–compatible prior
        without attempting to maximize or minimize any explicit
        objective beyond basic semantic correctness and coherence.
  \item The curvature tensors
        ${\GZLCurv}_{L^{\mathrm{base}}}^{(\mathrm{sym})}$,
        ${\GZLCurv}_{L^{\mathrm{base}}}^{(\mathrm{sem})}$,
        ${\GZLCurv}_{L^{\mathrm{base}}}^{(\mathrm{dec})}$
        are bounded in a way that reflects the intended ``no additional
        pressure'' property: they do not encode any extra structural
        preference beyond those required by the global PFB-GNLA
        background and minimal semantic correctness.
\end{enumerate}
We write
\[
  \GZLC^{\mathrm{base}}
  :=
  {\GZLC}_{L^{\mathrm{base}}}^{(\mathrm{sem})},
  \qquad
  \GZLCurv^{\mathrm{base}}
  :=
  {\GZLCurv}_{L^{\mathrm{base}}}^{(\mathrm{sem})},
\]
and call $(B_{L^{\mathrm{base}}},\GZLC^{\mathrm{base}},
\GZLCurv^{\mathrm{base}})$ the \emph{baseline semantic geometry} for
the agent.
\end{definition}

Given a concrete HACA law-object $L$ modeling a cognitive agent with
nontrivial preferences and objectives (for instance, a PACER engine with
specific goals and constraints), we can measure the deviation of its
law-space geometry from the baseline geometry of
$L^{\mathrm{base}}$ by comparing the semantic-layer connections and
curvatures.

\begin{definition}[Deviation from baseline geometry]
\label{def:baseline-deviation}
Let $L$ be a HACA law-object with law-support $B_L\subseteq\GZLS$,
semantic-layer connection ${\GZLC}_L^{(\mathrm{sem})}$, and curvature
${\GZLCurv}_L^{(\mathrm{sem})}$, and let $L^{\mathrm{base}}$ be a
baseline law-object as in
Definition~\ref{def:haca-baseline-law-object}. Assume that the
law-supports satisfy $B_L\subseteq B_{L^{\mathrm{base}}}$ (or we have
chosen an embedding of $B_L$ into $B_{L^{\mathrm{base}}}$). The
\emph{baseline deviation} of $L$ is the pair of tensors
\[
  \Delta{\GZLC}_L
  :=
  {\GZLC}_L^{(\mathrm{sem})} - \GZLC^{\mathrm{base}},
  \qquad
  \Delta{\GZLCurv}_L
  :=
  {\GZLCurv}_L^{(\mathrm{sem})} - \GZLCurv^{\mathrm{base}},
\]
defined pointwise on $B_L$ by subtracting the baseline connection and
curvature from those of $L$.
\end{definition}

In the subsequent sections, MDQ and value geometry for HACA agents will
be expressed in terms of these deviation tensors: semantic MDQ measure
how expensive it is to follow a law-space path with respect to the
deviant geometry, and value gradients encode how desirable it is, from
the standpoint of a given objective, to increase or decrease attention
to different parts of this deviated geometry. In this sense, all
quantities of interest are measured relative to the baseline geometry
$(\GZLC^{\mathrm{base}},\GZLCurv^{\mathrm{base}})$ rather than in an
absolute vacuum.

\section[Pack-Level Occupancy Measures and Value Geometry]{Pack-Level Occupancy Measures and Value Geometry}


In GRL, occupancy measures record how much time or probability mass a
policy spends in different regions of state–action space. For HACA
agents, we are interested in how much ``semantic attention'' a policy
devotes to different packs: word-packs, sentence segments, or paragraph
nodes in the operator universe $\Omega(L)$ of
Definition~\ref{def:haca-operator-universe}.

Let $L$ be a HACA law-object and let $\Pi(L)$ denote a class of
admissible policies for the agent (for instance, PACER or
SAC-PACER-HACA policies) defined on $L$. Fix a finite index set
\[
  I_L = \{1,\dots,N\}
\]
enumerating a collection of representative packs\footnote{For instance,
one can take $I_L$ to index a finite set of paragraph nodes, sentence
templates, or abstract pack patterns.} in $\Omega(L)$.

\begin{definition}[Pack-level occupancy measures]
\label{def:pack-occupancy-measures}
For a policy $\pi\in\Pi(L)$, the \emph{pack-level occupancy measure} is
a vector
\[
  \mu^\pi
  =
  (\mu_i^\pi)_{i\in I_L}
  \in\mathbb{R}_{\ge 0}^N
\]
where $\mu_i^\pi$ is the long-run fraction of semantic attention,
time, or usage frequency devoted to pack~$i$ under the policy~$\pi$.
Concretely, if $i$ indexes a measurable subset $P_i\subseteq\Omega(L)$
of operator configurations, one may set
\[
  \mu_i^\pi
  :=
  \lim_{T\to\infty}\frac{1}{T}\int_0^T
    \mathbf{1}_{\{\omega_t\in P_i\}}\,\mathrm{d}t,
\]
where $\omega_t$ is the operator trajectory induced by running $\pi$
over time and projecting to~$\Omega(L)$. When the underlying dynamics
is discrete, the integral is replaced by an average over time steps.
\end{definition}

An objective functional for a HACA agent is a map
\[
  \mathcal{J} : \Pi(L)\longrightarrow \mathbb{R}
\]
assigning to each policy $\pi$ a real-valued performance score
(which may encode a mixture of utility, safety, and other criteria).
Under mild regularity assumptions, $\mathcal{J}$ can be viewed as a
function of the occupancy measure $\mu^\pi$ together with additional
summary statistics of the policy, and one can consider its gradient
with respect to~$\mu^\pi$.

\begin{definition}[Pack-level value geometry vector]
\label{def:pack-value-geometry}
Assume that, for policies in a neighborhood of a reference policy
$\pi_0$, the functional $\mathcal{J}$ can be written as a differentiable
function of the occupancy measure,
\[
  \mathcal{J}(\pi)
  =
  \widehat{\mathcal{J}}(\mu^\pi)
  \quad\text{for }\pi\approx\pi_0.
\]
The \emph{value geometry vector} (or \emph{value baseline gradient
vector}) at~$\pi_0$ is the vector
\[
  \mathbf{v}
  =
  (v_i)_{i\in I_L},
  \qquad
  v_i
  :=
  \left.
    \frac{\partial \widehat{\mathcal{J}}}{\partial \mu_i}
  \right|_{\mu=\mu^{\pi_0}}.
\]
The component $v_i$ measures the local sensitivity of the objective to
an infinitesimal increase of occupancy at pack~$i$, holding other
components fixed.
\end{definition}

Geometrically, one may regard $\mu^\pi$ as a point in the simplex
\[
  \Delta^{N-1}
  =
  \Bigl\{
    \mu\in\mathbb{R}_{\ge 0}^N\;\Big|\;
    \sum_{i\in I_L}\mu_i = 1
  \Bigr\}
\]
(after suitable normalization), and $\mathbf{v}$ as a covector
assigning to each direction of movement in this simplex a first-order
change in~$\mathcal{J}$. When combined with the law-space geometry via
the map
\[
  \Phi_L : \Omega(L)\longrightarrow B_L\subseteq\GZLS,
\]
the occupancy measure and value vector together define a law-space–aware
value geometry: the policy does not merely allocate attention across
packs in a purely statistical sense, but across law-space sectors, and
$\mathbf{v}$ indicates which sectors are more or less desirable from
the standpoint of~$\mathcal{J}$.

\section[MDQ, Semantic Functors, and Baseline Constraints]{MDQ, Semantic Functors, and Baseline Constraints}


In the Applications Volume~I, differential dynamical quanta (MDQ)
were introduced as local quanta of law-aware dynamics along GRL path
integrals: at each point of a law-space path, the MDQ encode the
infinitesimal contribution of the PFB-GNLA data to the evolution of
observables. In Remark~\ref{rem:semantic-functors-vs-mdq} of the
present volume, semantic functors were identified as the semantic-sector
counterparts of MDQ: they assign to each operator configuration a
semantic evaluation that behaves as a differential quantum of meaning
along law-space paths.

Let $L$ be a HACA law-object equipped with a semantic gauge field and
baseline geometry as in
Definitions~\ref{def:haca-baseline-law-object}
and~\ref{def:semantic-gauge-field}. For each pack index $i\in I_L$,
a family of semantic functors $\{F_\alpha\}_\alpha$ and a choice of
law-space path ensemble under policy~$\pi$ determine a mean semantic MDQ
value $q_i^{\mathrm{sem}}(\pi)$ associated with pack~$i$: loosely
speaking, this is the average per-unit-time contribution of pack~$i$ to
the semantic MDQ density along the law-space trajectories induced by
$\pi$.

\begin{definition}[Semantic MDQ profile and joint cost functional]
\label{def:semantic-mdq-profile}
For a policy $\pi\in\Pi(L)$, the \emph{semantic MDQ profile} is the
vector
\[
  \mathbf{q}^{\mathrm{sem}}(\pi)
  =
  \bigl(q_i^{\mathrm{sem}}(\pi)\bigr)_{i\in I_L}
  \in\mathbb{R}_{\ge 0}^N,
\]
where $q_i^{\mathrm{sem}}(\pi)$ is the mean semantic MDQ assigned to
pack~$i$ under~$\pi$ as described above. Given a pack-level occupancy
vector $\mu^\pi$ and a value geometry vector $\mathbf{v}$ as in
Definition~\ref{def:pack-value-geometry}, we define a
\emph{joint semantic–value cost functional} by
\[
  \mathcal{J}_{\mathrm{joint}}(\pi)
  :=
  \sum_{i\in I_L}
    \mu_i^\pi\,q_i^{\mathrm{sem}}(\pi)
  \;-\;
  \lambda\,\sum_{i\in I_L}\mu_i^\pi\,v_i,
\]
where $\lambda\ge 0$ is a coupling parameter.
\end{definition}

The first term in $\mathcal{J}_{\mathrm{joint}}$ penalizes policies
that allocate attention to packs with large semantic MDQ contributions;
the second term rewards allocations that increase the objective
$\mathcal{J}$ as encoded by the value gradient~$\mathbf{v}$. The
baseline geometry enters through the definition of the MDQ profile: the
semantic MDQ values are computed relative to the baseline connection
and curvature, so that deviations from the neutral reading/answering
behavior are explicitly accounted for.

\begin{definition}[Baseline-constrained admissible region for MDQ]
\label{def:baseline-constrained-mdq-region}
Let $L$ be a HACA law-object with baseline deviation
$(\Delta{\GZLC}_L,\Delta{\GZLCurv}_L)$ as in
Definition~\ref{def:baseline-deviation}. A vector
$\mathbf{q}\in\mathbb{R}_{\ge 0}^N$ is called \emph{baseline-admissible}
if there exists a policy $\pi\in\Pi(L)$ such that:
\begin{enumerate}
  \item $\mathbf{q}^{\mathrm{sem}}(\pi)\le\mathbf{q}$ componentwise;
  \item the law-space trajectories induced by $\pi$ remain within safe
        sectors in the sense of
        Definition~\ref{def:safe-unsafe-sectors};
  \item the deviation tensors satisfy prescribed bounds along these
        trajectories, e.g.
        \[
          \bigl\|\Delta{\GZLC}_L\bigr\|\le C_{\mathrm{conn}},
          \qquad
          \bigl\|\Delta{\GZLCurv}_L\bigr\|\le C_{\mathrm{curv}}
        \]
        for given constants
        $C_{\mathrm{conn}},C_{\mathrm{curv}}\ge 0$.
\end{enumerate}
The set of all baseline-admissible MDQ profiles is denoted by
\[
  \mathcal{Q}_{\mathrm{adm}}(L)
  \subseteq \mathbb{R}_{\ge 0}^N.
\]
\end{definition}

From this point of view, the design problem for HACA agents equipped
with PACER or SAC-PACER-HACA engines can be phrased as: choose policies
$\pi$ whose occupancy measures $\mu^\pi$, value geometry vectors
$\mathbf{v}$, and semantic MDQ profiles
$\mathbf{q}^{\mathrm{sem}}(\pi)$ lie in desirable regions of the joint
space
\[
  \bigl\{(\mu,\mathbf{v},\mathbf{q})\bigr\}
\]
while respecting the baseline-admissible constraints encoded by
$\mathcal{Q}_{\mathrm{adm}}(L)$. The baseline geometry ensures that
these constraints are not arbitrary, but grounded in the underlying
law-space structure.

\section[Cross-Volume Alignment and Example Scenarios]{Cross-Volume Alignment and Example Scenarios}
The law-space baseline and value geometry constructions in this chapter
are deliberately parallel to those in the previous volumes. This
alignment allows us to reinterpret the domain-specific constructions of
Applications Volumes~I and~II in the HACA/PACER language and to import
their intuition into the semantic setting.

In the Applications Volume~I (physics and GRL), occupancy measures
describe how much time a policy spends in different regions of a
physical or abstract state–action space, and the value gradient
$\mathbf{v}$ indicates how desirable it is to increase or decrease
occupancy of those regions. MDQ quantify the local contribution of the
PFB-GNLA geometry to the evolution of observables along GRL path
integrals.

In the life-science–oriented Applications Volume~II (LBOPB), similar
structures appear in a different guise: occupancy measures quantify how
much time or probability mass is spent in different disease regimes,
drug regimens, or pharmacokinetic/pharmacodynamic (PK/PD)
configurations, while value gradients encode improvements or trade-offs
in efficacy, toxicity, and related criteria. MDQ control how the LBOPB
law-space geometry affects the evolution of pharmacological observables
and pharmacogenomics monoids (PGOM).

In the present HACA volume (Applications Volume~III), packs replace
physical states or biological regimes as the basic units of analysis:
occupancy measures $(\mu_i^\pi)$ describe how much semantic attention a
policy devotes to different packs, value gradients $(v_i)$ describe
which packs are more or less desirable from the standpoint of
$\mathcal{J}$, and semantic MDQ $(q_i^{\mathrm{sem}})$ describe how much
``semantic dynamical cost'' is associated with those packs. The
structural pattern is the same:
\[
\begin{array}{lll}
  \text{App.Vol.1:} &
    \text{states/paths in law-space for GRL-based physics} &
    \longrightarrow (\mu_i, v_i, q_i) \\[0.3em]
  \text{App.Vol.2 (LBOPB):} &
    \text{LBOPB law-objects (pathology, PK/PD, PGOM)} &
    \longrightarrow (\mu_i, v_i, q_i) \\[0.3em]
  \text{App.Vol.3 (HACA):} &
    \text{HACA packs and reading paths} &
    \longrightarrow (\mu_i, v_i, q_i^{\mathrm{sem}})
\end{array}
\]
The ambient interpretation of the indices $i$ changes from physical
states in Applications Volume~I, to LBOPB regimes in Applications
Volume~II, to HACA packs in Applications Volume~III, but the
mathematical structure (occupancy measures and value gradients defined
over law-space sectors) remains the same.

\begin{remark}[Additional construction difficulty for HACA vs.\ LBOPB]
\label{rem:haca-vs-lbopb-difficulty}
From the viewpoint of law-space geometry, the additional construction
difficulty of the HACA volume, compared to the LBOPB-based
Applications Volume~II, can be summarized as follows. In the LBOPB
setting, heterogeneous evolution is primarily realized as law-space
trajectories moving between different LBOPB law-objects (pathology
regimes, PK/PD configurations, PGOM sectors), while the underlying GRL
dynamics takes place on a single life-science law-space with a fixed
syntactic representation.

In the HACA setting, by contrast, heterogeneous evolution is driven
simultaneously at the KAT-based syntactic layer ($l_1$) and at the
higher semantic layers $l_n$ for $n\ge 2$. Full-document reading and
iterative paragraph updates are modeled as long KAT paths in the
operator universe $\Omega(L)$, which project via
$\operatorname{Sem}_L$ and $\operatorname{Law}_L$ into higher layers of
the law-space stratification. The triple-layer connection $\GZTLC_L$
and the associated semantic gauge fields must therefore propagate
local updates consistently across symbolic, semantic, and decision
layers, while respecting property-layer constraints and homotopy-level
safety conditions ($l_2$ for meaning, $l_3$ for safety and robustness,
and higher $l_n$ for metaphor and cross-domain structure).

In this sense, the HACA construction is not merely ``another''
application of the PFB-GNLA background, but one in which KAT syntax and
multi-layer semantic structure are tightly coupled by law-space
connections, making heterogeneous evolution (across domains, textual
segments, and reasoning modes) a first-class geometric object.
\end{remark}

As a concrete example, consider a cross-volume pipeline:
\begin{itemize}
  \item In Applications Volume~II (LBOPB), a law-space trajectory
        connects a pathological law-object to a reverse-design drug
        candidate via an LBOPB pipeline, with occupancy measures
        reflecting time spent in different PK/PD regimes and value
        gradients encoding efficacy and safety trade-offs.
  \item In Applications Volume~III (HACA), the same pipeline is
        represented at the HACA level as a sequence of packs:
        descriptions of the pathology, mechanisms of action, candidate
        structures, and trial designs; the HACA occupancy measures
        record how much attention a cognitive agent devotes to each of
        these packs, and the value gradients encode which packs
        contribute more to the success of the overall argument or
        decision.
\end{itemize}
Jacobiator-gap certificates and homotopy completeness, developed in the
next chapters, ensure that these cross-volume representations are not
merely informal analogies but are linked by controlled homotopies in
law-space: safe and effective pipelines in LBOPB can be lifted to safe
and well-structured reasoning pipelines in HACA/PACER, and vice versa,
all within a unified law-space geometric framework.



\part[Lexical KAT Action Monoids and the PACER Architecture]{Lexical KAT Action Monoids and the PACER Architecture\\
      for Pack-Aligned Compressive-Expansion Reasoning}

\chapter[Lexical KAT Action Monoids and Word-Pack Operators]{Lexical KAT Action Monoids and Word-Pack Operators:\\
         From Base Packs to Power Submonoid Spectra}
In Part~I we equipped HACA law-objects with multi-scale operator
universes, semantic structure spaces, and semantic gauge fields over
law-space. Part~II specializes these constructions to the lexical layer:
we isolate the free monoid of lexical sequences, describe endomorphism
monoids acting on it, and then organize the resulting KAT-based action
monoids into spectra of submonoids with distinct roles (memory,
prediction, termination, compliance). These structures provide the
syntactic backbone on which PACER operates: its compressive–alignment–
expansion–recompression cycle is realized as a controlled traversal of
KAT action monoids and their submonoid spectra.

\section[Free and Endomorphism Monoids over Lexical Sequences]{Free and Endomorphism Monoids over Lexical Sequences}


On the symbolic layer of a HACA law-object $L$, lexical data are
represented as sequences of tokens drawn from a finite alphabet
$\Sigma$. The natural algebraic structure on such sequences is the
free monoid.

\begin{definition}[Free monoid of lexical sequences]
\label{def:free-monoid-lexical}
Let $\Sigma$ be a finite alphabet (for instance, a set of characters,
subword units, or token identifiers). The \emph{free monoid} on
$\Sigma$ is the set
\[
  \Sigma^*
  :=
  \{\, w = a_1 a_2\cdots a_n \mid n\ge 0,\ a_i\in\Sigma \,\},
\]
equipped with concatenation
\[
  \circ : \Sigma^*\times\Sigma^*\\longrightarrow\Sigma^*,
  \qquad
  (u,v)\longmapsto u\circ v,
\]
and the empty word $\varepsilon$ as unit. Thus $(\Sigma^*,\circ,
\varepsilon)$ is a monoid: concatenation is associative and
$\varepsilon$ is a two-sided identity element.
\end{definition}

A HACA law-object $L$ has a symbolic layer $L^{(\mathrm{sym})}$ which
includes such a free monoid $(\Sigma^*,\circ,\varepsilon)$, together
with additional structure (e.g.\ lexical categories, token types,
positions in documents). To describe transformations of lexical data,
we consider monoids of endomorphisms of $\Sigma^*$.

\begin{definition}[Endomorphism monoid over lexical sequences]
\label{def:endomorphism-monoid}
The \emph{endomorphism monoid} of the free monoid $(\Sigma^*,\circ,
\varepsilon)$ is the set
\[
  \mathrm{End}(\Sigma^*)
  :=
  \{\, f:\Sigma^*\to\Sigma^* \,\big|\, f \text{ is a function} \,\},
\]
equipped with composition as multiplication,
\[
  (f,g)\longmapsto g\circ f,
\]
and the identity map $\mathrm{id}_{\Sigma^*}$ as unit. Thus
$(\mathrm{End}(\Sigma^*),\circ,\mathrm{id}_{\Sigma^*})$ is a monoid.
Elements of $\mathrm{End}(\Sigma^*)$ will be called \emph{lexical
operators}.
\end{definition}

\begin{example}[Simple lexical operators for $\Sigma=\{a,b\}$]
\label{ex:lexical-operators-ab}
Let $\Sigma=\{a,b\}$ and consider the following functions on $\Sigma^*$:
\begin{itemize}
  \item $T_a$: prefixing $a$ to a word,
        \(
          T_a(w) := a w;
        \)
  \item $D_b$: deleting all occurrences of $b$,
        \(
          D_b(w) := \text{the word obtained from $w$ by removing all $b$'s};
        \)
  \item $S_{ab}$: swapping $a$ and $b$,
        \(
          S_{ab}(a) = b,\ S_{ab}(b) = a,
        \)
        extended homomorphically to all of $\Sigma^*$.
\end{itemize}
Each of these maps lies in $\mathrm{End}(\Sigma^*)$, and compositions
such as $D_b\circ T_a$ or $S_{ab}\circ D_b$ illustrate how lexical
operators can be combined into more complex transformations.
\end{example}

In the HACA setting, we are not interested in all elements of
$\mathrm{End}(\Sigma^*)$, but in those that are compatible with the
semantic and law-space structures of $L$. The lexical KAT action
monoids introduced below will be defined as structured submonoids of
$\mathrm{End}(\Sigma^*)$ which preserve the relevant HACA structure.

\section[Lexical KAT Action Monoids and Base Packs]{Lexical KAT Action Monoids and Base Packs}


Kleene algebra with tests (KAT)~\cite{KozenSmithKAT1996} provides an algebraic framework for
reasoning about control flow and tests. In the present volume, we use
KAT not as a full-blown semantic theory, but as a convenient language
for encoding structured operations on lexical sequences and their
associated control conditions.

\begin{definition}[Lexical KAT action monoid]
\label{def:lex-kat-action-monoid}
Let $(\Sigma^*,\circ,\varepsilon)$ be the free monoid of lexical
sequences for a HACA law-object $L$, and let $\mathrm{End}(\Sigma^*)$
be its endomorphism monoid. A \emph{lexical KAT action monoid} on $L$
is a submonoid
\[
  \mathcal{M} \subseteq \mathrm{End}(\Sigma^*)
\]
generated by a finite or countable family of \emph{primitive lexical
operators} and \emph{tests}, together with an interpretation into a
KAT model $(K,T)$ such that:
\begin{enumerate}
  \item Each primitive operator $u\in\mathcal{M}$ is associated with an
        element $[u]\in K$, and each test is associated with an element
        $[p]\in T\subseteq K$, in a way that respects composition and
        KAT operations.
  \item The action of $\mathcal{M}$ on $\Sigma^*$, given by
        \[
          (\mathcal{M},\Sigma^*)\ni (u,w)
          \longmapsto u(w)\in\Sigma^*,
        \]
        is compatible with the HACA symbolic layer $L^{(\mathrm{sym})}$
        and the semantic projection $\operatorname{Sem}_L$ of
        Definition~\ref{def:haca-sem-law-projections}: applying $u$ to
        $w$ induces a well-defined transformation on the corresponding
        semantic structures and law-space points.
\end{enumerate}
\end{definition}

To describe the basic building blocks of $\mathcal{M}$, it is useful
to distinguish different levels of word-packs.

\begin{definition}[Base packs at multiple levels]
\label{def:base-packs}
Let $L$ be a HACA law-object with operator universe $\Omega(L)$ as in
Definition~\ref{def:haca-operator-universe}. A \emph{base pack} for
$L$ is a lexical operator $B\in\mathcal{M}$ that is designated as a
normalized representation of a semantically coherent unit at one of
the following levels:
\begin{itemize}
  \item \emph{Token-level base packs}: operators that act on individual
        tokens or very short token sequences (e.g.\ normalizing
        punctuation, case, or simple token substitutions).
  \item \emph{Phrase-level base packs}: operators that act on short
        phrases or lexical patterns (e.g.\ inserting or deleting a
        stock phrase, reordering modifiers, or expanding an idiom into
        an explicit phrase).
  \item \emph{Sentence-level base packs}: operators that act on whole
        sentences or sentence templates (e.g.\ replacing a sentence by
        an equivalent one with different word order, or inserting a
        standard disclaimer sentence).
\end{itemize}
We write $\mathcal{B}^{(0)}(L)$, $\mathcal{B}^{(1)}(L)$,
$\mathcal{B}^{(2)}(L)$ for the sets of token-, phrase-, and
sentence-level base packs, respectively, and
\[
  \mathcal{B}(L)
  :=
  \mathcal{B}^{(0)}(L)\cup\mathcal{B}^{(1)}(L)\cup\mathcal{B}^{(2)}(L)
  \subseteq \mathcal{M}
\]
for their union.
\end{definition}

\begin{example}[Sentence-level base pack as reorder/insert operator]
\label{ex:base-pack-sentence-reorder}
Consider a sentence of the form
\[
  w = \text{``If condition, then statement.''}
\]
A sentence-level base pack $B\in\mathcal{B}^{(2)}(L)$ may implement a
normalized reordering such as
\[
  B(w)
  =
  \text{``Statement holds if condition.''},
\]
possibly inserting or removing punctuation to fit the target style.
The corresponding KAT element $[B]\in K$ encodes both the control-flow
structure (condition–then) and the normalized lexical form. Applying
$B$ in the HACA setting induces a controlled change at the symbolic
layer, a corresponding change in the semantic structure space
$\mathcal{S}(L)$, and a consistent update at the decision layer via
the triple-layer connection.
\end{example}

In later chapters, PACER will manipulate sequences of base packs and
other elements of $\mathcal{M}$ to compress, align, expand, and
recompress lexical material, while respecting semantic and law-space
constraints.

\section[Power Submonoid Spectra and Compliance Pipelines]{Power Submonoid Spectra and Compliance Pipelines}


The lexical KAT action monoid $\mathcal{M}$ will typically be very
large. For PACER and certificate-based AI, we are not interested in
arbitrary subsets of $\mathcal{M}$, but in submonoids with specific
functional roles: memory, prediction, halting and termination,
compliance and audit, and so on. We group these submonoids into what
we call a power submonoid spectrum.

\begin{definition}[Power submonoid spectrum of a lexical KAT action monoid]
\label{def:power-submonoid-spectrum}
Let $\mathcal{M}$ be a lexical KAT action monoid as in
Definition~\ref{def:lex-kat-action-monoid}. A \emph{power submonoid
spectrum} for $\mathcal{M}$ is a distinguished family of submonoids
\[
  \mathsf{Spec}_{\mathrm{KAT}}(\mathcal{M})
  \subseteq
  \{\, N\subseteq\mathcal{M} \mid N\text{ is a submonoid} \,\},
\]
closed under finite intersections, such that each $N\in
\mathsf{Spec}_{\mathrm{KAT}}(\mathcal{M})$ is assigned a functional
role (memory, prediction, halting, compliance, etc.) compatible with
the HACA/PACER semantics. Typical examples include:
\begin{itemize}
  \item \emph{Historical-topology submonoids} $N_{\mathrm{hist}}$,
        generated by operators that store or retrieve historical
        context (e.g.\ reintroducing earlier packs, replaying previous
        sentences, or reconstructing prior states).
  \item \emph{Predictive-topology submonoids} $N_{\mathrm{pred}}$,
        generated by operators that insert predicted or hypothesized
        packs (e.g.\ forecasting missing segments, predicting
        follow-up clauses, or proposing candidate continuations).
  \item \emph{Halting-closure submonoids} $N_{\mathrm{halt}}$, whose
        elements ensure that certain KAT paths terminate or enter a
        designated stopping region (e.g.\ inserting explicit end-of-
        section markers or halting tokens).
  \item \emph{Compliance-pipeline submonoids} $N_{\mathrm{comp}}$,
        generated by operators that implement audit and compliance
        procedures (e.g.\ inserting mandatory disclaimers, normalizing
        terminology to a regulated vocabulary, or removing forbidden
        constructions).
\end{itemize}
\end{definition}

Each submonoid in $\mathsf{Spec}_{\mathrm{KAT}}(\mathcal{M})$ can be
thought of as a ``mode of operation'' for lexical transformations,
with a clear semantic interpretation at the HACA level. For example,
historical-topology submonoids organize the ways in which PACER may
safely revisit and reuse earlier material; predictive submonoids
capture lawful extrapolations or interpolations; halting-closure
submonoids encode termination guarantees; and compliance-pipeline
submonoids encode audit trails and normalization procedures.

\begin{remark}[Submonoid spectra and law-space categories]
\label{rem:submonoid-spectra-lawspace}
Through the semantic--law projection
$\Phi_L:\Omega(L)\to B_L\subseteq\GZLS$, each submonoid
$N\in\mathsf{Spec}_{\mathrm{KAT}}(\mathcal{M})$ determines a family of
law-space sectors and associated semantic property classes. In later
chapters, GRL path integrals, MDQ, and Jacobiator-gap certificates
will often be specified not just for arbitrary KAT paths, but for
KAT paths constrained to lie in particular submonoids. The power
submonoid spectrum thus serves as a semantic index over which we can
quantify and certify behavior.
\end{remark}

\section[Language Examples and Normative Subspaces]{Language Examples and Normative Subspaces}


The lexical KAT action monoid formalism is agnostic to the choice of
natural language: Chinese, English, domain-specific jargon, and mixed
registers can all be embedded into a common monoid of lexical
operators. What changes is the choice of normative subspaces and
baselines for audit, safety, and efficiency.

\begin{definition}[Normative subspaces in lexical KAT monoids]
\label{def:normative-subspace}
Let $\mathcal{M}$ be a lexical KAT action monoid associated with a
HACA law-object $L$. A \emph{normative subspace} for a given baseline
(e.g.\ audit, legal compliance, or style) is a submonoid
$N_{\mathrm{norm}}\subseteq\mathcal{M}$ together with a family of
base packs $\mathcal{B}_{\mathrm{norm}}\subseteq\mathcal{B}(L)$ such
that:
\begin{enumerate}
  \item Every operator in $N_{\mathrm{norm}}$ preserves the chosen
        baseline (for example, does not introduce forbidden terms and
        respects required stylistic or legal constraints).
  \item For every admissible lexical configuration that must be
        audited, there exists a sequence of operators in
        $N_{\mathrm{norm}}$ and canonical expansion operators that
        maps it into a configuration built exclusively from base
        packs in $\mathcal{B}_{\mathrm{norm}}$.
\end{enumerate}
We then say that $N_{\mathrm{norm}}$ is a (baseline-dependent)
normative subspace of $\mathcal{M}$.
\end{definition}

Remark~\ref{rem:mandarin-minimal-normative-subspace} below provides a
concrete example in the context of Chinese: standard written Mandarin
is taken as a minimal normative subspace for audit and alignment,
while dialects and idiomatic expressions are seen as compressed or
deformed versions that can be expanded back into Mandarin when
necessary.

\begin{remark}[Language Example: Mandarin as a Minimal Normative Subspace]
\label{rem:mandarin-minimal-normative-subspace}
As a concrete linguistic example, consider the law-space of Chinese language.
Different regional dialects, registers, and stylistic habits can be modeled as
distinct law-objects or regions in this space, while standard written Mandarin
may be viewed as a minimal normative subspace under an audit/alignment
baseline. Idioms, regional expressions, and highly compressed metaphors
correspond to higher-order ``compressed'' word-pack operators in the total
operator bundle, whose semantics can be expanded, when needed, into longer
normative word-pack sequences inside the Mandarin subspace.

Under the semantic homotopy framework of this volume, a dialectal path of
expressions is considered normatively equivalent to a canonical Mandarin path
if there exists a chain of property-changing and expansion operators that maps
the former into the Mandarin subspace while preserving non-empty property
intersections and admitting homotopy-filling of Jacobiator-gaps. In this
sense, the Mandarin subspace realizes a Level~C normative subspace for audit
and alignment: every admissible expression that must be audited can be
homotopically projected onto a normative path in Mandarin. However, with
respect to an efficiency baseline (for example, shortest utterances or
maximum stylistic compression), this subspace is not necessarily optimal:
dialectal or idiomatic forms may be more efficient. The LA/LB/LC levels
measure the degree of normalizability under a chosen baseline, rather than
enforcing a single gauge as universally optimal.
\end{remark}



\begin{example}[English legal texts and normative subspaces]
\label{ex:english-legal-normative-subspace}
As another example, consider English legal and regulatory texts. One
can distinguish a normative subspace generated by operators that
insert, normalize, and rearrange phrases according to a legal drafting
standard (for instance, a specific jurisdiction's legislative style
guide). Plain-language paraphrases, informal summaries, and domain-
specific jargon can then be viewed as compressed or relaxed versions
of these normative forms. For audit and compliance, lexical operators
in the compliance-pipeline submonoids
$N_{\mathrm{comp}}\in\mathsf{Spec}_{\mathrm{KAT}}(\mathcal{M})$ are
used to project arbitrary text back into the legal normative subspace
before semantic and law-space–level certificates are checked.
\end{example}

These examples illustrate how different natural languages and registers
can be embedded into a common lexical KAT action monoid, with different
normative subspaces chosen according to the baseline (audit, safety,
efficiency) of interest. The HACA/PACER framework abstracts away from
surface linguistic differences and focuses on the law-space–aware
structure of lexical operators and their submonoid spectra.

\chapter[From Context Iteration to Homotopy Networks of KAT Monoid Segments]{From Context Iteration to Homotopy Networks of KAT Monoid Segments:\\
         HACA/PACER as a Homotopy Network of KAT Action Monoids}


In the previous chapter, lexical KAT action monoids were introduced as
structured submonoids $\mathcal{M}\subseteq\mathrm{End}(\Sigma^*)$
acting on lexical sequences, together with a power submonoid spectrum
 encoding different functional roles (memory, prediction, halting,
 compliance, etc.). In this chapter we move from static monoids to
 \emph{context iteration}: we view PACER runs as sequences of local
 KAT submonoids generated by the operators actually used at each
context window, and then assemble these local pieces into a homotopy
 network of ``paragraph $\times$ submonoid'' nodes. This network will
 serve as the combinatorial skeleton on which reading-path certificates
 and law-space homotopy structures are built.

\begin{definition}[PACER Engine on a HACA Law-Object]
\label{def:pacer-engine-haca}
Let $L$ be a HACA law-object in the sense of
Definition~\ref{def:haca-law-objects-category}, equipped with:
\begin{itemize}
  \item a lexical KAT action monoid $\mathcal{M}$ as in
        Definition~\ref{def:lex-kat-action-monoid};
  \item paragraph-level submonoids $M(P_i)$ and paragraph nodes
        $v_i=(P_i,M(P_i))$ as in
        Definition~\ref{def:paragraph-submonoids};
  \item a triple-layer connection
        $\GZTLC_L$ and semantic gauge field $\mathcal{A}_L$ as in
        Definition~\ref{def:haca-triple-layer-connection} and
        Definition~\ref{def:semantic-gauge-field};
  \item safe law-space sectors $S\subseteq B_L$ and baseline-admissible
        MDQ region $\mathcal{Q}_{\mathrm{adm}}(L)$ as in
        Definition~\ref{def:safe-unsafe-sectors} and
        Definition~\ref{def:baseline-constrained-mdq-region}.
\end{itemize}
A \emph{PACER engine} on $L$ consists of the following data:
\begin{enumerate}
  \item A family of admissible policies $\Pi(L)$ acting on the lexical
        KAT action monoid $\mathcal{M}$ and on the paragraph-level
        submonoids $M(P_i)$.
  \item A context-iteration mechanism which, for each policy
        $\pi\in\Pi(L)$ and each input textual object (document or
        stream), produces a sequence of local generator sets
        $\{\mathcal{G}_t\}_t$ and corresponding local context submonoids
        $M_t=\langle\mathcal{G}_t\rangle\subseteq\mathcal{M}$ as in
        Definition~\ref{def:local-context-submonoid}, together with
        paragraph-level submonoids $M(P_i)$ and paragraph nodes
        $v_i=(P_i,M(P_i))$ as in
        Definition~\ref{def:paragraph-submonoids}.
  \item A reading-path construction which maps each run of $\pi$ to a
        PACER reading path
        \[
          \gamma^{(\pi)}:
          v_{i_0}\xrightarrow{F_{i_0}}v_{i_1}\xrightarrow{F_{i_1}}\cdots
          \xrightarrow{F_{i_{r-1}}}v_{i_r}
        \]
        in the paragraph--monoid network, where the morphisms
        $F_{i_j}:M(P_{i_j})\to M(P_{i_{j+1}})$ are paragraph-level KAT
        monoid morphisms in the sense of
        Definition~\ref{def:paragraph-monoid-morphism}, and
        $\gamma^{(\pi)}$ belongs to a homotopy class
        $[\gamma^{(\pi)}]$ as in
        Definition~\ref{def:pacer-reading-path-homotopy}.
  \item A GRL-based evaluation mechanism which associates to each run
        of $\pi$:
        \begin{itemize}
          \item a lexical GRL action $S_{\mathrm{lex}}[w]$ and path
                integral $\mathcal{Z}_{\mathrm{lex}}[w]$ on operator
                paths $w\in\mathsf{Paths}(\mathcal{M})$, as in
                Definition~\ref{def:lexical-grl-action-on-M};
          \item semantic GRL path integrals
                $\mathcal{Z}^{\mathrm{sem}}_\chi[\gamma]$ and predicates
                $\Xi_\chi(\gamma)$ on paragraph operators
                $\gamma\in\Omega^{(3)}(L)$ for certificate predicates
                $\chi\in\mathcal{C}_L$, as in
                Definition~\ref{def:semantic-grl-path-integral} and
                Definition~\ref{def:semantically-admissible-paragraph-operators};
          \item structural coordinates $(\mu^\pi,\alpha^\pi,\mathbf{v})$
                in the sense of
                Definition~\ref{def:mu-alpha-structural-coordinates}
                and Definition~\ref{def:pack-value-geometry},
                together with a semantic MDQ profile
                $\mathbf{q}^{\mathrm{sem}}(\pi)$ as in
                Definition~\ref{def:semantic-mdq-profile}.
        \end{itemize}
  \item Safety and admissibility conditions requiring that, for each
        run of $\pi$:
        \begin{itemize}
          \item the induced law-space path and its horizontal lift in
                the lexical principal bundle remain inside safe sectors
                $S$ and within the baseline-admissible MDQ region
                $\mathcal{Q}_{\mathrm{adm}}(L)$;
          \item lexical and semantic GRL actions satisfy the prescribed
                acceptance criteria for all relevant certificate
                predicates;
          \item the resulting reading-path certificate, in the sense of
                Definition~\ref{def:reading-path-certificate}, is
                well-defined.
        \end{itemize}
\end{enumerate}
We refer to $(L,\mathcal{M},\Pi(L))$ together with items (2)–(5) as a
\emph{PACER engine} (Pack-Aligned \\Compressive–Expansion Reasoner) on
the HACA law-object $L$.
\end{definition}



\section{Local Context Iteration as KAT Submonoids}


A PACER run proceeds in discrete context windows: at step $t$, the
engine receives a context (for instance, a sliding window of sentences
around the current position), applies a finite family of lexical
operators to compress, align, and expand content, and then moves on to
the next window. Each such step uses only a small subset of the global
lexical KAT action monoid~$\mathcal{M}$.

\begin{definition}[Local generator sets and context submonoids]
\label{def:local-context-submonoid}
Let $\mathcal{M}$ be a lexical KAT action monoid for a HACA law-object
$L$, as in Definition~\ref{def:lex-kat-action-monoid}. Consider a
PACER run indexed by discrete context steps $t=0,1,2,\dots$. For each
step $t$, let
\[
  \mathcal{G}_t
  :=
  \{\, G_{t,k}\in\mathcal{M} \mid k\in K_t \,\}
\]
be the finite set of lexical operators actually applied by PACER
during step~$t$ (including operators used for compression, alignment,
expansion, recompression, and any tests or control operators used in
this window). The \emph{local context submonoid} at step~$t$ is the
submonoid
\[
  M_t
  :=
  \langle \mathcal{G}_t\rangle
  \subseteq \mathcal{M},
\]
generated by~$\mathcal{G}_t$ under composition and containing the
identity. Thus $M_t$ consists of all finite compositions of operators
actually used (or admissible) at context step~$t$.
\end{definition}

\begin{remark}[Closure under composition and local KAT structure]
\label{rem:local-closure}
By construction, $M_t$ is a submonoid of $\mathcal{M}$ and hence of
$\mathrm{End}(\Sigma^*)$: it is closed under composition and contains
the identity operator. The KAT structure on $\mathcal{M}$ restricts to
$M_t$ in the sense that each $G_{t,k}$ carries an interpretation in the
ambient KAT model and the induced tests and control-flow operations are
still meaningful when restricted to the context window at step~$t$.
Thus $M_t$ can be viewed as the KAT action monoid describing the local
syntactic degrees of freedom available to PACER at that specific
context.
\end{remark}

\begin{example}[Local context submonoid for a definition window]
\label{ex:local-definition-window}
Suppose $P$ contains a technical definition paragraph. At a context
window $t$ covering this paragraph, PACER may use operators that:
\begin{itemize}
  \item normalize variable names and notation;
  \item insert missing parts of a definition schema;
  \item expand abbreviations into full phrases;
  \item align the structure to a canonical pattern (e.g.\ ``Term \dots
        means \dots'' or ``Definition. \dots'').
\end{itemize}
The corresponding generator set $\mathcal{G}_t$ consists of these
normalization, expansion, and alignment operators, and $M_t$ contains
all compositions of such operators. In later windows, $t'$ covering
proofs or examples, the generator sets and submonoids $M_{t'}$ will
be different, reflecting the different syntactic tasks at those
locations.
\end{example}

From the standpoint of this chapter, a PACER run is thus a sequence of
pairs $(\text{context}_t, M_t)$, where $\text{context}_t$ can be
encoded by the current positions in the document and $\mathcal{G}_t$
(and hence $M_t$) encodes the local lexical and control operations
available at that step.

\section{Paragraph-Level Segmentation and Node Construction}


Local context windows are finer than paragraphs: a single paragraph may
be visited multiple times, under different PACER modes (e.g.\ initial
reading, cross-reference resolution, summarization). For the global
structure of a document, however, paragraphs are natural units. We
therefore aggregate local context submonoids into paragraph-level
submonoids.

Let a fixed document be segmented into paragraphs
\[
  P_1, P_2,\dots, P_N.
\]
For each paragraph $P_i$, consider the set $T_i$ of context steps $t$
during a PACER run that intersect $P_i$ in a prescribed way (for
example, windows whose center lies in $P_i$ or windows that include a
sentence from $P_i$).

\begin{definition}[Paragraph-level generator sets and submonoids]
\label{def:paragraph-submonoids}
For each paragraph $P_i$, define the paragraph-level generator set
\[
  \mathcal{G}_i
  :=
  \bigcup_{t\in T_i}\mathcal{G}_t,
\]
where $\mathcal{G}_t$ is the local generator set at step~$t$ as in
Definition~\ref{def:local-context-submonoid}. The corresponding
paragraph-level submonoid is
\[
  M(P_i)
  :=
  \langle \mathcal{G}_i\rangle
  \subseteq\mathcal{M}.
\]
The \emph{paragraph node} associated with $P_i$ is the pair
\[
  v_i
  :=
  \bigl(P_i, M(P_i)\bigr).
\]
\end{definition}

\begin{remark}[Document skeleton as a paragraph–monoid graph]
\label{rem:paragraph-monoid-graph}
The collection of nodes
\[
  \{v_i=(P_i,M(P_i))\}_{i=1}^N
\]
can be viewed as a skeleton of the document: each node encodes not
only a textual unit but also the KAT submonoid of lexical operators
that PACER tends to use around it. Directed edges between nodes will
be introduced in the next section, reflecting transitions between
paragraphs along PACER reading paths. In this way, the document is
represented as a graph of ``paragraph $\times$ KAT submonoid'' nodes,
rather than as a mere sequence of text blocks.
\end{remark}

The paragraph-level submonoids $M(P_i)$ are not arbitrary: they are
typically contained in or generated by intersections of the global
power submonoid spectrum
$\mathsf{Spec}_{\mathrm{KAT}}(\mathcal{M})$ introduced earlier (for
instance, the intersection of historical, predictive, and compliance
submonoids relevant to paragraph~$P_i$). This provides a semantic
classification of node types.

\section{PACER Reading Paths as Homotopy-Equivalent Networks}


We now introduce edges between paragraph nodes. Intuitively, an edge
from $v_i$ to $v_j$ should encode the admissible ways for PACER to
transform lexical material associated with $P_i$ into lexical material
associated with $P_j$, under the constraints of the HACA/PACER
semantics and law-space geometry.

\begin{definition}[Paragraph-level KAT monoid morphisms]
\label{def:paragraph-monoid-morphism}
Let $P_i$ and $P_{i+1}$ be two consecutive paragraphs (in reading
order), with corresponding submonoids $M(P_i)$ and $M(P_{i+1})$. A
\emph{paragraph-level KAT monoid morphism} from $P_i$ to $P_{i+1}$ is
a monoid homomorphism
\[
  F_i : M(P_i)\longrightarrow M(P_{i+1})
\]
that is compatible with the HACA semantics and law-space projections
in the following sense:
\begin{enumerate}
  \item For each $u\in M(P_i)$ and each lexical configuration $w$ in
        the support of $P_i$, the transformed configuration
        $F_i(u)(w)$ lies in the lexical support of $P_{i+1}$ (or in a
        designated buffer associated with it).
  \item The induced map on semantic structures and law-space sectors,
        via $\operatorname{Sem}_L$ and $\operatorname{Law}_L$, is
        compatible with the semantic gauge field and safe sectors, so
        that admissible reading trajectories remain within safe
        law-space regions.
\end{enumerate}
\end{definition}

Given a family of morphisms $\{F_i\}_{i=1}^{N-1}$, a PACER reading
path can be represented as a path in the paragraph–monoid graph.

\begin{definition}[PACER reading paths]
\label{def:pacer-reading-paths}
A \emph{PACER reading path} for a document with nodes
$v_i=(P_i,M(P_i))$ is a sequence
\[
  \gamma^{(\ell)}:
  v_{i_0}\xrightarrow{F_{i_0}}v_{i_1}\xrightarrow{F_{i_1}}\cdots
  \xrightarrow{F_{i_{r-1}}}v_{i_r},
\]
where each $F_{i_j}:M(P_{i_j})\to M(P_{i_{j+1}})$ is a
paragraph-level KAT monoid morphism as in
Definition~\ref{def:paragraph-monoid-morphism}. The index~$\ell$
labels different runs or different policies (e.g.\ different PACER
parameter settings).
\end{definition}

Two reading paths may visit the same paragraphs in different orders,
or use different submonoid morphisms between the same nodes. To
capture the idea that such paths can represent the same semantic and
law-space evolution, we introduce a notion of homotopy between
reading paths.

\begin{definition}[Homotopy between PACER reading paths]
\label{def:pacer-reading-path-homotopy}
Let $\gamma^{(1)}$ and $\gamma^{(2)}$ be two PACER reading paths
between the same start and end nodes, say
\[
  \gamma^{(1)}: v_a\to\cdots\to v_b,
  \qquad
  \gamma^{(2)}: v_a\to\cdots\to v_b.
\]
We say that $\gamma^{(1)}$ and $\gamma^{(2)}$ are \emph{homotopy
equivalent} (written $\gamma^{(1)}\sim\gamma^{(2)}$) if there exists a
finite sequence of elementary homotopy moves, each supported on a
small subdiagram of nodes
\[
  v_i\to v_j\to v_k
  \quad\text{or}\quad
  v_i\to v_j\leftarrow v_k,
\]
such that:
\begin{enumerate}
  \item At the level of lexical configurations and KAT monoid
        morphisms, the corresponding squares or triangles commute up
        to controlled error inside safe sectors.
  \item At the semantic and law-space levels, the two paths induce
        law-space trajectories that are related by a homotopy
        satisfying the safety and MDQ constraints of
        Definitions~\ref{def:safe-unsafe-sectors}
        and~\ref{def:baseline-constrained-mdq-region}.
\end{enumerate}
The resulting equivalence classes $[\gamma]$ are called
\emph{reading-path homotopy classes}.
\end{definition}

\begin{remark}[2-cells and homotopy networks of KAT monoid segments]
\label{rem:2-cells-homotopy-networks}
The elementary homotopy moves in
Definition~\ref{def:pacer-reading-path-homotopy} can be visualized as
2-cells (faces) in a higher-dimensional network whose 1-cells are
paragraph-level morphisms $F_i$ and whose 0-cells are paragraph
nodes~$v_i$. Gluing these 2-cells along shared edges, we obtain a
homotopy network of KAT monoid segments: different PACER runs
correspond to different paths in this network, and homotopy-equivalent
runs correspond to paths that bound a finite number of 2-cells. Long-
range topic cycles and semantic loops in a document can thus be seen
as nontrivial loops in this network, possibly carrying nontrivial
Jacobiator-gap or MDQ data.
\end{remark}

\section{Law-Space Interpretation and Certificate Structures}


The paragraph–monoid homotopy network constructed above has a natural
law-space interpretation. The lexical KAT action monoid $\mathcal{M}$
and its submonoids act on the free monoid $\Sigma^*$ of lexical
sequences, which is related to the operator universe $\Omega(L)$ of a
HACA law-object via an embedding of lexical structures. The semantic–
law projection $\Phi_L:\Omega(L)\to B_L\subseteq\GZLS$ then sends KAT
paths to law-space paths.

\begin{definition}[Lexical KAT principal bundle over law-space]
\label{def:lex-kat-principal-bundle}
Let $L$ be a HACA law-object with law-support $B_L\subseteq\GZLS$ and
lexical KAT action monoid $\mathcal{M}$. Consider the bundle
\[
  P_{\mathrm{lex}}
  :=
  B_L\times(\Sigma^*,\mathcal{M})
\]
whose fiber over $x\in B_L$ consists of lexical configurations and
KAT operators acting on them. Equipped with the natural right action
of $\mathcal{M}$ on the $\Sigma^*$-component, $P_{\mathrm{lex}}$ is a
principal $\mathcal{M}$-bundle over $B_L$ in the trivialized sense.
The semantic gauge field and triple-layer connection of $L$ induce a
connection on $P_{\mathrm{lex}}$, allowing us to define horizontal
lifts of law-space paths.
\end{definition}

Each PACER reading path $\gamma$ in the paragraph–monoid network can
be lifted to a path in $P_{\mathrm{lex}}$ by choosing, for each
paragraph node, representative lexical configurations and KAT operators
consistent with the connection. Homotopy-equivalent reading paths
induce law-space homotopies satisfying the safety and MDQ constraints
of the previous chapters. This suggests a natural notion of
reading-path certificate.

\begin{definition}[Reading-path certificates]
\label{def:reading-path-certificate}
Let $L$ be a HACA law-object with safe sectors as in
Definition~\ref{def:safe-unsafe-sectors} and baseline-constrained MDQ
region $\mathcal{Q}_{\mathrm{adm}}(L)$ as in
Definition~\ref{def:baseline-constrained-mdq-region}. A
\emph{reading-path certificate} for a PACER reading path $\gamma$ is
the data of:
\begin{enumerate}
  \item A reading-path homotopy class $[\gamma]$ in the paragraph–
        monoid network.
  \item A representative law-space path $\Gamma:[0,1]\to B_L$ and a
        horizontal lift $\tilde{\Gamma}$ in $P_{\mathrm{lex}}$
        compatible with the connection induced by the semantic gauge
        field.
  \item A semantic MDQ profile and value geometry pair
        $(\mathbf{q}^{\mathrm{sem}},\mathbf{v})$ lying in the
        admissible region defined by
        $\mathcal{Q}_{\mathrm{adm}}(L)$ and the law-space baseline.
\end{enumerate}
Two certificates are considered equivalent if they represent the same
reading-path homotopy class and the same law-space homotopy class,
with MDQ and value geometry data differing only within prescribed
tolerances.
\end{definition}

\begin{remark}[Towards certificate-based AI on homotopy networks]
\label{rem:cert-based-ai-homotopy-network}
Reading-path certificates provide a bridge between the combinatorial
structure of PACER runs (as paths in a paragraph–monoid network) and
the geometric structure of HACA law-space (as paths in $B_L$ equipped
with semantic gauge fields, safe sectors, and MDQ constraints). In the
certificate-based AI framework of later chapters, such certificates
can be used to:
\begin{itemize}
  \item compare different runs of PACER on the same document or
        related documents;
  \item audit whether a given PACER run remained within acceptable
        safe sectors and MDQ regions;
  \item construct homotopy-invariant summaries of ``how'' a result was
        obtained, not just ``what'' the result is.
\end{itemize}
In this sense, HACA/PACER is not only a mechanism for generating
language and reasoning steps, but also a generator of law-space
homotopy data that can be certified, compared, and governed.
\end{remark}

\begin{remark}[Horizontal and vertical higher-order homotopy in HACA/PACER]
\label{rem:horizontal-vertical-homotopy}
It is conceptually useful to distinguish two orthogonal directions in
which higher-order homotopy appears in the HACA/PACER framework.

\emph{Horizontally}, PACER moves along reading paths through the
paragraph--monoid network of
Definition~\ref{def:paragraph-submonoids} and
Definition~\ref{def:pacer-reading-paths}. Different runs (or different
policies) on the same document correspond to different paths
\[
  \gamma^{(\ell)}:
    v_{i_0}\xrightarrow{F_{i_0}}\cdots\xrightarrow{F_{i_{r-1}}}v_{i_r}
\]
in this network. Homotopies between such paths, as in
Definition~\ref{def:pacer-reading-path-homotopy}, generate 2-cells
encoding permissible reroutings of context iteration. This horizontal
direction captures the \emph{chapter/section topology} of a document:
how a HACA/PACER agent is allowed to traverse and revisit paragraphs,
sections, and chapters under the constraints of semantic gauge fields,
MDQ bounds, and safe sectors.

\emph{Vertically}, at each node $v_i=(P_i,M(P_i))$ there is a
multi-layer HACA fiber: the lexical KAT layer ($l_1$) of operators on
$\Sigma^*$, together with the higher semantic layers $l_n$ for $n\ge 2$
(property layers, homotopy/safety layers, metaphor and value layers)
over the law-space sector $B_L\subseteq\GZLS$. This fiber carries its
own higher-order homotopy structure, controlled by the triple-layer
connection $\GZTLC_L$, the CH-layering framework, Jacobiator-gaps, and
the LA/LB/LC homotopy completeness levels. Vertical homotopies deform
lexical, semantic, and decision-level data \emph{within} a fixed
paragraph node while preserving the required property intersections,
gap-fillability, and baseline constraints.

Taken together, these two directions form a ``double'' higher-order
homotopy picture: horizontally, along the chapter/paragraph topology
of context iteration; vertically, inside the multi-layer HACA fiber at
each node. Reading-path certificates from
Definition~\ref{def:reading-path-certificate} implicitly live in this
two-dimensional structure: they record both the horizontal homotopy
class of a PACER run in the paragraph--monoid network and the vertical
homotopy data in law-space (semantic properties, Jacobiator-gaps, MDQ
and value geometry) along the lifted path. This is the precise sense
in which HACA/PACER realizes a bidirectional higher-order homotopy
structure over documents.
\end{remark}



\section{Homotopy Completeness and Logical Chain Metrics}


\begin{definition}[Homotopy-Complete Segmentation and Jacobiator-Gap-Free Book]
\label{def:homotopy-complete-book}
Let \\$\{P_i\}_{i=1}^N$ be a segmentation of a book into minimal context units.
For each segment $P_i$ we associate:
\begin{itemize}
  \item a KAT action submonoid $M(P_i) \subset \mathcal{M}$ generated by the
        PACER operators effectively used on $P_i$;
  \item a non-empty set of semantic properties $\mathrm{Prop}(P_i)$ encoding
        the local thematic and logical commitments of $P_i$.
\end{itemize}
We say that the segmentation $\{P_i\}$ is \emph{homotopy-complete} and that
the book is \emph{Jacobiator-gap free under this segmentation} if the
following conditions hold:
\begin{enumerate}
  \item For every adjacent pair $(P_i,P_{i+1})$ there exists
        \begin{itemize}
          \item a KAT monoid morphism
                $F_i : M(P_i) \to M(P_{i+1})$ describing PACER-style
                context migration, and
          \item a property-changing operator $S_{i \to i+1}$ on the property
                layer
        \end{itemize}
        such that
        \[
          \mathrm{Prop}(P_i)
          \cap S_{i \to i+1}^{-1}\bigl(\mathrm{Prop}(P_{i+1})\bigr)
          \neq \varnothing.
        \]
        In particular, the property layer is intersection-closed along each
        local transition.
  \item There exists a designated ``terminal'' property
        $\mathrm{End}$ carried by at least one segment (for instance, in the
        last chapter, encoding that a given line of argument is concluded),
        and every PACER reading path from the first chapter to the last
        chapter can be realized, up to law-space homotopy, as a path whose
        property labels vary continuously and end at a node that carries
        $\mathrm{End}$.
\end{enumerate}
Whenever these conditions are satisfied, we say that the book is
\emph{Jacobiator-gap free} with respect to the segmentation
$\{P_i\}$ and the chosen property layer.
\end{definition}



\begin{example}[Finite Path Patterns with and without Logical Completion]
\label{ex:finite-path-patterns}
Consider nodes $1,2,3,4,5,6,7,8,9,10$ representing minimal context units,
each equipped with a submonoid $M(i)$ and a property set
$\mathrm{Prop}(i)$. Assume that node $10$ carries the terminal property
$\mathrm{End} \in \mathrm{Prop}(10)$.

Suppose that the following node sequences are PACER reading paths:
\[
\begin{tikzcd}[column sep=1.6em]
|[circle,draw]|{1} \arrow[r] & |[circle,draw]|{3} \arrow[r] & |[circle,draw]|{5} \arrow[r] & |[circle,draw]|{7} \arrow[r] & |[circle,draw]|{9} \arrow[r] & |[circle,draw]|{10}
\end{tikzcd}
\qquad
\begin{tikzcd}[column sep=1.6em]
|[circle,draw]|{1} \arrow[r] & |[circle,draw]|{2} \arrow[r] & |[circle,draw]|{10}
\end{tikzcd}
\]
\[
\begin{tikzcd}[column sep=1.6em]
|[circle,draw]|{1} \arrow[r] & |[circle,draw]|{4} \arrow[r] & |[circle,draw]|{6} \arrow[r] & |[circle,draw]|{10}
\end{tikzcd}
\qquad
\begin{tikzcd}[column sep=1.6em]
|[circle,draw]|{4} \arrow[r] & |[circle,draw]|{8} \arrow[r] & |[circle,draw]|{10}
\end{tikzcd}
\]
If for every consecutive pair $(i,j)$ appearing in these sequences there
exists a KAT monoid morphism $F_{ij}:M(i)\to M(j)$ and a property-changing
operator $S_{i\to j}$ such that
$\mathrm{Prop}(i) \cap S_{i\to j}^{-1}(\mathrm{Prop}(j)) \neq \varnothing$,
then each of the four sequences above is a homotopy-admissible reading path
from its starting node to the terminal node $10$.

By contrast, consider a candidate path of the form
\[
  \begin{tikzcd}[column sep=1.6em]
  |[circle,draw]|{1} \arrow[r] & |[circle,draw]|{4} \arrow[r] & |[circle,draw]|{6} \arrow[r, dashed, "\nrightarrow" description] & |[circle,draw]|{\varnothing} \arrow[r, dashed] & |[circle,draw]|{10}
  \end{tikzcd}
\]
If the prefix $1 \!-\! 4 \!-\! 6$ cannot be extended, under the same
constraints on KAT morphisms and property intersections, to any path that
actually reaches a node carrying $\mathrm{End}$, then this prefix is a
\emph{non-extendable path prefix} and witnesses a local breakdown of
homotopy completeness for this segmentation.
\end{example}



\begin{remark}[Logical Chain Scores and Jacobiator-Gap Residuals]
\label{rem:logical-chain-scores}
In the framework of Definition~\ref{def:homotopy-complete-book}, every
non-extendable path prefix that cannot reach a segment carrying the terminal
property $\mathrm{End}$ contributes a negative term to any quantitative
score of logical continuity for the book. Such prefixes can be interpreted
as residual Jacobiator-gap contributions on the homotopy network: they mark
local failures of property-intersection closure or of triple-wise
Jacobiator-style consistency.

Aggregated quantities built from these residuals---for example, the number
of non-extendable prefixes, or a weighted ``gap energy'' that sums their
severities across all triples of segments---provide quantitative support
for assessing the homotopy completeness of the entire book. In this sense,
path breakages and three-step inconsistencies can be used as negative scores
in a logical-chain metric, making the notion of ``Jacobiator-gap free''
operational for long-form texts.
\end{remark}



\begin{definition}[Homotopy Completeness Levels A/B/C]
\label{def:homotopy-completeness-levels}
Let $\gamma$ range over PACER/HACA reading paths from a designated start
law-object $L_{\mathrm{start}}$ to a designated target law-object
$L_{\mathrm{target}}$ carrying a terminal property $\mathrm{End}$. We
define three levels of homotopy completeness.

As further clarified in
Proposition~\ref{prop:gaugeable-operator-bundle-from-property-intersections},
these levels can also be understood as measuring the degree to which a
gaugeable operator bundle can be constructed and normalized over the union
$\mathcal{U}$ of property intersections induced by heterogeneous
property-changing morphisms and their higher-homotopy smoothing.

\begin{description}
  \item[Level A (existential completeness).]
    The system is Level A homotopy-complete if there exists at least one
    path $\gamma$ such that:
    \begin{enumerate}
      \item $\gamma$ is homotopy-admissible in the sense of
            Definition~\ref{def:homotopy-complete-book}
            (property-layer constraints hold and local Jacobiator-gaps can
            be homotopy-filled), and
      \item $\gamma$ ends at a node carrying $\mathrm{End}$.
    \end{enumerate}
    Intuitively, there is at least one ``explainable'' and logically
    complete trajectory from $L_{\mathrm{start}}$ to $L_{\mathrm{target}}$.

  \item[Level B (common-core completeness).]
    The system is Level B homotopy-complete if it is Level A, and in
    addition there exists a non-empty set of law-objects or property
    configurations $\mathcal{C}$ (a \emph{homotopy common core}) such that
    every good path $\gamma$ (i.e.\ every homotopy-admissible path ending
    at $\mathrm{End}$ with acceptable quantitative scores) passes through
    $\mathcal{C}$ on the property layer. In symbols, for the set of good
    paths $\Gamma_{\mathrm{good}}$ we require
    \[
      \bigcap_{\gamma \in \Gamma_{\mathrm{good}}}
      \mathrm{Im}_{\mathrm{prop}}(\gamma)
      \;\supseteq\;
      \mathcal{C}
      \;\neq\;\varnothing.
    \]
    Different trajectories may take different routes, but they all share a
    common semantic core region.

  \item[Level C (normative completeness).]
    The system is Level C homotopy-complete if it is Level B, and in
    addition every admissible path $\gamma$ that remains within a prescribed
    ``normative'' regime (for example, a standard proof system, a safety
    pipeline, or a regulated workflow) is homotopic, in the semantic sense
    of this volume, to a fixed canonical path $\gamma_{\ast}$ from
    $L_{\mathrm{start}}$ to $L_{\mathrm{target}}$. In other words, within
    the normative subspace, all good paths are homotopy-equivalent and
    share not only a common core but a common homotopy class.
\end{description}
\end{definition}



\begin{definition}[Homotopy Common Core]
\label{def:homotopy-common-core}
In the setting of Level B completeness above, a \emph{homotopy common
core} is any non-empty set $\mathcal{C}$ of law-objects or property
configurations such that every good path $\gamma$ from
$L_{\mathrm{start}}$ to $L_{\mathrm{target}}$ intersects $\mathcal{C}$ on
the property layer. The common core $\mathcal{C}$ can be interpreted as a
region of law-space where the essential semantic commitments of the system
are realized: different trajectories may diverge before or after
$\mathcal{C}$, but they all traverse the same core area when they are
logically complete and admissible.
\end{definition}

\begin{remark}[Negative Scoring for Non-Extendable Prefixes]
\label{rem:negative-scoring-nonextendable}
The framework of Levels A/B/C is compatible with the negative scoring
scheme for non-extendable prefixes discussed in
Example~\ref{ex:finite-path-patterns}. In practice, one can:
\begin{itemize}
  \item treat Level A as a minimal pass condition: at least one path from
        $L_{\mathrm{start}}$ to $L_{\mathrm{target}}$ is homotopy-complete;
  \item penalize every non-extendable prefix that cannot reach a node
        carrying $\mathrm{End}$ under the given admissibility constraints,
        by adding to a logical-chain ``gap energy'' $\mathsf{GapEnergy}$;
  \item use Level B and Level C as increasingly strong certificate
        conditions: Level B requires the existence of a homotopy common
        core; Level C requires that, within a prescribed normative regime,
        all good paths are semantically homotopy-equivalent to a canonical
        baseline path.
\end{itemize}
In this way, homotopy completeness becomes a graded notion, with Level A/B/C
corresponding to different strengths of homotopy commonality and different
levels of auditability and interpretability.
\par
In the general law-space setting developed in the pure mathematics
volume~\cite{GaoZheng2025GFramework}, the \\LA/LB/LC conditions can be
rephrased in terms of the vanishing of the corresponding obstruction
classes for law-space bundles and reductions of structure group. In
particular, the bundle-construction argument carried out in
Proposition~\ref{app:prop-operator-bundle-from-property-intersections}
should be read as the HACA/PACER specialization of that general
framework: LA/LB/LC provide precisely the homotopy completeness needed
to trivialize the Jacobiator-gap obstruction and to realize the
operator bundle over the relevant law-space region.
\end{remark}



\begin{remark}[Interpretation and Information-Compensation View]
\label{rem:homotopy-commonality-interpretation}
The strength of homotopy commonality can be understood as a graded measure
of ``how much different paths really instantiate the same logic'' in a
given law-space and under a chosen baseline. At Level A, one only knows
that there exists at least one path that is homotopy-complete and reaches
the terminal property: qualitatively, there is at least one way to ``tell
the story correctly'', but other paths may diverge, break, or remain
unaudited. At Level B, all good paths share a homotopy common core
$\mathcal{C}$: explanations or strategies may differ in surface form, but
structurally they all traverse the same core region in law-space and on
the property layer. At Level C, within a prescribed normative subspace,
all good paths are semantically homotopy-equivalent to a canonical
baseline path $\gamma_\ast$, yielding a rigid form of norm-governed
behavior.

From the standpoint of higher homotopies and Jacobiator-gap repair, these
levels also reflect different demands on information-dimension
compensation. With weak commonality (Level A), one may need to introduce
significant path-specific higher-homotopy data to repair gaps along
individual trajectories. With stronger commonality (Level B), many repairs
can be centralized around the homotopy core $\mathcal{C}$, turning
path-specific patches into shared ``infrastructure''. In the strongest
case (Level C), a single normative repair around $\gamma_\ast$ suffices to
support an entire family of good paths inside the normative subspace,
minimizing the structural and informational overhead needed to maintain
reversible and solvable law-space dynamics.
\end{remark}





\part[GRL-SAC, MDQ, and Law-Space Learning in HACA/PACER]{GRL-SAC, MDQ, and Law-Space Learning in HACA/PACER\\
      with KAT Action Monoids}

\chapter[GRL Path Integrals on KAT Action Monoids]{GRL Path Integrals on KAT Action Monoids:\\
         Law-Space Learning over Lexical Operator Monoids}


\section[Lexical GRL Path Integrals on \texorpdfstring{$\mathcal{M}$}{M}]{Lexical GRL Path Integrals on \texorpdfstring{$\mathcal{M}$}{M}}


Let $\mathcal{M}$ be a lexical KAT action monoid for a HACA law-object
$L$, as in Definition~\ref{def:lex-kat-action-monoid}. Elements
$u\in\mathcal{M}$ act on lexical sequences $w\in\Sigma^*$ and, via the
embedding into the operator universe $\Omega(L)$, induce transformations
on semantic structures and law-space sectors.

\begin{definition}[Operator paths on lexical KAT action monoids]
\label{def:operator-paths-on-M}
A \emph{lexical operator path} on $\mathcal{M}$ is a finite or
continuous sequence
\[
  w = (u_0,u_1,\dots,u_T)
  \quad\text{or}\quad
  w:[0,T]\to\mathcal{M},
\]
depending on whether we work in discrete or continuous time. Given an
initial lexical configuration $x_0\in\Sigma^*$, the path induces a
lexical trajectory
\[
  x_{t+1} = u_t(x_t)
  \quad\text{(discrete time)},
\]
or
\[
  \frac{\mathrm{d}}{\mathrm{d}t}x_t
  = \mathcal{F}(x_t,u_t)
  \quad\text{(continuous time)},
\]
for a suitable dynamical law $\mathcal{F}$ derived from the underlying
GRL model. We write $\mathsf{Paths}(\mathcal{M})$ for the set of such
operator paths.
\end{definition}

In the Applications Volume~I, an action functional
$S_{\mathrm{lex}}[\gamma^{\mathrm{lex}}]$ is defined over law-space
paths (or operator paths) and used to construct GRL path integrals. In
the present lexical setting, we pull this structure down to
$\mathcal{M}$ by associating to each operator path $w$ a law-space
path via the semantic--law projection $\Phi_L$ and then evaluating the
GRL action along that path.

\begin{definition}[Lexical GRL action functional on \texorpdfstring{$\mathcal{M}$}{M}]
\label{def:lexical-grl-action-on-M}
Let $w\in\mathsf{Paths}(\mathcal{M})$ be an operator path with induced
lexical trajectory $\{x_t\}$ and associated law-space path
$\Gamma:[0,T]\to B_L$ obtained by composing the embedding of lexical
configurations into $\Omega(L)$ with $\Phi_L$. A \emph{lexical GRL
action functional} on $\mathcal{M}$ is a map
\[
  S_{\mathrm{lex}} : \mathsf{Paths}(\mathcal{M})\longrightarrow\mathbb{R}
\]
of the form
\[
  S_{\mathrm{lex}}[w]
  =
  \int_0^T \mathcal{L}_{\mathrm{lex}}\bigl(x_t,u_t,\Gamma(t)\bigr)\,\mathrm{d}t
\]
(or the corresponding discrete-time sum), where
$\mathcal{L}_{\mathrm{lex}}$ is a local Lagrangian encoding the
contribution of lexical MDQ (differential dynamical quanta) along the
path. The associated GRL path integral is
\[
  \mathcal{Z}_{\mathrm{lex}}[w]
  :=
  \int \exp\bigl(-S_{\mathrm{lex}}[w]\bigr)\,\mathcal{D}w,
\]
defined in the sense of the Applications Volume~I.
\end{definition}

\begin{remark}[Legal KAT paths as lexical minimal-action trajectories]
\label{rem:legal-kat-paths-minimal-action}
In this framework, ``legal'' KAT paths—for instance, paths that
respect syntactic and lexical constraints, or that implement specific
normalization schemes—are those for which $S_{\mathrm{lex}}[w]$ (or
$\mathcal{Z}_{\mathrm{lex}}[w]$) lies in an admissible sector of the
lexical action space, as in the Applications Volume~I. The lexical
GRL minimal-action principle thus governs how PACER may legally
manipulate word-packs and sentence structures at the purely lexical
level, before semantic and higher-layer constraints are taken into
account.
\end{remark}

\section[Semantic GRL Path Integrals and Admissible Paragraph Operators]{Semantic GRL Path Integrals and Admissible Paragraph Operators}
\label{subsec:semantic-grl-path-integral-haca}


In the Applications Volume~I, GRL path integrals were used to encode
law-aware dynamics as generalized action functionals over operator
paths. In the HACA/PACER setting of the present volume, we use the
semantic functor field
(Definition~\ref{def:semantic-gauge-field}) to construct analogous
path integrals whose role is to capture \emph{long semantic
expressions} along paragraph-level KAT paths. These path integrals
induce truth-valued predicates on paragraph operators and thus
determine which paragraphs are semantically admissible.



\begin{definition}[Semantic GRL Path Integral Based on Semantic Functors]
\label{def:semantic-grl-path-integral}
Let $L$ be a HACA law-object with semantic structure space
$\mathcal{S}(L)$, law-space sector $B_L\subseteq\GZLS$, and semantic
gauge field $\{F_\alpha\}_\alpha$ as in
Definition~\ref{def:semantic-gauge-field}. Let
$\Omega^{(3)}(L)$ denote the space of paragraph-level operator
paths (KAT path integrals) for $L$.

For a fixed truth-valued semantic predicate
$\chi:\mathcal{S}(L)\to\{\mathrm{true},\mathrm{false}\}$ (as in
Definition~\ref{def:haca-semantic-predicates}), we define a
\emph{semantic GRL path-integral predicate}
\[
  \Xi_\chi : \Omega^{(3)}(L)\longrightarrow
  \{\mathrm{true},\mathrm{false}\},
\]
as follows. For each paragraph operator (paragraph path)
$\gamma\in\Omega^{(3)}(L)$, regarded as a discrete or continuous
parametrized path of multi-scale operators,
\[
  \gamma : [0,1]\to \Omega(L),
\]
we consider the associated semantic path in $\mathcal{S}(L)$,
\[
  t\longmapsto s_t := \operatorname{Sem}_L(\gamma(t)),
\]
and form the GRL-style path-integral functional
\[
  \mathcal{Z}^{\mathrm{sem}}_\chi[\gamma]
\]
by substituting the semantic functor field (and, in particular,
the predicate $\chi$) into the GRL path-integral kernel as defined in
the Applications Volume~I. We then choose a designated admissible sector
$C_\chi$ in the value space of $\mathcal{Z}^{\mathrm{sem}}_\chi$ and
set
\[
  \Xi_\chi(\gamma)
  :=
  \begin{cases}
    \mathrm{true}, & \text{if }
    \mathcal{Z}^{\mathrm{sem}}_\chi[\gamma]\in C_\chi,\\[2pt]
    \mathrm{false}, & \text{otherwise.}
  \end{cases}
\]
The functional $\mathcal{Z}^{\mathrm{sem}}_\chi[\gamma]$ is called the
\emph{semantic GRL path integral} of $\gamma$ with respect to $\chi$,
and $\Xi_\chi$ is the induced truth-valued predicate on paragraph
operators.

\end{definition}

\begin{definition}[Semantically Admissible Paragraph Operators]
\label{def:semantically-admissible-paragraph-operators}
Let $\mathcal{C}_L$ be a distinguished family of certificate
predicates for $L$, i.e.\ a subset of truth-valued semantic predicates
$\chi\in\mathsf{Obs}(L)$ selected as semantic constraints or
certificates. A paragraph operator (paragraph path)
$\gamma\in\Omega^{(3)}(L)$ is called \emph{semantically admissible}
(with respect to $\mathcal{C}_L$) if and only if
\[
  \Xi_\chi(\gamma)=\mathrm{true}
\]
for all $\chi\in\mathcal{C}_L$, where $\Xi_\chi$ is the semantic GRL
path-integral predicate from
Definition~\ref{def:semantic-grl-path-integral}. The collection of all
semantically admissible paragraph operators is denoted by
\[
  \Omega^{(3)}_{\mathrm{adm}}(L)
  :=
  \bigl\{\,
    \gamma\in\Omega^{(3)}(L)\ \big|\ 
    \Xi_\chi(\gamma)=\mathrm{true}
    \text{ for all }\chi\in\mathcal{C}_L
  \,\bigr\}.
\]

\end{definition}

\begin{remark}[Long Semantic Expressions as Semantic GRL Path Integrals]
\label{rem:long-semantic-expression-as-grl-pi}
From the \\HACA/PACER perspective, a paragraph operator
$\gamma\in\Omega^{(3)}(L)$ can be viewed as a long KAT path built
from word-packs and sentence-level operators. The semantic GRL
path integrals $\mathcal{Z}^{\mathrm{sem}}_\chi[\gamma]$ for
certificate predicates $\chi\in\mathcal{C}_L$ are then the
law-space–aware \emph{long semantic expressions} associated with
$\gamma$: they aggregate, along the entire paragraph path, the local
contributions of the semantic functors and of the underlying
PFB-GNLA connection.

The requirement that $\Xi_\chi(\gamma)=\mathrm{true}$ for all
$\chi\in\mathcal{C}_L$ expresses the idea that a paragraph is not
merely a syntactically well-formed KAT path, but must also satisfy
global semantic and law-space constraints once evaluated through the
semantic functor field. In this precise sense, paragraph operators
are admissible only when their associated semantic GRL path integrals
satisfy the designated certificate conditions; otherwise, the
corresponding long semantic expressions are rejected as semantically
invalid.

\end{remark}

\begin{remark}[Two GRL Path Integrals and Two Minimal-Action Principles]
\label{rem:two-grl-pi-two-actions}
From the perspective of the GaoZheng G-Framework across volumes, the
semantic GRL path integral introduced in
Definition~\ref{def:semantic-grl-path-integral} should be understood as a
\emph{second}, semantic, minimal-action principle, distinct from (but
compatible with) the lexical GRL minimal-action principle of Applications
Volume I.

In the Applications Volume~I, differential dynamical quanta (MDQ) were
attached to \\character- or token-level operator paths and used to define a
GRL action functional
\[
  S_{\mathrm{lex}}[\gamma^{\mathrm{lex}}],
\]
whose associated GRL path integral
$\mathcal{Z}_{\mathrm{lex}}[\gamma^{\mathrm{lex}}]$ encodes the
law-aware dynamics of lexical organization. Legal word-packs, sentence
forms, and KAT-syntactic/grammatical configurations are precisely those
operator paths for which $S_{\mathrm{lex}}$ (or $\mathcal{Z}_{\mathrm{lex}}$)
satisfies the prescribed acceptance criteria: in this sense, the
character-level MDQ and the lexical GRL path integral realize a
\emph{minimal-action principle for language organization} at the KAT
word-pack level.

In the present volume, the semantic functor field
$\{F_\alpha\}_\alpha$ plays the role of a second family of MDQ, now
attached to the semantic structure space $\mathcal{S}(L)$. For each
certificate predicate $\chi\in\mathcal{C}_L$, the semantic GRL path
integral $\mathcal{Z}^{\mathrm{sem}}_\chi[\gamma]$ and its induced
predicate $\Xi_\chi(\gamma)$ define a \emph{minimal-action principle for
semantic dynamics}: a paragraph operator $\gamma\in\Omega^{(3)}(L)$ is
semantically admissible only if the semantic action associated with
$\gamma$ lies in the accepted sector $C_\chi$ for all semantic
certificates $\chi$.

Thus, there are in effect two coupled GRL path integrals and two
minimal-action principles:
\begin{enumerate}
  \item A \emph{lexical} GRL path integral, based on character-level MDQ,
        with minimal action $S_{\mathrm{lex}}[\gamma^{\mathrm{lex}}]$,
        governing the formation of legally admissible KAT word-packs and
        syntactic/grammatical structures.
  \item A \emph{semantic} GRL path integral, based on semantic functor
        MDQ, with semantic action functionals
        $\mathcal{Z}^{\mathrm{sem}}_\chi[\gamma]$ (or the underlying
        action $S^{\mathrm{sem}}_\chi[\gamma]$), governing the formation
        of legally admissible semantic content and semantic connections
        (law-level transformations under the principal-bundle structure).
\end{enumerate}
A paragraph operator $\gamma$ is fully admissible only when it passes
\emph{both} layers: it must be a legal lexical path according to the
character-based GRL minimal-action principle, and simultaneously a legal
semantic path according to the semantic GRL minimal-action principle
driven by the functor field. In this way, the operator word-pack universe
and the semantic connection universe are coupled by two GRL path
integrals sharing the same PFB-GNLA background but realizing distinct,
yet compatible, minimal-action constraints.

\end{remark}
 
\section[Occupancy Measures and Strategy Weights as Structural Coordinates]{Occupancy Measures and Strategy Weights as Structural Coordinates}


In the chapter on law-space baselines and pack-level value geometry,
we introduced occupancy measures $\mu_i^\pi$ and value geometry vectors
$\mathbf{v}=(v_i)_i$ indexing packs or paragraph-level units
(Definition~\ref{def:pack-occupancy-measures} and
Definition~\ref{def:pack-value-geometry}). In the present context,
these quantities arise naturally as structural coordinates for GRL
path integrals on lexical operator monoids.

Let $\{P_i\}$ denote a family of paragraph-level segments or pack
clusters, and let $I_L$ index a finite collection of KAT-based segments
(for example, the paragraph-level submonoids $M(P_i)$ from
Definition~\ref{def:paragraph-submonoids}). For a PACER/SAC-PACER-HACA
policy $\pi$, we can measure:
\begin{itemize}
  \item the \emph{occupancy measure} $\mu_i^\pi$, i.e.\ the fraction
        of time or semantic attention the policy spends in the
        $i$-th segment or submonoid;
  \item the \emph{strategy weight} $\alpha_i^\pi$, i.e.\ how strongly
        the policy prefers to use operators from the $i$-th segment or
        submonoid when multiple choices are available.
\end{itemize}

\begin{definition}[Occupancy measures and strategy weights on KAT segments]
\label{def:mu-alpha-structural-coordinates}
Let $I_L$ index a finite family of KAT-based segments. For a policy
$\pi$, we define:
\begin{itemize}
  \item the \emph{occupancy measure} vector
        $\mu^\pi=(\mu_i^\pi)_{i\in I_L}$, where $\mu_i^\pi$ is the
        long-run fraction of time or attention devoted to the $i$-th
        segment, as in Definition~\ref{def:pack-occupancy-measures};
  \item the \emph{strategy weight} vector
        $\alpha^\pi=(\alpha_i^\pi)_{i\in I_L}$, where $\alpha_i^\pi$
        is a normalized weight (for example, a probability or mixing
        coefficient) specifying how often PACER chooses operators from
        the $i$-th segment when a choice is available.
\end{itemize}
We interpret $(\mu^\pi,\alpha^\pi)$ as structural coordinates for the
behavior of $\pi$ on the lexical KAT action monoid~$\mathcal{M}$.
\end{definition}

Given a lexical GRL action functional $S_{\mathrm{lex}}[w]$ on
$\mathcal{M}$ and semantic GRL path integrals
$\mathcal{Z}^{\mathrm{sem}}_\chi[\gamma]$ on $\Omega^{(3)}(L)$, a
global performance objective for a policy $\pi$ can be expressed as
\[
  \mathcal{J}(\pi)
  =
  \mathbb{E}_\pi
  \Bigl[
    S_{\mathrm{lex}}[w]
    + \lambda_{\mathrm{sem}}
      \sum_{\chi\in\mathcal{C}_L}\Phi_\chi\bigl(\mathcal{Z}^{\mathrm{sem}}_\chi[\gamma]\bigr)
  \Bigr],
\]
where $\Phi_\chi$ extracts a scalar semantic penalty or reward from
the semantic path integral and $\lambda_{\mathrm{sem}}\ge 0$ is a
coupling parameter. Under mild regularity conditions, this objective
can be seen as a function of $(\mu^\pi,\alpha^\pi,\mathbf{v})$,
\[
  \mathcal{J}(\pi)
  \approx
  \widehat{\mathcal{J}}(\mu^\pi,\alpha^\pi,\mathbf{v}),
\]
with gradients
\[
  \frac{\partial\mathcal{J}}{\partial\mu_i},
  \qquad
  \frac{\partial\mathcal{J}}{\partial\alpha_i}
\]
serving as structural directions for policy updates in SAC-PACER-HACA
engines.

\section[Connection to SAC-PACER-HACA Engines]{Connection to SAC-PACER-HACA Engines}


In the SAC-PACER-HACA framework developed in later chapters, the
objective functionals and constraints introduced so far are combined
into a joint optimization problem. The lexical GRL action
$S_{\mathrm{lex}}[w]$ enforces minimal-action behavior at the KAT
word-pack level; the semantic GRL path integrals
$\mathcal{Z}^{\mathrm{sem}}_\chi[\gamma]$ enforce minimal-action
behavior at the semantic level; the occupancy measures $\mu_i^\pi$,
strategy weights $\alpha_i^\pi$, and value geometry vector
$\mathbf{v}$ capture how a policy $\pi$ distributes its effort across
KAT segments and semantic roles; and the semantic MDQ profile
$\mathbf{q}^{\mathrm{sem}}(\pi)$
(Definition~\ref{def:semantic-mdq-profile})
quantifies the semantic dynamical cost associated with each segment.

At a high level, a SAC-PACER-HACA engine can be viewed as performing
soft actor–critic updates on these structural coordinates:
\begin{itemize}
  \item the \emph{actor} adjusts $\alpha^\pi$ (and, implicitly, the
        policy on $\mathcal{M}$) to reduce the expected GRL action and
        semantic penalties, while respecting occupancy constraints and
        safe sectors;
  \item the \emph{critic} estimates value functions or advantage
        signals as functions of $(\mu^\pi,\alpha^\pi)$ and law-space
        state, using the MDQ and value geometry data as features.
\end{itemize}
In this interpretation, GRL path integrals on KAT action monoids
provide the law-space learning backbone for SAC-PACER-HACA engines,
and $(\mu_i,\alpha_i)$ serve as the structural coordinates in which
both policy and value updates are expressed.

\section[HACA-Based GRL Cognitive Scheme]{HACA-Based GRL Cognitive Scheme}


In the preceding sections, we have introduced:
\begin{itemize}
  \item lexical GRL path integrals on the lexical KAT action monoid
        $\mathcal{M}$ (Definition~\ref{def:lexical-grl-action-on-M});
  \item semantic GRL path integrals on paragraph operators and the
        notion of semantically admissible paragraph paths
        (Definition~\ref{def:semantic-grl-path-integral} and
        Definition~\ref{def:semantically-admissible-paragraph-operators});
  \item pack-level occupancy measures $\mu_i^\pi$, strategy weights
        $\alpha_i^\pi$, and value geometry vectors $\mathbf{v}$ as
        structural coordinates
        (Definition~\ref{def:pack-occupancy-measures},
         Definition~\ref{def:pack-value-geometry},
         Definition~\ref{def:mu-alpha-structural-coordinates});
  \item homotopy networks of paragraph-level KAT monoid segments,
        PACER reading paths, and reading-path certificates
        (Definitions~\ref{def:paragraph-submonoids},
         \ref{def:pacer-reading-paths},
         \ref{def:reading-path-certificate}, and
         Remark~\ref{rem:horizontal-vertical-homotopy}).
\end{itemize}
We now summarize these ingredients as a unified HACA-based GRL
cognitive scheme.

\begin{definition}[HACA-Based GRL Cognitive Scheme]
\label{def:haca-grl-cognitive-scheme}
Let $L$ be a HACA law-object with lexical KAT action monoid
$\mathcal{M}$, paragraph–monoid network, triple-layer connection
$\GZTLC_L$, semantic gauge field $\mathcal{A}_L$, safe sectors
$S\subseteq B_L$, and baseline-admissible MDQ region
$\mathcal{Q}_{\mathrm{adm}}(L)$. A \emph{HACA-based GRL cognitive
scheme} on $L$ is a procedure which, for each input lexical object
(e.g.\ a document represented as a raw string in $\Sigma^*$), carries
out the following steps:
\begin{enumerate}
  \item \textbf{Lexical heterogeneous KAT decomposition (syntax level).}
        The raw input string is first embedded into the free monoid
        $(\Sigma^*,\circ,\varepsilon)$ and the lexical KAT action
        monoid $\mathcal{M}$. Using lexical GRL action
        $S_{\mathrm{lex}}[w]$ and the associated path integrals
        $\mathcal{Z}_{\mathrm{lex}}[w]$, together with PACER's
        compressive–alignment–expansion–recompression cycle on
        $\mathcal{M}$, the scheme selects a family of operator paths
        and segmentations that decompose the input into heterogeneous
        KAT segments and paragraph paths
        $\gamma\in\Omega^{(3)}(L)$, subject to lexical minimal-action
        and syntactic admissibility constraints.
  \item \textbf{Multi-layer semantic GRL evaluation (semantic layers).}
        For each candidate paragraph path
        $\gamma\in\Omega^{(3)}(L)$ arising from the lexical
        decomposition, the scheme evaluates semantic GRL path
        integrals $\mathcal{Z}^{\mathrm{sem}}_\chi[\gamma]$ for a
        family of certificate predicates $\chi\in\mathcal{C}_L$
        covering the relevant law-space layers (meaning at $l_2$,
        safety and robustness at $l_3$, and higher-layer structures at
        $l_{n\ge 4}$). The resulting predicates
        $\Xi_\chi(\gamma)$ are used to accept or reject candidate
        decompositions and to provide semantic scores or penalties for
        each paragraph path.
  \item \textbf{Joint optimization over structural coordinates.}
        Using the pack-level occupancy measures $\mu_i^\pi$, strategy
        weights $\alpha_i^\pi$, value geometry vector $\mathbf{v}$,
        and semantic MDQ profile $\mathbf{q}^{\mathrm{sem}}(\pi)$, the
        scheme defines a joint objective
        $\mathcal{J}(\pi)\approx\widehat{\mathcal{J}}(\mu^\pi,\alpha^\pi,
        \mathbf{v})$ combining lexical and semantic GRL actions,
        subject to safe-sector and MDQ constraints. A SAC-PACER-HACA
        engine then adjusts the policy $\pi$ (and consequently the
        structural coordinates $(\mu^\pi,\alpha^\pi)$) so as to
        minimize this joint objective while remaining inside the
        admissible region determined by $S$ and
        $\mathcal{Q}_{\mathrm{adm}}(L)$.
  \item \textbf{Cognitive universe and bidirectional higher-order homotopy.}
        The optimized behavior of the scheme is represented by
        reading-path homotopy classes $[\gamma^{(\pi)}]$ in the
        paragraph–monoid network and by corresponding law-space
        homotopy classes in $B_L$:
        \begin{itemize}
          \item \emph{horizontally}, different PACER runs correspond
                to different paths in the\\ paragraph–monoid network,
                related by controlled homotopies as in
                Definition~\ref{def:pacer-reading-path-homotopy} and
                Remark~\ref{rem:horizontal-vertical-homotopy};
          \item \emph{vertically}, at each node the multi-layer HACA
                fiber (lexical layer $l_1$ plus higher semantic layers
                $l_{n\ge 2}$) carries its own higher-order homotopy
                structure (CH-layering, Jacobiator-gaps, LA/LB/LC),
                and the scheme uses the law-space connection
                $\GZTLC_L$ and semantic gauge field $\mathcal{A}_L$ to
                update the $n$-th order connections in response to the
                input.
        \end{itemize}
        The pair consisting of the horizontal homotopy class
        $[\gamma^{(\pi)}]$ and the vertical law-space homotopy data
        thus encodes the resulting ``cognitive state'' of the scheme
        within the HACA law-space universe.
\end{enumerate}
\end{definition}

\begin{remark}[Input $\to$ KAT decomposition $\to$ multi-layer GRL $\to$ optimized cognition]
\label{rem:haca-grl-cognitive-scheme-summary}
Informally, the HACA-based GRL cognitive scheme can be summarized as
the following pipeline:
\[
\begin{aligned}
  &\text{input (raw lexical data)}
    \;\longrightarrow\;
    \text{KAT-based heterogeneous decomposition}
    \;\longrightarrow\;\\
  &\qquad\text{multi-layer GRL path integrals (lexical + semantic)}
    \;\longrightarrow\;\\
  &\qquad\text{joint optimization over }(\mu_i,\alpha_i,\mathbf{v})
    \;\longrightarrow\;\\
  &\qquad\text{cognitive universe with bidirectional higher-order homotopy}.
\end{aligned}
\]
The first arrow corresponds to KAT-driven segmentation and PACER-based
lexical operations under a lexical minimal-action principle; the
second arrow corresponds to semantic GRL evaluation across law-space
layers; the third arrow corresponds to SAC-PACER-HACA optimization in
structural coordinate space; and the final arrow corresponds to the
update of the HACA law-space fiber bundle, with horizontal (chapter/
section) and vertical (multi-layer semantic) higher-order homotopies
as described in Remark~\ref{rem:horizontal-vertical-homotopy}. This is
the dynamic cognitive paradigm realized by HACA/PACER on top of the
structural HACA law-object definition in this volume.
\end{remark}

\chapter[SAC-PACER-HACA Engines and Value Geometry]{SAC-PACER-HACA Engines and Value Geometry:\\
         Objectives, Occupancy Measures, and MDQ}


In standard Soft Actor-Critic (SAC) architectures, the objective
combines expected return with an entropy bonus, encouraging policies
that are both high-reward and sufficiently exploratory. In the
SAC-PACER-HACA setting of this volume, the same basic philosophy is
lifted into HACA law-space: lexical GRL actions, semantic GRL path
integrals, occupancy measures, and value geometry vectors define an
objective $\mathcal{J}(\pi)$ and constraints that govern how policies
evolve on KAT action monoids under the guidance of HACA law-objects,
semantic gauge fields, and safe law-space sectors. This chapter
formalizes these objectives and constraints and embeds them into the
law-space geometry developed in previous parts.

\section[Objectives and Constraints for SAC-PACER-HACA]{Objectives and Constraints for SAC-PACER-HACA \\Engines}


In standard Soft Actor-Critic (SAC)~\cite{Haarnoja2018SAC}, the objective for a policy $\pi$ is of the form
\[
  \mathcal{J}_{\mathrm{SAC}}(\pi)
  =
  \mathbb{E}_\pi
  \Bigl[\,
    \sum_{t} r(s_t,a_t)
    + \alpha_{\mathrm{ent}}\,\mathcal{H}\bigl(\pi(\cdot\mid s_t)\bigr)
  \,\Bigr],
\]
where $r$ is a reward function, $\mathcal{H}$ is an entropy term, and
$\alpha_{\mathrm{ent}}\ge 0$ controls the exploration bonus. State
space, action space, and reward are typically defined in a task-
specific, black-box manner.

In the SAC-PACER-HACA engine, the functional roles of $r$ and
$\mathcal{H}$ are replaced by quantities defined on HACA law-space,
KAT action monoids, and semantic certificates. For a HACA law-object
$L$ with lexical KAT action monoid $\mathcal{M}$, semantic GRL path
integrals, occupancy measures $\mu_i^\pi$, strategy weights
$\alpha_i^\pi$, value geometry vector $\mathbf{v}$, and semantic MDQ
profile $\mathbf{q}^{\mathrm{sem}}(\pi)$, a typical SAC-PACER-HACA
objective can be written schematically as
\[
\begin{aligned}
  \mathcal{J}(\pi)
  &=
  \mathbb{E}_\pi
  \Bigl[
    S_{\mathrm{lex}}[w]
    + \lambda_{\mathrm{sem}}
      \sum_{\chi\in\mathcal{C}_L}
      \Phi_\chi\bigl(\mathcal{Z}^{\mathrm{sem}}_\chi[\gamma]\bigr)\\
  &\qquad
    + \lambda_{\mathrm{MDQ}}
      \sum_{i}\mu_i^\pi\,q_i^{\mathrm{sem}}(\pi)
    - \lambda_{\mathrm{val}}
      \sum_{i}\mu_i^\pi\,v_i
    - \alpha_{\mathrm{ent}}\,
      \mathcal{H}_{\mathrm{struct}}(\pi)
  \Bigr].
\end{aligned}
\]
Here:
\begin{itemize}
  \item $S_{\mathrm{lex}}[w]$ is the lexical GRL action on operator
        paths $w\in\mathsf{Paths}(\mathcal{M})$, as in
        Definition~\ref{def:lexical-grl-action-on-M};
  \item $\mathcal{Z}^{\mathrm{sem}}_\chi[\gamma]$ are semantic GRL
        path integrals on paragraph paths
        $\gamma\in\Omega^{(3)}(L)$ for certificate predicates
        $\chi\in\mathcal{C}_L$, with $\Phi_\chi$ extracting scalar
        penalties or rewards;
  \item $\mu_i^\pi$ are occupancy measures on KAT-based segments,
        $\mathbf{q}^{\mathrm{sem}}(\pi)=(q_i^{\mathrm{sem}}(\pi))_i$
        is the semantic MDQ profile, and
        $\mathbf{v}=(v_i)_i$ is the value geometry vector, as in
        Definition~\ref{def:semantic-mdq-profile} and
        Definition~\ref{def:pack-value-geometry};
  \item $\mathcal{H}_{\mathrm{struct}}(\pi)$ is a structural entropy
        or diversification term over KAT segments and semantic roles
        (for example, encouraging a spread of $\alpha_i^\pi$);
  \item $\lambda_{\mathrm{sem}},\lambda_{\mathrm{MDQ}},\lambda_{\mathrm{val}},
        \alpha_{\mathrm{ent}}\ge 0$ are hyperparameters controlling
        the trade-offs between these contributions.
\end{itemize}

In addition to minimizing $\mathcal{J}(\pi)$, SAC-PACER-HACA engines
must satisfy structural and safety constraints:
\begin{enumerate}
  \item \textbf{Law-space safety:} the induced law-space paths and
        their horizontal lifts in the lexical principal bundle must
        remain within safe sectors $S\subseteq B_L$ and within the
        baseline-admissible MDQ region
        $\mathcal{Q}_{\mathrm{adm}}(L)$
        (Definitions~\ref{def:safe-unsafe-sectors} and
        \ref{def:baseline-constrained-mdq-region}).
  \item \textbf{Certificate admissibility:} paragraph operators must
        belong to $\Omega^{(3)}_{\mathrm{adm}}(L)$, i.e.\ satisfy
        $\Xi_\chi(\gamma)=\mathrm{true}$ for all semantic certificates
        $\chi\in\mathcal{C}_L$
        (Definition~\ref{def:semantically-admissible-paragraph-operators}).
  \item \textbf{Baseline deviation bounds:} the deviation of the
        semantic-layer connection and curvature from the baseline
        geometry, $(\Delta{\GZLC}_L,\Delta{\GZLCurv}_L)$ as in
        Definition~\ref{def:baseline-deviation}, must remain within
        prescribed bounds along trajectories induced by~$\pi$.
\end{enumerate}



\section[Occupancy Measures and Gradient-Based Baselines]{Occupancy Measures and Gradient-Based Baselines}


In the law-space baseline chapter, we introduced occupancy measures
$\mu_i^\pi$ and value geometry vectors $\mathbf{v}=(v_i)_i$ for
HACA agents (Definitions~\ref{def:pack-occupancy-measures} and
\ref{def:pack-value-geometry}). In the SAC-PACER-HACA context, these
quantities play the role of gradient-based baselines for policy
optimization.

Recall that, for a policy $\pi$, the occupancy measure $\mu_i^\pi$
records how much time, attention, or frequency is devoted to the
$i$-th pack-level segment (for instance, a paragraph node or a KAT
submonoid), while $v_i$ measures the local sensitivity of the objective
$\mathcal{J}(\pi)$ to an infinitesimal increase in $\mu_i^\pi$.

\begin{definition}[Gradient-based value baselines]
\label{def:gradient-based-baseline}
Let $\mathcal{J}(\pi)$ be a differentiable objective for policies
$\pi\in\Pi(L)$, and assume that in a neighborhood of a reference
policy $\pi_0$ we can write
\[
  \mathcal{J}(\pi)
  \approx
  \widehat{\mathcal{J}}(\mu^\pi),
\]
where $\mu^\pi=(\mu_i^\pi)_{i\in I_L}$ is the occupancy vector on a
finite index set $I_L$ of segments. The \emph{gradient-based value
baseline} at $\pi_0$ is the vector
\[
  \mathbf{v}
  =
  (v_i)_{i\in I_L},
  \qquad
  v_i
  :=
  \left.
    \frac{\partial \widehat{\mathcal{J}}}{\partial \mu_i}
  \right|_{\mu=\mu^{\pi_0}}.
\]
\end{definition}

\begin{example}[Finite-dimensional toy model]
\label{ex:toy-model-mu-v}
Consider a simplified model with three segments $i=1,2,3$ (for
example, three paragraph clusters), and assume that the objective
depends only on their occupancies via
\[
  \widehat{\mathcal{J}}(\mu_1,\mu_2,\mu_3)
  =
  c_1\mu_1^2 + c_2\mu_2^2 + c_3\mu_3^2,
\]
with constants $c_i>0$ reflecting the ``cost'' of spending attention
in each segment. Then
\[
  v_i
  =
  \left.\frac{\partial\widehat{\mathcal{J}}}{\partial\mu_i}\right|_{\mu=\mu^{\pi_0}}
  =
  2c_i\,\mu_i^{\pi_0},
\]
and the gradient-based baseline points in the direction of reducing
occupancy in segments with high $c_i\,\mu_i^{\pi_0}$ values. In a
full SAC-PACER-HACA engine, the coefficients $c_i$ and occupancies
$\mu_i^{\pi_0}$ would themselves be derived from GRL path integrals,
MDQ, and semantic certificates.
\end{example}

In practice, SAC-PACER-HACA updates can be organized so that, at each
step, the actor receives gradient signals proportional to
\[
  -\frac{\partial\mathcal{J}}{\partial\mu_i},
  \qquad
  -\frac{\partial\mathcal{J}}{\partial\alpha_i},
\]
where the derivatives with respect to $\alpha_i^\pi$ capture how
changes in strategy weights affect the objective. These derivatives
are computed by backpropagating through the GRL path integrals and
MDQ profiles and combining the contributions of lexical actions,
semantic penalties, and value geometry.



\section[Embedding SAC-PACER-HACA into Law-Space Geometry]{Embedding SAC-PACER-HACA into Law-Space Geometry}


Policy updates in SAC-PACER-HACA can be viewed as inducing flows on
law-space. Let $\{\pi_\theta\}$ be a parametric family of policies
indexed by parameters $\theta\in\Theta\subseteq\mathbb{R}^d$. The
gradient flow
\[
  \frac{\mathrm{d}\theta}{\mathrm{d}t}
  = -\nabla_\theta\mathcal{J}(\pi_\theta)
\]
induces a flow on the space of occupancy measures, strategy weights,
and law-space trajectories, and, via $\Phi_L$, a flow on the law-space
sector $B_L\subseteq\GZLS$.

\begin{definition}[Law-space flow induced by SAC-PACER-HACA updates]
\label{def:lawspace-flow-sac-pacer-haca}
Let $\pi_\theta$ be a differentiable family of policies, and let
$(\mu^\theta,\alpha^\theta)$ denote the corresponding occupancy
measures and strategy weights. The SAC-PACER-HACA update
\[
  \frac{\mathrm{d}\theta}{\mathrm{d}t}
  = -\nabla_\theta\mathcal{J}(\pi_\theta)
\]
induces:
\begin{itemize}
  \item a flow on structural coordinates:
        \(
          \frac{\mathrm{d}\mu^\theta}{\mathrm{d}t},
          \frac{\mathrm{d}\alpha^\theta}{\mathrm{d}t};
        \)
  \item a flow on law-space paths $\Gamma_\theta:[0,T]\to B_L$ via the
        dependence of trajectories on $\pi_\theta$ and the projection
        $\Phi_L$;
  \item a flow on the HACA fiber bundle over $B_L$ through the
        triple-layer connection $\GZTLC_L$ and semantic gauge field
        $\mathcal{A}_L$.
\end{itemize}
We refer to the resulting vector field on $B_L$ (together with its
lift to the HACA bundle) as the \emph{law-space flow} induced by the
SAC-PACER-HACA engine.
\end{definition}

Soft constraints in SAC-PACER-HACA (for example, penalties on semantic
MDQ or deviations from value baselines) correspond to adding
smooth terms to $\mathcal{J}(\pi)$, thereby modifying the flow
vector field without introducing explicit hard boundaries. Hard
constraints (such as staying inside safe sectors or respecting strict
certificate admissibility) are enforced by restricting attention to
flows that remain within a feasible region of $B_L$ and of the HACA
fiber, or by adding barrier-type terms that blow up near the boundary.

\begin{remark}[Feasible regions and safe sectors]
\label{rem:feasible-region-safe-sectors}
Geometrically, we can think of the admissible region for SAC-PACER-HACA
as a subset
\[
  \mathcal{F}
  \subseteq
  \bigl(B_L\times\mathcal{M}\times\mathcal{P}(\Omega^{(3)}(L))\bigr),
\]
defined by safety, certificate, and baseline constraints. The law-space
flow induced by the SAC-PACER-HACA updates must remain within
$\mathcal{F}$; when a trajectory approaches the boundary of $\mathcal{F}$,
either the dynamic is adjusted (e.g.\ via projection or barrier
terms) or the policy is rejected. This picture provides the geometric
background for the certificate-based AI constructions in subsequent
chapters, where homotopy classes of safe flows and reading paths play
 a central role.
\end{remark}



\section[Connections to Standard Actor-Critic Architectures]{Connections to Standard Actor-Critic Architectures}


Structurally, SAC-PACER-HACA can be seen as a refinement of standard
actor–critic and SAC architectures along three axes:
\begin{enumerate}
  \item \textbf{State representation:} replacing black-box state
        spaces with HACA law-objects and law-space sectors $B_L$,
        together with their multi-layer semantic structure.
  \item \textbf{Action representation:} replacing generic actions with
        KAT operators in lexical action monoids $\mathcal{M}$ and
        paragraph-level KAT monoid morphisms $F_i:M(P_i)\to M(P_{i+1})$.
  \item \textbf{Objective and constraints:} replacing simple reward +
        entropy objectives with objectives built from lexical and
        semantic GRL path integrals, MDQ profiles, value geometry, and
        certificate-based safety constraints.
\end{enumerate}

At the implementation level, many underlying components of a standard
SAC system can remain unchanged: neural networks parameterizing
policies and value functions, replay buffers, and gradient-based
optimization procedures. The main difference lies in how states and
actions are interpreted and how the objective and constraints are
constructed.

\begin{itemize}
  \item In standard SAC, the state $s_t$ is typically a vector of
        observations, and actions $a_t$ are drawn from a parametric
        distribution $\pi_\theta(a_t\mid s_t)$.
  \item In SAC-PACER-HACA, the ``state'' is a point in the HACA
        law-space fiber: it includes information about the current
        paragraph node $v_i$, the current law-space sector, and the
        semantic certificate status. The ``action'' is a choice of
        KAT operator (or composition of operators) in $\mathcal{M}$
        and a choice of transition between paragraph nodes.
  \item The critic no longer estimates value purely as a function of
        state and action, but as a function of structural coordinates
        $(\mu^\pi,\alpha^\pi)$ and law-space position, with MDQ and
        value geometry serving as key features.
\end{itemize}

\begin{remark}[Incremental integration into existing systems]
\label{rem:incremental-integration}
From an engineering standpoint, it is possible to introduce the
SAC-PACER-HACA structure gradually into existing actor–critic systems:
\begin{enumerate}
  \item Start by interpreting a subset of actions as KAT-like
        operators (e.g.\ edit operations on text) and a subset of
        states as HACA-like law-space slices (e.g.\ paragraph-level
        representations), while keeping the rest of the system
        unchanged.
  \item Introduce semantic certificates and safe sectors as additional
        constraints on trajectories, but initially only as soft
        penalties in the objective.
  \item Gradually extend the representation to a full HACA law-object,
        refine the power submonoid spectrum of $\mathcal{M}$, and
        strengthen safety constraints from soft penalties to full
        certificate-based admissibility checks.
\end{enumerate}
This incremental route allows one to leverage existing actor–critic
infrastructure while progressively incorporating the law-space and
certificate-based structure developed in this volume.
\end{remark}



\chapter[Jacobiator-Gap Certificates for HACA and PACER]{Jacobiator-Gap Certificates for HACA and PACER:\\
         Heterogeneous Property-Changing Morphisms and Law-Space Invariants}


In the pure mathematics volume of the GaoZheng G-Framework, Jacobiator-
gaps arise as obstructions to strict Lie-type identities in generalized
noncommutative Lie algebras (PFB-GNLA) and related $L_\infty$-structures:
the Jacobiator measures the three-fold failure of associativity in
brackets, and its cohomology class provides a law-space invariant.
In the HACA/PACER setting of the present volume, we reinterpret this
notion on the property layer for HACA law-objects and use it as the
foundation for certificate structures: Jacobiator-gap data, together
with their higher homotopies, become the basic invariants behind
law-space certificates for cognitive and application systems.

\section[Jacobiator-Gap as a Law-Space Invariant in HACA]{Jacobiator-Gap as a Law-Space Invariant in HACA}


In the pure mathematics volume, the Jacobiator-gap is defined at the
structural level for generalized Lie brackets and $L_\infty$-operations:
given a triple of operations $(\cdot,\cdot,\cdot)$, the Jacobiator
measures the failure of the Jacobi identity and determines a class in
a third cohomology group. When this class is trivial (i.e.\ the
Jacobiator-gap is a coboundary), higher homotopies can be chosen to
``fill'' the gap, and the resulting structures admit well-behaved
higher connections and bundles.

In the present HACA volume, we mainly use the \emph{property-layer}
version of this notion, introduced in
Definition~\ref{def:haca-jacobiator-gap-property-layer}. For HACA
law-objects $L_i$ and heterogeneous property-changing morphisms
$(f_{ij},S_{ij})$ as in
Definition~\ref{def:property-changing-operators}, the property-layer
Jacobiator-gap
\[
  J_{\mathrm{prop}}
  =
  S_{23}\circ S_{12} - S_{13}
\]
measures the failure of property-level data to be compatible under
three-step compositions. This is a HACA-specific reinterpretation of
the Jacobiator-gap: instead of measuring non-closure of brackets, it
measures non-closure of heterogeneous property-changing data.

\begin{definition}[Property-layer Jacobiator-gap and law-space invariance]
\label{def:haca-jacobiator-gap-invariant}
Let $L_1,L_2,L_3$ be HACA law-objects with law-supports
$B_{L_1},B_{L_2},B_{L_3}\subseteq\GZLS$, and consider heterogeneous\\
property-changing morphisms
\[
  (f_{12},S_{12}):L_1\to L_2,
  \quad
  (f_{23},S_{23}):L_2\to L_3,
  \quad
  (f_{13},S_{13}):L_1\to L_3.
\]
The \emph{property-layer Jacobiator-gap} at $(L_1,L_2,L_3)$ is the
defect
\[
  J_{\mathrm{prop}}
  :=
  S_{23}\circ S_{12} - S_{13}
\]
as in Definition~\ref{def:haca-jacobiator-gap-property-layer}. Two
such triples of heterogeneous morphisms are called \emph{law-space
equivalent} if they are related by homotopies and refinements that
preserve the induced law-space paths in
$B:=B_{L_1}\cup B_{L_2}\cup B_{L_3}\subseteq\GZLS$ and respect the
HACA semantic and safety constraints (in particular, remain within
safe sectors and baseline-admissible MDQ domains). The law-space
equivalence class of $J_{\mathrm{prop}}$ under such transformations
is called the \emph{property-layer Jacobiator-gap invariant} of the
triple $(L_1,L_2,L_3)$.
\end{definition}

Intuitively, the Jacobiator-gap invariant records the obstruction to
reconciling the two ways of transporting property data from $L_1$ to
$L_3$: via the two-step chain $L_1\to L_2\to L_3$ or via the direct
morphism $L_1\to L_3$. When the invariant is trivial (i.e.\
law-space equivalent to zero), higher homotopy data can be chosen to
fill the gap, and the corresponding heterogeneous property-changing
data can be organized into well-behaved bundles.

\begin{definition}[Homotopy admissibility and Jacobiator-gap coboundary]
\label{def:homotopy-admissible-jgap}
A finite chain of HACA law-objects
\[
  L_0\to L_1\to\cdots\to L_n
\]
with heterogeneous property-changing morphisms is called
\emph{homotopy admissible} if, for every triple $(L_i,L_{i+1},L_{i+2})$
along the chain, the corresponding property-layer Jacobiator-gap
$J_{\mathrm{prop}}$ admits homotopy-filling data $H_i$ satisfying the
conditions of
Definition~\ref{def:haca-jacobiator-gap-property-layer} and
Proposition~\ref{prop:gaugeable-operator-bundle-from-property-intersections}
(non-empty intersections, stability, and coboundary triviality). In
this case, the collection $\{J_{\mathrm{prop}},H_i\}$ defines a
trivial class in the relevant third cohomology group over the union
of law-sectors, and we say that the Jacobiator-gap along the chain is
a \emph{coboundary}.
\end{definition}

In the homotopy admissible case,
Proposition~\ref{prop:gaugeable-operator-bundle-from-property-intersections}
shows that the union of property intersections supports a principal
operator bundle $\pi_{\mathrm{op}}:P_{\mathrm{op}}\to\mathcal{U}$,
whose fibers encode admissible operators at each property configuration
and whose gaugeability is controlled by the LA/LB/LC homotopy
completeness levels. The existence of such a bundle, together with
the triviality of the Jacobiator-gap invariant, is what allows us to
define law-space certificates for HACA and PACER.



\section[Certificate Bundles for HACA/PACER]{Certificate Bundles for HACA/PACER}


Building on the operator bundle $P_{\mathrm{op}}\to\mathcal{U}$ from
Proposition~\ref{prop:gaugeable-operator-bundle-from-property-intersections},
we now introduce certificate bundles for HACA/PACER. Roughly, these
bundles collect all admissible homotopy-filling data for Jacobiator-
gaps along law-space paths, together with their compatibility
conditions, and organize them as fibers over a base of law-objects,
properties, and path data.

\begin{definition}[Certificate base space for HACA/PACER]
\label{def:certificate-base-space}
Let $L$ be a HACA law-object and let $\Omega^{(3)}_{\mathrm{adm}}(L)$
denote the set of semantically admissible paragraph paths as in
Definition~\ref{def:semantically-admissible-paragraph-operators}. The
\emph{certificate base space} $B_{\mathrm{cert}}(L)$ is defined as the
set of tuples
\[
  b = \bigl(L_0,\dots,L_n; \gamma; \mathcal{P}\bigr)
\]
where:
\begin{itemize}
  \item $L_0,\dots,L_n$ is a homotopy admissible chain of HACA law-
        objects (Definition~\ref{def:homotopy-admissible-jgap});
  \item $\gamma\in\Omega^{(3)}_{\mathrm{adm}}(L)$ is a paragraph
        reading path compatible with this chain in the sense that its
        law-space projection lies in the union
        $B:=\bigcup_i B_{L_i}\subseteq\GZLS$ and satisfies the
        safety and MDQ constraints of
        Definition~\ref{def:safe-unsafe-sectors} and
        Definition~\ref{def:baseline-constrained-mdq-region};
  \item $\mathcal{P}$ denotes the relevant property-layer data along
        the chain (property sets, property intersections, and induced
        Jacobiator-gaps).
\end{itemize}
\end{definition}

\begin{definition}[Certificate bundles and fibers]
\label{def:certificate-bundle}
For each $b\in B_{\mathrm{cert}}(L)$, let $\mathcal{H}_b$ denote the
set of all higher-homotopy and smoothing data along the chain
$(L_0,\dots,L_n)$ that:
\begin{itemize}
  \item fill all local property-layer Jacobiator-gaps as in
        Definition~\ref{def:haca-jacobiator-gap-property-layer};
  \item stabilize the property intersections and admissible operator
        fibers so that a principal operator bundle
        $P_{\mathrm{op}}\to\mathcal{U}$ exists over the union of
        intersections (as in
        Proposition~\ref{prop:gaugeable-operator-bundle-from-property-intersections});
  \item are compatible with the semantic certificates $\mathcal{C}_L$
        and with the reading-path admissibility of $\gamma$.
\end{itemize}
Define
\[
  P_{\mathrm{cert}}
  :=
  \bigsqcup_{b\in B_{\mathrm{cert}}(L)} \mathcal{H}_b,
\]
with projection
\[
  \pi_{\mathrm{cert}} : P_{\mathrm{cert}}\longrightarrow B_{\mathrm{cert}}(L),
  \qquad
  h\mapsto b \text{ if }h\in\mathcal{H}_b.
\]
We call $\pi_{\mathrm{cert}}:P_{\mathrm{cert}}\to B_{\mathrm{cert}}(L)$
the \emph{certificate bundle} for HACA/PACER on~$L$. Elements of
$\mathcal{H}_b$ are called \emph{certificates} at~$b$.
\end{definition}

The LA/LB/LC homotopy completeness levels
(Definition~\ref{def:homotopy-completeness-levels}) admit a natural
interpretation in this framework:
\begin{itemize}
  \item \textbf{Level A (local completeness):} along a single
        homotopy-complete path, the restriction of the certificate
        bundle is trivial, and certificates can be chosen locally
        along the path without obstruction.
  \item \textbf{Level B (common-core completeness):} when there exists
        a homotopy common core $\mathcal{C}\subseteq B_{\mathrm{cert}}$
        through which all relevant paths pass, the restriction of
        $P_{\mathrm{cert}}$ to $\mathcal{C}$ admits a common local
        trivialization; certificates can be compared and glued in this
        region, providing a shared ``core chart'' for auditing and
        alignment.
  \item \textbf{Level C (normative completeness):} when there exists a
        normative subspace $\mathcal{N}\subseteq B_{\mathrm{cert}}$
        such that all admissible paths can be projected into
        $\mathcal{N}$, the restriction $P_{\mathrm{cert}}|_{\mathcal{N}}$
        is globally trivializable (up to the chosen baseline). One can
        then choose a global certificate or canonical baseline path
        to which all good paths are semantically homotopy-equivalent.
\end{itemize}

\begin{remark}[Interaction with paragraph--KAT homotopy networks]
\label{rem:certificate-bundle-homotopy-network}
The certificate bundle interacts naturally with the paragraph--KAT
homotopy network constructed in the previous chapter. Each PACER
reading path $\gamma$ in the paragraph–monoid network determines a
point $b\in B_{\mathrm{cert}}(L)$ via its associated law-objects and
property-layer data, and a certificate in $\mathcal{H}_b$ specifies
the Jacobiator-gap homotopies and operator-bundle gauges needed to
regard $\gamma$ as a certified path. In particular, the
reading-path certificates of
Definition~\ref{def:reading-path-certificate} can be seen as sections
(or partial sections) of the certificate bundle along families of
reading paths. This provides the geometric backbone for
certificate-based AI in the HACA/PACER setting.
\end{remark}



\section[Example: Pathology-to-Reverse-Design Pipelines Revisited]{Example: Pathology-to-Reverse-Design Pipelines Revisited\\
         (LBOPB-Based Law-Objects from Volume II)}


\paragraph{Cross-Volume Setup.}
Let
\[
  L_{\mathrm{path}},\;
  L_{\mathrm{cell}},\;
  L_{\mathrm{met}},\;
  L_{\mathrm{tox}},\;
  L_{\mathrm{PK/PD}},\;
  L_{\mathrm{rev\mbox{-}drug}}
\]
denote the LBOPB-based law-objects for pathology, cellular/immune response,
metabolic and signalling processes, toxicity constraints, PK/PD envelopes
and reverse drug design, respectively, as constructed in Applications
Volume~II (see the chapters devoted to LBOPB and PGOM for detailed
definitions). In the present volume we only use these $L_{\bullet}$ as
law-objects inside the HACA law-space, together with property layers
and property-changing operators.

\paragraph{Property Layers and Heterogeneous Morphisms.}
For each $L_\bullet$ above, we consider a property set
$\mathrm{Prop}(L_\bullet)$ encoding the medically relevant invariants and
constraints (pathological markers, target mechanisms, structural and
toxicity bounds, PK/PD envelopes, etc.), as in
Definition~\ref{def:haca-property-layer}. Between consecutive stages we
choose heterogeneous property-changing morphisms
\[
  (f_{\mathrm{path}\to\mathrm{cell}}, S_{\mathrm{path}\to\mathrm{cell}}),
  \quad
  (f_{\mathrm{cell}\to\mathrm{met}}, S_{\mathrm{cell}\to\mathrm{met}}),
  \quad
  \dots,
  \quad
  (f_{\mathrm{PK/PD}\to\mathrm{rev\mbox{-}drug}},
   S_{\mathrm{PK/PD}\to\mathrm{rev\mbox{-}drug}})
\]
in the sense of Definition~\ref{def:property-changing-operators}. These
morphisms couple law-space evolution with property-level transformations
for the LBOPB pipeline.

\paragraph{Local Jacobiator-Gaps and Homotopy Filling.}
For each triple of consecutive stages, for instance
\[
  (L_{\mathrm{path}},L_{\mathrm{cell}},L_{\mathrm{met}}),
  \quad
  (L_{\mathrm{cell}},L_{\mathrm{met}},L_{\mathrm{tox}}),
  \quad
  \dots,
\]
we can form the property-layer Jacobiator-gap
$J_{\mathrm{prop}}$ as in
Definition~\ref{def:haca-jacobiator-gap-property-layer}. Homotopy data
$H_k$ on the property layer fill these local gaps precisely when the
corresponding triple intersections
\[
  \mathrm{Prop}(L_i)
  \cap S_{i\to i+1}^{-1}\bigl(\mathrm{Prop}(L_{i+1})\bigr)
  \cap (S_{i+2\leftarrow i+1}\circ S_{i\to i+1})^{-1}
       \bigl(\mathrm{Prop}(L_{i+2})\bigr)
  \neq \varnothing
\]
are non-empty for each triple $(L_i,L_{i+1},L_{i+2})$ in the pipeline.
Thus, local Jacobiator-gaps are removed by homotopy-filling on the property
layer.

\paragraph{Global Homotopy Paths and GRL Optimization.}
Collecting the LBOPB-related operator monoids from Volume~II into a finite
family
\[
  M_1,\dots,M_7
\]
(for example, corresponding to pathology, immunity, metabolism, signalling,
toxicity, PK and PD), we can regard the entire LBOPB pipeline as a
law-space chain inside the HACA framework. A higher-homotopy path
\[
  \gamma:\;
  L_{\mathrm{path}}
  \xRightarrow[\;M_1,\dots,M_7\;]{\text{homotopy evolution}}
  L_{\mathrm{rev\mbox{-}drug}}
\]
that lifts this chain to the property layer and remains compatible with all
local constraints provides a globally coherent evolution from pathology to
reverse drug design.

In a GRL formulation, such a path $\gamma$ can be equipped with an action
functional $\mathcal{S}(\gamma)$; dually, one can consider a policy
$\pi$ and a performance functional $\mathcal{J}(\pi)$ defined on LBOPB
operator paths. Solving the higher-homotopy equation for $\gamma$ or
equivalently optimizing the GRL path integral
\[
  \gamma^\star = \arg\min_{\gamma}\, \mathcal{S}(\gamma)
  \quad\text{or}\quad
  \pi^\star = \arg\max_{\pi}\, \mathcal{J}(\pi)
\]
yields a homotopy-optimal law-space trajectory from pathology to reverse
drug design. This trajectory:
\begin{itemize}
  \item fills the local Jacobiator-gaps along the LBOPB chain on the
        property layer; and
  \item realizes a globally consistent evolution across the LBOPB operator
        monoids imported from Volume~II.
\end{itemize}

\paragraph{Certificates and Cross-Volume Interpretation.}
In the present volume, we interpret such a homotopy-optimal trajectory as a
\emph{certificate} for the LBOPB pipeline: it certifies that the underlying
life-science law-objects and operators (constructed in Volume~II) admit a
HACA-compatible, property-layer-consistent evolution from pathology to
reverse drug design. This makes the notion of ``Jacobiator-gap free'' and
homotopy completeness operational not only for cognitive/linguistic
systems (HACA and PACER), but also for life-science applications via
cross-volume law-space coupling.

\begin{remark}[Relation to certificate bundles and LA/LB/LC levels]
\label{rem:lbopb-cert-bundle-levels}
From the viewpoint of the certificate bundle
$\pi_{\mathrm{cert}}:P_{\mathrm{cert}}\to B_{\mathrm{cert}}(L)$ in
Definition~\ref{def:certificate-bundle}, the LBOPB pipeline corresponds
to a subset $B_{\mathrm{cert}}^{\mathrm{LBOPB}}\subseteq B_{\mathrm{cert}}(L)$
consisting of tuples encoding the pathology-to-reverse-design chain
and its property-layer data. A homotopy-optimal trajectory $\gamma^\star$
as above provides a section (or partial section) of $P_{\mathrm{cert}}$ over
$B_{\mathrm{cert}}^{\mathrm{LBOPB}}$, i.e.\ a consistent choice of
Jacobiator-gap fillings and gauges along the entire pipeline. When this
section satisfies Level~A/B/C completeness conditions (local along the
pipeline, around a common core, or on a normative subspace), the LBOPB
pipeline is said to enjoy LA/LB/LC certificate completeness respectively.
In particular, a Level~C-complete LBOPB corridor can serve as a normative
template for future pipelines, anchoring cross-volume law-space reasoning
between Volumes~II and~III.
\end{remark}





\part[Certificate-Based AI, Interpretability, and Governance]{Certificate-Based AI, Interpretability, and Governance\\
      in HACA and PACER Systems}

\chapter[Certificate-Based AI for Hierarchical Cognitive Architectures]{Certificate-Based AI for Hierarchical Cognitive Architectures:\\
         Law-Space Certificates and MDQ-Based Evaluation}


\section[Law-Space Certificates and MDQ-Based Evaluation]{Law-Space Certificates and MDQ-Based Evaluation}


In the preceding parts we have introduced all the basic ingredients of
law-space certificates for HACA/PACER:
\begin{itemize}
  \item property-layer Jacobiator-gaps and their higher homotopies,
        together with homotopy admissible chains of HACA law-objects
        (Definitions~\ref{def:haca-jacobiator-gap-property-layer},
        \ref{def:haca-jacobiator-gap-invariant},
        \ref{def:homotopy-admissible-jgap});
  \item principal operator bundles over property intersections and
        their gaugeability as controlled by LA/LB/LC levels
        (Proposition~\ref{prop:gaugeable-operator-bundle-from-property-intersections}
        and Definition~\ref{def:homotopy-completeness-levels});
  \item certificate bundles $\pi_{\mathrm{cert}}:P_{\mathrm{cert}}\to
        B_{\mathrm{cert}}(L)$ whose fibers consist of admissible
        Jacobiator-gap fillings and gauge choices
        (Definition~\ref{def:certificate-bundle});
  \item pack-level occupancy measures $\mu_i^\pi$, value geometry
        vectors $\mathbf{v}$, strategy weights $\alpha_i^\pi$, and
        semantic MDQ profiles $\mathbf{q}^{\mathrm{sem}}(\pi)$
        (Definitions~\ref{def:pack-occupancy-measures},
         \ref{def:pack-value-geometry},
         \ref{def:mu-alpha-structural-coordinates},
         \ref{def:semantic-mdq-profile});
  \item semantic GRL path integrals
        $\mathcal{Z}^{\mathrm{sem}}_\chi[\gamma]$ and predicates
        $\Xi_\chi(\gamma)$ on paragraph paths, together with lexical
        GRL actions $S_{\mathrm{lex}}[w]$ on KAT action monoids
        (Definitions~\ref{def:lexical-grl-action-on-M},
         \ref{def:semantic-grl-path-integral},
         \ref{def:semantically-admissible-paragraph-operators});
  \item reading-path certificates that tie together law-space paths,
        occupancy measures, value geometry, and semantic MDQ
        (Definition~\ref{def:reading-path-certificate}).
\end{itemize}
This section explains how these elements can be organized into a
certificate-based AI evaluation pipeline for HACA/PACER systems.

\begin{definition}[Law-Space Certificate Extraction for a PACER Run]
\label{def:lawspace-certificate-extraction}
Let $L$ be a HACA law-object equipped with a PACER engine in the sense
of Definition~\ref{def:pacer-engine-haca}, and let $\pi$ be a policy
in the associated policy family $\Pi(L)$. A \emph{law-space certificate
extraction} for a run of $\pi$ on an input document (or stream) is a
procedure that, given a runtime trace $\mathsf{Log}(\pi)$, constructs
a \emph{certificate record}
\[
  \mathsf{Cert}(\pi)
\]
consisting of the following components:
\begin{enumerate}
  \item A reconstructed lexical operator path $w\in\mathsf{Paths}(\mathcal{M})$
        and paragraph reading path $\gamma\in\Omega^{(3)}(L)$ associated
        with the run, together with their law-space projection
        $\Gamma:[0,T]\to B_L$ via $\Phi_L$.
  \item A reading-path certificate in the sense of
        Definition~\ref{def:reading-path-certificate}, i.e.\ law-space
        homotopy data and structural coordinates
        $(\mathbf{q}^{\mathrm{sem}}(\pi),\mathbf{v},\mu^\pi,\alpha^\pi)$
        along $\gamma$ that satisfy the admissibility conditions
        (safe sectors, MDQ bounds, semantic certificates).
  \item A collection of property-layer Jacobiator-gap invariants and
        their homotopy-filling data along the law-space path $\Gamma$,
        viewed as elements of the certificate bundle
        $P_{\mathrm{cert}}$ over appropriate points
        $b\in B_{\mathrm{cert}}(L)$.
  \item Aggregated statistics derived from the above, such as LA/LB/LC
        completeness levels (on which segments the certificate bundle
        is locally, regionally, or globally trivializable), MDQ
        budgets per segment, and value geometry comparisons between
        different runs or versions.
\end{enumerate}
The resulting $\mathsf{Cert}(\pi)$ is called a \emph{law-space
certificate record} for the given run.
\end{definition}

Conceptually, the certificate extraction pipeline can be viewed as
having four stages:
\begin{enumerate}
  \item \textbf{From logs to paths:} process runtime logs to reconstruct
        KAT operator paths $w$, paragraph paths $\gamma$, and law-space
        paths $\Gamma$.
  \item \textbf{From paths to certificates:} compute lexical and
        semantic GRL path integrals, occupancy measures, value
        gradients, MDQ profiles, and reading-path certificates.
  \item \textbf{From certificates to bundle data:} lift the path-level
        certificates to the certificate bundle $P_{\mathrm{cert}}$ by
        recording Jacobiator-gap fillings and LA/LB/LC completeness
        information.
  \item \textbf{From bundle data to evaluation:} aggregate the above
        into human-readable and \\machine-actionable summaries for
        evaluation tasks such as version comparison, auditing, and
        rollback decisions.
\end{enumerate}

\begin{definition}[Certificate-Based Evaluation Tasks]
\label{def:certificate-based-evaluation-tasks}
Given law-space certificate records \\$\mathsf{Cert}(\pi)$ for one or
more runs of HACA/PACER engines on a law-object $L$, we define the
following archetypal evaluation tasks:
\begin{itemize}
  \item \emph{Version-to-version comparison.} Given two policies
        $\pi^{(1)}$ and $\pi^{(2)}$ (for example, two model versions)
        and their certificate records $\mathsf{Cert}(\pi^{(1)})$ and
        $\mathsf{Cert}(\pi^{(2)})$, compare:
        \begin{enumerate}
          \item their reading-path homotopy classes;
          \item the distribution of occupancy measures and strategy
                weights over segments;
          \item changes in value geometry vectors and semantic MDQ
                profiles;
          \item changes in Jacobiator-gap certificates and LA/LB/LC
                completeness levels.
        \end{enumerate}
        The outcome is a structured report indicating where the newer
        policy behaves identically up to controlled homotopy and where
        genuine changes in law-space behavior occur.
  \item \emph{Audit evidence generation.} For a given run of $\pi$ and
        a given output (e.g.\ a specific answer or decision), extract
        the segment of $\mathsf{Cert}(\pi)$ corresponding to the
        relevant law-space path and reading-path homotopy class, and
        present:
        \begin{enumerate}
          \item the KAT operator sequence and paragraph path involved;
          \item the semantic certificates $\Xi_\chi(\gamma)$ that were
                satisfied;
          \item the local Jacobiator-gap fillings and MDQ/resource
                usage along the path.
        \end{enumerate}
        This provides a reproducible, law-space–based explanation for
        the behavior of the system.
  \item \emph{Rollback and gating decisions.} Given certificate records
        for policies deployed in A/B testing or progressive rollouts,
        use changes in certificate data (e.g.\ increased MDQ costs,
        new unsafe sectors visited, weakened LA/LB/LC completeness) as
        objective signals for gating or rolling back deployments.
\end{itemize}
\end{definition}

\begin{remark}[Automation vs Human Oversight in Certificate Pipelines]
\label{rem:automation-vs-human-oversight}
In principle, most of the computational steps in
Definition~\ref{def:lawspace-certificate-extraction} can be automated:
reconstructing paths from logs, computing GRL path integrals and MDQ
profiles, lifting data to certificate bundles, and generating summary
statistics are all algorithmic tasks. However, interpreting the
resulting certificates and deciding how to act on them typically
require human expertise:
\begin{itemize}
  \item Domain experts may decide which semantic certificates
        $\mathcal{C}_L$ are relevant for a given deployment context
        (e.g.\ safety, fairness, regulatory compliance).
  \item Governance bodies may decide what thresholds on MDQ, value
        geometry changes, or LA/LB/LC completeness levels should
        trigger alarms, audits, or rollback actions.
  \item Investigators may examine individual certificate records to
        understand rare or high-stakes behaviors in detail.
\end{itemize}
Thus, the certificate-based AI framework proposed in this volume is
intentionally hybrid: it provides a mathematically structured pipeline
for extracting and comparing law-space certificates, while leaving
normative judgments and policy decisions to human governance layers.
\end{remark}



\section[Semantic Law-Space vs Corpus Semantics in Neural LMs]{Semantic Law-Space in HACA/PACER versus \\Corpus-Based Semantics in Neural Language Models}
\label{subsec:haca-vs-neural-lm}

In order to position the semantic law-space geometry of HACA/PACER within
the broader landscape of contemporary machine learning, it is useful to
contrast it with the semantics implicitly realized by large-scale
autoregressive transformer-based language models. The aim of this
subsection is not to provide a survey of specific architectures, but to
clarify, at a structural level, how a law-space–based semantics differs
from corpus-driven distributional semantics.



\paragraph{Operator Universes versus Token Sequences.}

In the HACA/PACER setting, each law-object $L$ comes equipped with a
multi-scale operator universe
\[
  \Omega(L)
  =
  \Omega^{(0)}(L)\ \sqcup\ \Omega^{(1)}(L)\ \sqcup\ \Omega^{(2)}(L)\ \sqcup\ \Omega^{(3)}(L),
\]
as in Definition~\ref{def:haca-operator-universe}, where
$\Omega^{(0)}(L)$ collects character- or token-level operators,
$\Omega^{(1)}(L)$ word- or pack-level operators, $\Omega^{(2)}(L)
$ sentence-level operators, and $\Omega^{(3)}(L)$ paragraph- or
reading-path operators (KAT path integrals). Each layer carries an
explicit algebraic structure: free monoids with Kleene algebra with
tests (KAT) constraints at the word level, logical and dependency
structures at the sentence level, and graph-based path monoids at the
paragraph level.

By contrast, in large-scale autoregressive transformer-based language
models, the basic object space is the set $\mathcal{X}^\ast$ of all
finite sequences over a fixed token set $\mathcal{X}$. In principle, the
same sequence space can represent characters, words, sentences, and
paragraphs, but the distinctions between these levels are not built into
the object layer as different operator types. Rather, they are absorbed
implicitly into position encodings and learned statistical patterns over
$\mathcal{X}^\ast$. No explicit multi-scale operator hierarchy with
separate algebraic laws is exposed at the model interface.



\paragraph{Semantic Structure versus Distributed Representations.}

On top of $\Omega(L)$, \\HACA/PACER assumes an explicit semantic structure
space $\mathcal{S}(L)$ together with a semantic map
\[
  \operatorname{Sem}_L:\Omega(L)\longrightarrow \mathcal{S}(L).
\]
The space $\mathcal{S}(L)$ carries a family of explicitly specified
semantic operations
\[
  \mathrm{op}_S : \mathcal{S}(L)^n \longrightarrow \mathcal{S}(L),
\]
including, but not limited to, KAT composition, control-flow joins, and
graph-based composition of paragraph nodes and reading paths. Semantic
predicates are modeled as truth-valued functors
\[
  \chi:\mathcal{S}(L)\to\{\mathrm{true},\mathrm{false}\}
\]
that are functorially compatible with these operations and whose
``true'' fibers are closed under them
(Definition~\ref{def:haca-semantic-predicates}). Consequently, the
semantic property classes
\[
  P_\chi
  =
  \bigl\{\, \omega\in\Omega(L)\ \big|\ \chi(\operatorname{Sem}_L(\omega))
  = \mathrm{true} \,\bigr\}
\]
are algebraically and homotopically stable under the admissible
operations of HACA/PACER.

In transformer-based language models, by contrast, the semantic content
of a sequence $x\in\mathcal{X}^\ast$ is encoded in a high-dimensional
representation $h_\theta(x)\in E_\theta$ obtained via a deep stack of
linear and nonlinear layers. Various semantic tests (such as sentiment
classification, topic detection, or safety filters) are implemented via
functions
\[
  g:E_\theta \to Y,
\]
where $Y$ may be a finite label set or a probability simplex. However,
the internal operations on $E_\theta$ are not given by an explicit
algebra with axioms akin to KAT; instead, they result from the
composition of learned affine maps and pointwise nonlinearities. There is
no requirement that semantic tests $g$ be functorially compatible with a
fixed, externally specified family of semantic operations, nor that the
``true'' sets $g^{-1}(\{\mathrm{true}\})$ be closed under any
interpretable semantic composition.



\paragraph{Law-Space Geometry versus Implicit Objective Distributions.}

A key structural distinction lies at the level of law-space. In the
HACA/PACER framework, the semantic map $\operatorname{Sem}_L$ is
followed by a law-space projection
\[
  \operatorname{Law}_L:\mathcal{S}(L)\longrightarrow B_L\subseteq\GZLS,
\]
where $B_L$ is a sector of the GaoZheng law-space $\GZLS$ supporting the
given HACA/PACER configuration. The composite
\[
  \Phi_L:=\operatorname{Law}_L\circ\operatorname{Sem}_L
  :\Omega(L)\to B_L
\]
associates to each operator configuration an explicit law-space footprint.
Structural properties are realized as pullbacks of law-space sectors
\[
  P_\Sigma
  :=
  \Phi_L^{-1}(\Sigma)\subseteq\Omega(L),
\]
and observable property families are precisely those subsets of
$\Omega(L)$ that are both semantically observable and realizable as such
law-space pullbacks (Definition~\ref{def:haca-observable-properties}).
Under the conditions spelled out earlier, structural properties and
structural laws become two equivalent descriptions of the same GZ-LS data
(Remark~\ref{rem:structural-properties-as-laws}), and homotopy-based
notions such as Jacobiator-gaps and LA/LB/LC levels can be attached to
property classes as certificate data.

In contrast, contemporary large-scale autoregressive transformer-based
language models do not make use of an explicit law-space. Their training
objective is typically to approximate a conditional distribution
\[
  p_\theta(x_t\mid x_{<t})
  \approx
  p_{\text{data}}(x_t\mid x_{<t})
\]
over token sequences, where $p_{\text{data}}$ is implicitly defined by a
corpus. Any ``rules'' or ``constraints'' that emerge are encoded as
statistical regularities within this distribution and its parameter
space. There is no separate geometric space of laws, no explicit notion
of a law-space sector $B_L\subseteq\GZLS$, and no associated connection
or curvature to which one could attach homotopy-based certificates.
Properties such as safety, consistency, or task adherence are then
assessed via additional discriminative heads or post-hoc evaluation,
rather than being encoded as law-space sectors with associated
certificates.



\paragraph{Implications for Certificate-Based AI.}

From the standpoint of certificate-based AI, the HACA/PACER construction
thus differs from corpus-based semantic modeling in three structural
respects:
\begin{enumerate}
  \item The operator layer is multi-scale and algebraically typed
        ($\Omega^{(0)}\sqcup\Omega^{(1)}\sqcup\Omega^{(2)}\sqcup\Omega^{(3)}$),
        rather than a single undifferentiated sequence space.
  \item The semantic layer carries explicit operations and
        truth-valued functors whose ``true'' sectors are closed under
        these operations, allowing semantic properties to be realized as
        homotopy-stable classes.
  \item The law-space layer provides a geometric realization of these
        properties as sectors of $\GZLS$ equipped with connection and
        curvature data, so that structural properties and structural laws
        can be certified via Jacobiator-gap and LA/LB/LC invariants.
\end{enumerate}
Large-scale neural language models share with HACA/PACER the general
pattern of mapping token sequences to internal representations and then
evaluating properties, but they do so without an explicit multi-scale
operator hierarchy, without externally specified semantic operations
with closure conditions, and without a separate law-space geometry. The
HACA/PACER formalism can therefore be viewed as a proposal for a
law-space–aware semantic architecture in which property evaluation is
intrinsically tied to geometric and homotopy-theoretic certificates.



\begin{remark}[Heuristic and Conjectural Relation to Neural Language Models]
\label{rem:haca-heuristic-neural-lm}
From a heuristic and conjectural point of view, one may \emph{interpret}
the empirical success of contemporary large-scale autoregressive
transformer-based language models as suggesting that they implicitly
approximate, in a statistical and non-rigorous manner, certain aspects of
the HACA/PACER architecture developed in this volume. The discussion in
this remark should be read purely as an explanatory picture, and not as a
formal identification of existing architectures with HACA.

On the input side, such models operate on token sequences
$\mathcal{X}^\ast$ and learn internal representations that \emph{appear}
to exhibit multi-scale structure (from token-level patterns to sentence-
 and paragraph-level coherence). It is tempting to view this as echoing,
at least qualitatively, the multi-scale operator universe
\[
  \Omega(L)
  =
  \Omega^{(0)}(L)\sqcup\Omega^{(1)}(L)
  \sqcup\Omega^{(2)}(L)\sqcup\Omega^{(3)}(L)
\]
and the semantic map
$\operatorname{Sem}_L:\Omega(L)\to\mathcal{S}(L)$ in HACA. Likewise,
various downstream classifiers and reward-model heads can be \emph{loosely}
interpreted as playing roles reminiscent of semantic functors and
truth-valued predicates on $\mathcal{S}(L)$: they assign scores or
acceptance probabilities to operator configurations in ways that correlate
with semantic and pragmatic adequacy.

However, in contrast to the explicit HACA construction, all of this
remains implicit and distributional. The models, as currently formulated,
do not expose a typed operator hierarchy, an external KAT structure, or a
semantic functor field with closure and homotopy-stability conditions; nor
do they provide a separate law-space geometry with connections, curvatures,
and Jacobiator-gap certificates. In this sense, and only at a heuristic
level, current neural language models may be viewed as statistical, fuzzy
approximations to an HACA-like architecture: they learn functions that
often behave \emph{as if} there were an underlying multi-scale operator
and semantic law-space, but without making this structure explicit or
enforcing it as a set of axioms.

Under this conjectural reading, the HACA/PACER framework can be regarded
as a law-space–first, certificate-capable completion of patterns that
present-day neural language models seem to capture in a data-first,
statistically learned form. It replaces implicit, distributional
regularities with explicit operators, functors, and geometric constraints,
thereby offering a setting in which minimal-action principles and semantic
admissibility can be stated and, at least in principle, verified in a
mathematically controlled way. This remark is not a claim about the actual
design of existing models, but a speculative explanation of why their
observed behaviour may align with the structures developed here.

\end{remark}

\begin{remark}[Artistic Law-Spaces and Cross-Dialect Expressions]
\label{rem:artistic-law-spaces-cross-dialects}
The LA/LB/LC hierarchy also sheds light on domains where heterogeneity and
cross-dialect behavior are a feature rather than a bug. Consider an
``artistic law-space'' in which different styles, media, and cultural
vocabularies play the role of dialects. A single artwork can be modeled as a
path from an initial intention to a realized configuration that may traverse
several stylistic dialects (for instance, blending classical and deconstructive
elements, textual and visual motifs, or multiple cultural references).

In such a space, Level~A completeness corresponds to the existence of at least
one path from artist intention to audience understanding that is
homotopy-admissible and semantically coherent, even if many other paths are
fragmentary or opaque. Level~B completeness corresponds to the presence of a
homotopy common core: high-quality interpretations, possibly articulated in
different critical dialects (formalism, affective reading, political critique,
etc.), nonetheless converge on a shared structural or thematic nucleus (for
example, a specific tension, contradiction, or motif). Level~C, in artistic
contexts, is usually not realized as a single ``standard dialect'' (which would
risk suppressing creative diversity), but rather as a minimal
structural-symbolic layer: a set of narrative, compositional, or motivic
structures to which all discussable or archivable works can be projected for
purposes of analysis and preservation.

This example illustrates that the notion of a normative subspace in this volume
is flexible: in some domains (for example, regulatory pipelines, technical
documentation) it may resemble a single standard dialect such as Mandarin; in
others (for example, art) it may be better understood as a weak structural core
that supports many cross-dialect trajectories while still enabling meaningful
certificates and comparative evaluation.
\end{remark}



\chapter[Homotopy Networks of Packs and Reading-Path Certificates]{Homotopy Networks of Packs and Reading-Path Certificates:\\
         From Paragraph Graphs to Law-Space Governance}


In previous chapters, we have described how PACER runs can be
represented as paths in a paragraph–monoid network and how law-space
certificates can be attached to these paths. In this chapter, we
focus on the homotopy structure of such networks and on the role of
reading-path certificates in governance: how to use homotopy classes
of reading paths as the basic units for auditing, comparing, and
governing large textual systems such as technical reports, standards,
and regulatory corpora.

\section[Reading-Path Certificates in KAT Homotopy Networks]{Reading-Path Certificates in KAT Homotopy Networks}


Recall from Part~II that a document can be represented as a graph of
paragraph nodes
\[
  v_i = (P_i,M(P_i)),
\]
where $P_i$ is a paragraph and $M(P_i)$ is its paragraph-level KAT
submonoid (Definition~\ref{def:paragraph-submonoids}). Paragraph-level
KAT monoid morphisms
$F_i:M(P_i)\to M(P_{i+1})$ (Definition~\ref{def:paragraph-monoid-morphism})
define directed edges, and PACER reading paths are sequences
\[
  \gamma^{(\ell)}:
  v_{i_0}\xrightarrow{F_{i_0}}v_{i_1}\xrightarrow{F_{i_1}}\cdots
  \xrightarrow{F_{i_{r-1}}}v_{i_r}
\]
in this paragraph–monoid network (Definition~\ref{def:pacer-reading-paths}).
Elementary homotopy moves between such paths give rise to a homotopy
relation $\sim$ on reading paths
(Definition~\ref{def:pacer-reading-path-homotopy}), and the resulting
equivalence classes are referred to as reading-path homotopy classes.

In the law-space–certificate framework, we now refine this notion to
capture the idea that two reading paths may be different at the
syntactic level but \emph{equivalent for governance purposes}.

\begin{definition}[Governance-equivalent reading paths]
\label{def:governance-equivalent-reading-paths}
Let $\gamma^{(1)}$ and $\gamma^{(2)}$ be two PACER reading paths on a
HACA law-object $L$ with the same start and end nodes. We say that
$\gamma^{(1)}$ and $\gamma^{(2)}$ are \emph{governance-equivalent},
written $\gamma^{(1)}\approx_{\mathrm{gov}}\gamma^{(2)}$, if:
\begin{enumerate}
  \item They are homotopy equivalent in the sense of
        Definition~\ref{def:pacer-reading-path-homotopy}, i.e.\
        $\gamma^{(1)}\sim\gamma^{(2)}$ via a finite sequence of
        elementary homotopy moves in the paragraph–monoid network.
  \item Their associated law-space certificate records
        $\mathsf{Cert}(\gamma^{(1)})$ and
        $\mathsf{Cert}(\gamma^{(2)})$ (as constructed by the
        law-space certificate extraction pipeline in
        Definition~\ref{def:lawspace-certificate-extraction}) agree
        up to prescribed tolerances in the following sense:
        \begin{itemize}
          \item the same semantic certificate predicates
                $\mathcal{C}_L$ are satisfied along both paths;
          \item the differences in occupancy measures, strategy weights,
                value geometry vectors, and semantic MDQ profiles stay
                within governance-specified bounds;
          \item the induced Jacobiator-gap certificate data in the
                certificate bundle $P_{\mathrm{cert}}$ lie in the same
                LA/LB/LC completeness regime.
        \end{itemize}
\end{enumerate}
When $\gamma^{(1)}\approx_{\mathrm{gov}}\gamma^{(2)}$, we say that they
belong to the same \emph{governance reading-path class}, denoted
$[\gamma]_{\mathrm{gov}}$.
\end{definition}

\begin{definition}[Reading-path certificate objects]
\label{def:reading-path-certificate-objects}
A \emph{reading-path certificate object} for a document on a HACA
law-object $L$ is a triple
\[
  \bigl([\gamma]_{\mathrm{gov}},\;\mathsf{Cert}(\gamma),\;\kappa\bigr),
\]
where:
\begin{itemize}
  \item $[\gamma]_{\mathrm{gov}}$ is a governance reading-path class
        as in Definition~\ref{def:governance-equivalent-reading-paths};
  \item $\mathsf{Cert}(\gamma)$ is a representative law-space
        certificate record (Definition~\ref{def:lawspace-certificate-extraction})
        for any $\gamma$ in this class;
  \item $\kappa$ is a governance-level metadata object collecting
        policy-relevant annotations, such as the deployment context,
        applicable regulatory regimes, and thresholds used to define
        $\approx_{\mathrm{gov}}$.
\end{itemize}
Reading-path certificate objects form the basic units of analysis for
governance, auditing, and lifecycle management of HACA/PACER systems.
\end{definition}

\begin{remark}[Homotopy-complete books and governance stability]
\label{rem:homotopy-complete-book-gov}
In the homotopy framework of Part~II, a ``homotopy-complete book''
(Definition~\ref{def:homotopy-complete-book}) is one in which
paragraph-level KAT paths can be homotopically deformed within a
controlled network without breaking structural constraints. From a
governance perspective, such a book admits a stable set of reading-
path certificate objects: most admissible reading paths between
designated entry and exit sections fall into a small number of
governance classes $[\gamma]_{\mathrm{gov}}$, and certificate
records within each class differ only within tolerances. This
stability is crucial for predictable behavior: it ensures that
different runs of HACA/PACER on the same document do not produce
governance-relevant surprises, even if they differ in low-level
lexical trajectories.
\end{remark}



\section[Law-Space Governance over Document Collections]{Law-Space Governance over Document Collections}


Let $\{D_j\}_{j\in J}$ be a family of documents (for example, a set
of regulations, standards, or technical reports). We associate to each
$D_j$ a HACA law-object $L_j$ with its own paragraph–monoid network,
certificate bundle $P_{\mathrm{cert}}^{(j)}$, and set of reading-path
certificate objects. To reason about governance at the scale of the
entire collection, we must organize these into a larger law-space
structure.

\begin{definition}[Law-space union for document collections]
\label{def:lawspace-union-doc-collection}
For a family of documents $\{D_j\}_{j\in J}$ with associated HACA
law-objects $L_j$ and law-sectors $B_{L_j}\subseteq\GZLS$, the
\emph{law-space union} of the collection is the subset
\[
  B_{\mathrm{coll}}
  :=
  \bigcup_{j\in J} B_{L_j}
  \subseteq\GZLS,
\]
equipped with the induced restriction of the law-space connection
and curvature. The certificate base space for the collection is the
disjoint union
\[
  B_{\mathrm{cert}}^{\mathrm{coll}}
  :=
  \bigsqcup_{j\in J} B_{\mathrm{cert}}(L_j),
\]
and the certificate bundle is the disjoint union of the individual
bundles
\[
  \pi_{\mathrm{cert}}^{\mathrm{coll}}:
  P_{\mathrm{cert}}^{\mathrm{coll}}
  :=
  \bigsqcup_{j\in J} P_{\mathrm{cert}}^{(j)}
  \longrightarrow
  B_{\mathrm{cert}}^{\mathrm{coll}}.
\]
\end{definition}

In practice, document collections are often interlinked: cross-
references, shared definitions, and consistency requirements induce
additional structure on top of the disjoint union. These can be
encoded as additional edges and 2-cells in the aggregated paragraph–
monoid networks, or as compatibility conditions on certificates
across different $L_j$.

\begin{definition}[Cross-document homotopy and coherence]
\label{def:cross-document-homotopy}
Let $D_j$ and $D_k$ be two documents in the collection, with HACA
law-objects $L_j$ and $L_k$. A \emph{cross-document reading path} is
a sequence of paragraph nodes and morphisms that may traverse both
documents, for example
\[
  v_{i_0}^{(j)}\to\cdots\to v_{i_r}^{(j)}
  \xrightarrow{G_{j\to k}}
  v_{k_0}^{(k)}\to\cdots\to v_{k_s}^{(k)},
\]
where $G_{j\to k}$ represents a cross-reference or logical linkage.
Two cross-document reading paths are said to be \emph{coherently
homotopy-equivalent} if they are related by homotopies that preserve
cross-document constraints (e.g.\ consistency of definitions) and
yield governance-equivalent certificates in the sense of
Definition~\ref{def:governance-equivalent-reading-paths}, now
computed over $B_{\mathrm{coll}}$ and
$P_{\mathrm{cert}}^{\mathrm{coll}}$.
\end{definition}

\begin{remark}[Governance tasks over document collections]
\label{rem:governance-tasks-collections}
At the scale of document collections, typical governance questions
include:
\begin{itemize}
  \item \emph{Coverage and gaps.} Are there parts of the law-space
        union $B_{\mathrm{coll}}$ or of the certificate base space
        $B_{\mathrm{cert}}^{\mathrm{coll}}$ that are never visited by
        any admissible reading paths (under PACER) and thus poorly
        tested or under-specified?
  \item \emph{Coherence and drift.} Do cross-document reading paths
        that implement similar tasks (for example, compliance checks
        across related regulations) belong to the same governance
        classes $[\gamma]_{\mathrm{gov}}$? If not, where do their
        certificates diverge (occupancy patterns, MDQ costs, LA/LB/LC
        completeness levels)?
  \item \emph{Impact of revisions.} When a subset of documents
        $\{D_j\}_{j\in J'}$ is revised, how do the corresponding
        certificate bundles and governance classes change? Can we
        localize the impact to particular segments of the network,
        or do revisions propagate widely through cross-document
        homotopies?
\end{itemize}
These questions can be studied using aggregated reading-path certificate
objects, \\cross-document homotopy relations, and statistics derived from
the collection-level certificate bundle
$\pi_{\mathrm{cert}}^{\mathrm{coll}}$.
\end{remark}



\section[Examples: Regulations, Standards, and Technical Reports]{Examples: Regulations, Standards, and Technical Reports}


In this final section we sketch several abstract scenarios where
homotopy networks of packs and reading-path certificates can inform
governance and system design. The aim is not to propose specific
regulations or institutional arrangements, but to illustrate how the
structures developed in this volume could be instantiated in practice.

\paragraph{Example 1: Financial Regulation Rulebooks.}

Consider a collection of financial regulation rulebooks and guidance
documents, indexed as $\{D_j\}_{j\in J}$. Each $D_j$ is mapped to a
HACA law-object $L_j$ with a paragraph–monoid network, certificate
bundle $P_{\mathrm{cert}}^{(j)}$, and reading-path certificate objects
for typical compliance tasks (e.g.\ anti-manipulation checks, market-
abuse detection, disclosure requirements). PACER reading paths
correspond to ways in which a compliance agent (human or machine)
navigates these rulebooks when answering questions or assessing
scenarios.

Governance questions include:
\begin{itemize}
  \item Do different agents or different system versions follow
        governance-equivalent paths $[\gamma]_{\mathrm{gov}}$ for the
        same task, or do they exploit different parts of the rulebook
        network?
  \item Are there segments of the law-space union $B_{\mathrm{coll}}$
        that are frequently visited by high-stakes tasks but have weak
        LA/LB/LC completeness, suggesting the need for clarification
        or consolidation?
  \item When rules are updated, can reading-path certificate objects
        be used to determine whether historical decisions remain
        defensible under the new regime, or whether they must be
        revisited?
\end{itemize}

\paragraph{Example 2: Pharmaceutical Guidelines and Clinical Protocols.}

A second example is a corpus of pharmaceutical guidelines, clinical
trial protocols, and post-marketing surveillance documents. Each
document is again mapped to a HACA law-object, and typical tasks (e.g.\
protocol design, adverse-event reporting, risk–benefit assessment)
are realized as PACER reading paths with associated law-space
certificates. LBOPB pipelines from Volume~II provide the underlying
life-science law-space, while Volume~III supplies the HACA/PACER
semantic and certificate layers.

Reading-path certificates can help:
\begin{itemize}
  \item document which textual sources and law-space sectors were
        effectively used when designing a protocol or evaluating a
        risk;
  \item identify where different institutions follow diverging
        reading-path governance classes $[\gamma]_{\mathrm{gov}}$ for
        nominally similar tasks, highlighting potential harmonization
        needs;
  \item support audits and investigations by providing structured
        evidence about how guidelines were interpreted in context.
\end{itemize}

\paragraph{Example 3: Technical Standards and Engineering Specifications.}

A third example is a family of technical standards and engineering
specifications (for example, in telecommunications, safety-critical
systems, or interoperability protocols). Here, the paragraph–monoid
networks and reading-path certificates can be used to:
\begin{itemize}
  \item define canonical reading-path governance classes for
        certification and conformity assessment processes;
  \item track how different implementations or vendors traverse the
        same standards: whether they rely on the same structural
        clauses or interpretive paths, or whether their reading-path
        certificates indicate divergent interpretations;
  \item evaluate the impact of standard revisions by comparing
        certificate data for pre- and post-revision reading paths.
\end{itemize}

\begin{remark}[Conceptual blueprint for engineering implementations]
\label{rem:conceptual-blueprint}
Across these examples, the same conceptual blueprint recurs:
\begin{enumerate}
  \item Represent documents as HACA law-objects with paragraph–monoid
        networks.
  \item Instrument PACER-based systems so that each run yields
        reading-path certificate objects
        $([\gamma]_{\mathrm{gov}},\mathsf{Cert}(\gamma),\kappa)$.
  \item Use law-space homotopy, certificate bundles, and LA/LB/LC
        levels to compare and aggregate these objects across tasks,
        versions, and institutions.
\end{enumerate}
This blueprint does not prescribe specific algorithms or institutional
structures; rather, it provides a mathematically structured space of
options within which concrete engineering and governance designs can
be chosen. In this sense, the homotopy networks of packs and reading-
path certificates developed in this volume are proposed as a general
toolkit for certificate-based AI in complex textual and regulatory
environments.
\end{remark}



\chapter[System Architectures, Safety, and Governance in HACA/PACER]{System Architectures, Safety, and Governance in HACA/PACER:\\
         From Engineering Implementations to Policy Interfaces}


Throughout this chapter we work under the law-space hypotheses
introduced earlier: \\HACA/PACER engines are assumed to operate inside
safe law-space sectors in the sense of
Definition~\ref{def:safe-unsafe-sectors}, and the underlying
law-geometry is assumed to satisfy the homotopy completeness levels
LA/LB/LC of Definition~\ref{def:homotopy-completeness-levels}.
Accordingly, the system architectures and governance patterns below
should be read as engineering realizations of the certificate-based
law-space geometry developed in the preceding chapters on
Jacobiator-gap certificates and reading-path certificates, rather than
as purely ad hoc software patterns.

\section[System Architectures for HACA/PACER Engines]{System Architectures for HACA/PACER Engines}


In practical deployments, HACA/PACER engines sit inside larger system
architectures. A typical deployment can be organized into five layers:
\begin{enumerate}
  \item \textbf{Input layer} (corpora and environment observations);
  \item \textbf{Law-space representation layer} (HACA law-objects and
        lexical KAT action monoids);
  \item \textbf{Engine layer} (PACER engines and SAC-PACER-HACA
        learning);
  \item \textbf{Certificate extraction and audit layer};
  \item \textbf{External service and policy interface layer}.
\end{enumerate}
We describe these layers and their interfaces in turn.

\paragraph{Input Layer.}
The input layer ingests textual corpora (documents, standards,
guidelines, regulatory texts) and, when relevant, non-textual
observations from an external environment. Raw textual inputs are
tokenized into sequences in $\Sigma^*$ and embedded into the free
monoid $(\Sigma^*,\circ,\varepsilon)$ and the lexical KAT action
monoid $\mathcal{M}$ associated with a HACA law-object $L$ (see
Definitions~\ref{def:free-monoid-lexical},
\ref{def:lex-kat-action-monoid}). At this layer, the system may reuse
existing infrastructure for data ingestion, preprocessing, and
feature extraction from conventional ML pipelines.

\paragraph{Law-Space Representation Layer.}
The second layer houses the HACA law-space representation:
\begin{itemize}
  \item law-objects $L$ with law-sectors $B_L\subseteq\GZLS$;
  \item multi-scale operator universes $\Omega(L)$ and semantic
        structure spaces $\mathcal{S}(L)$ with maps
        $\operatorname{Sem}_L$ and $\operatorname{Law}_L$;
  \item lexical KAT action monoids $\mathcal{M}$ and paragraph
        submonoids $M(P_i)$.
\end{itemize}
This layer also maintains baseline geometry and safe sectors (as in
Definitions~\ref{def:haca-baseline-law-object},
\ref{def:baseline-deviation}, and
\ref{def:safe-unsafe-sectors}). Interfaces between the input layer
and the law-space representation layer include:
\begin{itemize}
  \item mapping from token sequences to KAT operators and packs;
  \item assignment of inputs to appropriate law-objects $L$ and
        law-sectors $B_L$;
  \item updates to baseline geometry and safe sectors as the
        deployment evolves.
\end{itemize}

\paragraph{Engine Layer: PACER and SAC-PACER-HACA.}
The engine layer runs PACER engines and SAC-PACER-HACA learning
(Definitions~\ref{def:pacer-engine-haca} and
\ref{def:lawspace-flow-sac-pacer-haca}). Its responsibilities include:
\begin{itemize}
  \item executing PACER-based context iteration and reading paths on
        the paragraph–monoid network;
  \item computing lexical and semantic GRL path integrals
        ($S_{\mathrm{lex}}[w]$ and $\mathcal{Z}^{\mathrm{sem}}_\chi[\gamma]$)
        as part of policy evaluation;
  \item adjusting policy parameters $\theta$ via SAC-style updates to
        optimize the law-space objective $\mathcal{J}(\pi_\theta)$
        under the constraints described in
        Section~\ref{def:lawspace-flow-sac-pacer-haca};
  \item maintaining structural coordinates
        $(\mu^\pi,\alpha^\pi,\mathbf{v},\mathbf{q}^{\mathrm{sem}}(\pi))$
        and providing them to downstream modules.
\end{itemize}
This layer can be built on top of existing actor–critic or SAC
implementations by reinterpreting states, actions, and objectives in
HACA law-space, as discussed in the SAC-PACER-HACA chapter.

\paragraph{Certificate Extraction and Audit Layer.}
The fourth layer implements the certificate extraction pipeline
(Definition~\ref{def:lawspace-certificate-extraction}) and the
certificate bundles (Definitions~\ref{def:certificate-base-space},
\ref{def:certificate-bundle}). Its functions include:
\begin{itemize}
  \item logging PACER runs at the level of KAT operators, paragraph
        nodes, and law-space paths;
  \item reconstructing reading paths and computing reading-path
        certificates $\mathsf{Cert}(\gamma)$;
  \item lifting certificates to the certificate bundle
        $P_{\mathrm{cert}}$ and tracking LA/LB/LC completeness;
  \item aggregating certificate data into records suitable for
        version-to-version comparison, audit evidence, and rollback
        decisions (Definition~\ref{def:certificate-based-evaluation-tasks}).
\end{itemize}
Interfaces to the engine layer include log streams, policy identifiers,
and structural coordinate snapshots; interfaces to the service and
policy layer include certificate queries, dashboards, and alerts.

\paragraph{External Service and Policy Interface Layer.}
The top layer exposes HACA/PACER capabilities to external users and
governance bodies. It includes:
\begin{itemize}
  \item application-facing APIs (e.g.\ question answering, drafting,
        summarization, analysis) that internally route requests to
        PACER engines and law-space modules;
  \item safety and guardrail services that consult law-space
        certificates and safe sectors before returning outputs
        (filtering, red-teaming, or gating on certificate checks);
  \item policy interfaces through which human governance structures
        can configure certificate families $\mathcal{C}_L$, thresholds
        on MDQ and value geometry, and rules for gating, rollback, and
        audit.
\end{itemize}
While the underlying numerical implementation may resemble existing
ML/RL service stacks, the presence of an explicit law-space
representation layer and a certificate bundle layer is the key
architectural difference introduced by this volume.

\begin{remark}[Interfaces with existing ML/RL stacks]
\label{rem:interfaces-ml-rl-stacks}
In many settings, HACA/PACER engines need not replace existing ML/RL
components outright. Instead, they can be integrated as additional
layers that interpret existing neural models as approximate realizations
of parts of the HACA structure (see
Section~\ref{subsec:haca-vs-neural-lm} and
Section~\ref{sec:haca-positioning-lm}). For example:
\begin{itemize}
  \item a pretrained transformer may serve as a feature extractor for
        semantic functors or as a source of candidate packs and paths;
  \item existing reward models can be reinterpreted as noisy proxies
        for semantic GRL path integrals or MDQ-based cost terms;
  \item existing logging and monitoring infrastructure can be extended
        to capture the additional law-space and certificate data
        required by the HACA/PACER framework.
\end{itemize}
In this way, the system architecture can evolve gradually from a purely
data-first stack toward a law-space–aware, certificate-based design.
\end{remark}



\section[Safety and Governance Mechanisms for HACA/PACER Systems]{Safety and Governance Mechanisms for HACA/PACER Systems}


HACA/PACER systems incorporate multiple safety and governance mechanisms
that operate across the system lifecycle. At a high level, these
mechanisms can be organized along two axes:
\begin{itemize}
  \item by \emph{mechanism type}: safe sectors, Jacobiator-gap
        certificates, MDQ bounds, certificate bundles, output
        filtering, and so on;
  \item by \emph{lifecycle phase}: design and specification, training
        and validation, deployment, monitoring, and retirement.
\end{itemize}

\paragraph{Mechanism Types.}
The core mechanisms include:
\begin{enumerate}
  \item \textbf{Safe sectors and baseline geometry.} Law-space safe
        sectors $S\subseteq B_L$ and baseline geometries
        $(\GZLC^{\mathrm{base}},\GZLCurv^{\mathrm{base}})$ bound the
        regions in which policies are allowed to operate
        (Definitions~\ref{def:haca-baseline-law-object},
         \ref{def:safe-unsafe-sectors}, and
         \ref{def:baseline-deviation}).
  \item \textbf{Certificate bundles and Jacobiator-gap checks.}
        Certificate bundles $\pi_{\mathrm{cert}}:P_{\mathrm{cert}}\to
        B_{\mathrm{cert}}(L)$ and Jacobiator-gap invariants provide
        higher-order homotopy checks for structural consistency
        (Definitions~\ref{def:haca-jacobiator-gap-invariant},
         \ref{def:certificate-bundle}).
  \item \textbf{MDQ and occupancy bounds.} Semantic MDQ profiles
        $\mathbf{q}^{\mathrm{sem}}(\pi)$ and occupancy measures
        $\mu^\pi$ are constrained to lie inside admissible regions
        (Definitions~\ref{def:semantic-mdq-profile},
         \ref{def:baseline-constrained-mdq-region}). Exceeding MDQ
        budgets or occupancy thresholds triggers alarms or enforcement
        actions.
  \item \textbf{Output screening and guardrails.} Outputs are screened
        using semantic certificates $\mathcal{C}_L$ and reading-path
        certificates: outputs whose generation involves paths leaving
        safe sectors or violating key certificates are blocked,
        flagged, or routed for human review.
  \item \textbf{Logging and audit trails.} Detailed logs of PACER
        operations (KAT operators, paragraph paths, law-space paths,
        certificate evaluations) are maintained and used to construct
        law-space certificate records for later audit and analysis
        (Definition~\ref{def:lawspace-certificate-extraction}).
\end{enumerate}

\paragraph{Lifecycle Phases.}
Safety and governance mechanisms must be integrated across the
lifecycle of HACA/PACER systems:

\subparagraph{Design and Specification.}
In the design phase, domain experts and governance bodies:
\begin{itemize}
  \item specify the relevant HACA law-objects $L$ and law-space
        sectors $B_L$ for the intended application domain;
  \item select semantic certificate predicates $\mathcal{C}_L$ and
        define safe sectors, normative subspaces, and LA/LB/LC
        completeness criteria appropriate for the domain;
  \item identify critical tasks and reading-path classes that require
        especially strong certificates (e.g.\ high-stakes decisions).
\end{itemize}

\subparagraph{Training and Validation.}
During training and offline validation:
\begin{itemize}
  \item SAC-PACER-HACA engines are trained to minimize $\mathcal{J}(\pi)$
        under constraints, with penalties for leaving safe sectors or
        violating semantic certificates;
  \item certificate data (MDQ budgets, LA/LB/LC levels, certificate
        coverage of tasks) are used as additional validation metrics,
        alongside standard predictive accuracy or reward-based metrics;
  \item homotopy completeness and certificate coverage analyses are
        used to identify under-tested regions of law-space or weak
        parts of the paragraph–monoid network.
\end{itemize}

\subparagraph{Deployment and Monitoring.}
In deployment:
\begin{itemize}
  \item runtime monitoring continuously tracks reading-path
        certificates, MDQ usage, and visits to safe/unsafe sectors;
  \item certificate-based alerts are generated when trajectories
        approach boundaries of safe sectors, exceed MDQ budgets, or
        deviate from established governance reading-path classes
        $[\gamma]_{\mathrm{gov}}$;
  \item aggregated certificate statistics are used to detect concept
        drift, gradual degradation of LA/LB/LC completeness, or
        emerging patterns of unsafe behavior.
\end{itemize}

\subparagraph{Rollback, Update, and Retirement.}
Finally, in the update and retirement phases:
\begin{itemize}
  \item certificate comparisons between successive policy versions
        (Definition~\ref{def:certificate-based-evaluation-tasks}) guide
        decisions about promotions, rollbacks, or staged deployments;
  \item reading-path certificates are used to determine whether past
        outputs remain acceptable under new policies and law-space
        configurations, or whether they need re-evaluation;
  \item retirement policies can require that certificate records and
        audit trails be preserved for a specified horizon to support
        future investigations or regulatory reviews.
\end{itemize}

\begin{remark}[Alignment with existing AI safety and compliance frameworks]
\label{rem:alignment-ai-safety-frameworks}
The mechanisms described above are compatible with many existing AI
safety and compliance frameworks, which typically emphasize:
\begin{itemize}
  \item risk identification and classification;
  \item evaluation and testing procedures;
  \item monitoring, logging, and incident response;
  \item documentation and transparency requirements.
\end{itemize}
HACA/PACER does not attempt to replace these frameworks, but to
provide a mathematically structured substrate on which they can be
instantiated. Law-space safe sectors, certificate bundles, MDQ budgets,
and reading-path certificates offer concrete, machine-checkable
counterparts to abstract notions of ``risk boundaries'', ``test
coverage'', and ``auditability''. This allows governance bodies to
formulate policies in conceptual terms while still enabling technical
implementations that can be analyzed and verified using the tools
developed in this volume.
\end{remark}



\section[Positioning HACA/PACER vs Neural LM Paradigms]{Positioning HACA/PACER Relative to Contemporary \\Language Modeling Paradigms}
\label{sec:haca-positioning-lm}



A detailed structural comparison between the semantic law-space of
HACA/PACER and corpus-based semantics in neural language models was
given in Section~\ref{subsec:haca-vs-neural-lm}. In this concluding
section we revisit that comparison from an architectural and
governance-oriented perspective, emphasizing how HACA/PACER systems
can coexist with, and potentially structure, contemporary language
modeling paradigms.

A recurring theme throughout this volume has been the distinction between
law-space–first constructions and data-first constructions of semantics.
The HACA/PACER framework adopts a law-space–first perspective: it starts
from an explicit law-space $\GZLS$ with connections and curvatures, then
builds multi-scale operator universes, semantic structure spaces, and
truth-valued functors in such a way that structural properties are
realized as geometric sectors of law-space equipped with homotopy-based
certificates.

In contrast, contemporary large-scale autoregressive transformer-based
language models~\cite{Vaswani2017Attention} instantiate a data-first paradigm. Their semantics is
primarily distributional: internal representations are shaped by the
statistics of large text corpora, and any emergent notion of a ``rule''
is encoded implicitly in the parameterization of conditional probability
distributions. There is no separate law-space geometry in which laws can
be articulated, deformed, and certified; instead, both ``laws'' and
``properties'' are absorbed into the same learned mapping from token
sequences to output distributions.

From the viewpoint developed here, HACA/PACER can be interpreted as a
proposal to disentangle these roles. Operator universes and semantic
structures are treated as \emph{law-aware} objects: their admissible
transformations are constrained by explicit semantic operations and their
properties are required to be observable, homotopy-stable, and
geometrically realizable in $\GZLS$. This makes it possible, at least in
principle, to attach certificates---in the form of Jacobiator-gap
invariants and LA/LB/LC levels---to semantic properties of reading
paths, compressed expansions, and hierarchical cognitive architectures.

By comparison, current neural language modeling paradigms can be seen as
providing highly effective approximations to certain projections of this
structure: they implement powerful statistical maps from operator-level
inputs (token sequences) to internal representations that correlate with
semantic and pragmatic features, and they realize various property tests
via discriminative heads. What is missing is an explicit law-space layer
and a systematic correspondence between semantic properties and
geometry-level laws.

This gap suggests a natural direction for further work. One may regard
HACA/PACER as a law-space template, and ask to what extent existing
neural architectures can be enriched with law-space structure: for
example, by constraining internal representations to live over explicit
operator bundles, by training truth-valued functors that satisfy
closure and homotopy-stability conditions, or by coupling the training
objective to geometric invariants in $\GZLS$. Such developments would not
replace corpus-based training, but would reorganize it around an
explicit notion of law-space, bringing neural language modeling closer
to the certificate-based AI paradigm developed in this volume.





\appendix

\chapter{Abstract Example Schemes for HACA/PACER Operator Paths and Context Networks}


In this appendix we collect several abstract example schemes that
illustrate how the structures developed in this volume---HACA law-
objects, lexical KAT action monoids, paragraph--monoid networks,
PACER reading paths, GRL path integrals, and certificate bundles---can
be instantiated in typical tasks. The aim is not to specify concrete
algorithms or datasets, but to provide schematic patterns of operator
paths and context networks that may guide engineering implementations.
The examples are intentionally minimal and use only the notation and
structures already defined in the main text.

\section{Example Scheme I: Single-Document Question Answering}


\paragraph{Setup.}
Let $D$ be a technical document (for example, a specification or
research report) segmented into paragraphs
\[
  P_1,\dots,P_N.
\]
We associate to $D$ a HACA law-object $L$ with law-sector
$B_L\subseteq\GZLS$, multi-scale operator universe $\Omega(L)$,
semantic structure space $\mathcal{S}(L)$, and semantic--law map
$\Phi_L:\Omega(L)\to B_L$ as in
Definitions~\ref{def:haca-law-objects-category}
and~\ref{def:haca-sem-law-projections}. The lexical KAT action monoid
$\mathcal{M}$ and paragraph-level submonoids $M(P_i)$ are constructed
as in Part~II, yielding paragraph nodes
\[
  v_i = (P_i,M(P_i)).
\]

A user query $q$ (e.g.\ ``Does this document permit X under
condition Y?'') is presented at the external service layer and passed
to the HACA/PACER engine. We regard $q$ as specifying:
\begin{itemize}
  \item an initial set of paragraph nodes $\{v_{i_0}\}$ deemed
        relevant as entry points (for instance, sections indexed as
        ``scope'', ``definitions'', or ``conditions''); and
  \item a semantic certificate subset
        $\mathcal{C}_L^{(q)}\subseteq\mathcal{C}_L$ capturing the
        relevant semantic constraints (e.g.\ internal consistency with
        definitions, correct use of legal terms, safety conditions).
\end{itemize}

\paragraph{Operator Paths and Context Iteration.}
The PACER engine runs context iteration starting from the initial
nodes. At each step $t$:
\begin{itemize}
  \item a context window (e.g.\ a few paragraphs around $P_{i_t}$) is
        selected;
  \item a finite set of lexical operators
        $\mathcal{G}_t=\{G_{t,k}\}\subseteq\mathcal{M}$ is chosen;
  \item the local context submonoid $M_t=\langle\mathcal{G}_t\rangle$
        is formed (Definition~\ref{def:local-context-submonoid});
  \item lexical operations $G_{t,k}$ are applied to construct or
        update candidate answer segments (e.g.\ extracting relevant
        clauses, normalizing or rephrasing text).
\end{itemize}
The entire run yields an operator path $w\in\mathsf{Paths}(\mathcal{M})$
and a paragraph reading path
\[
  \gamma:
  v_{i_0}\xrightarrow{F_{i_0}}v_{i_1}\xrightarrow{F_{i_1}}\cdots
  \xrightarrow{F_{i_{r-1}}}v_{i_r},
\]
where each $F_{i_j}:M(P_{i_j})\to M(P_{i_{j+1}})$ encodes how KAT
operations transport context from $P_{i_j}$ to $P_{i_{j+1}}$.

\paragraph{Semantic GRL Evaluation and Answer Construction.}
For the resulting reading path $\gamma$ and associated operator
configurations $\omega_t\in\Omega(L)$, the engine:
\begin{itemize}
  \item computes lexical GRL actions $S_{\mathrm{lex}}[w]$ along the
        operator path;
  \item computes semantic GRL path integrals
        $\mathcal{Z}^{\mathrm{sem}}_\chi[\gamma]$ and predicates
        $\Xi_\chi(\gamma)$ for $\chi\in\mathcal{C}_L^{(q)}$;
  \item checks whether $\Xi_\chi(\gamma)=\mathrm{true}$ for all
        certificates in the query-specific set
        $\mathcal{C}_L^{(q)}$.
\end{itemize}
If the certificate checks pass and the combined objective
$\mathcal{J}(\pi)$ (from the SAC-PACER-HACA chapter) falls below a
governance-specified threshold, the engine constructs an answer by
assembling relevant packs from the visited paragraphs (e.g.\ clauses
and definitions), possibly compressing and rephrasing them while
preserving semantic admissibility.

\paragraph{Context Network and Certificate Record.}
At the certificate extraction layer, the run is summarized as a
reading-path certificate record $\mathsf{Cert}(\pi)$ containing:
\begin{itemize}
  \item the paragraph-level path $\gamma$ and the sequence of nodes
        visited;
  \item the semantic certificate outcomes
        $\Xi_\chi(\gamma)$ for $\chi\in\mathcal{C}_L^{(q)}$;
  \item MDQ usage and occupancy measures $\mu_i^\pi$ along the path;
  \item Jacobiator-gap certificate data for the relevant law-objects
        and property layers.
\end{itemize}
Two different runs answering the same query $q$ are considered
governance-equivalent if they fall into the same governance
reading-path class $[\gamma]_{\mathrm{gov}}$ and yield certificate
records that agree within tolerances
(Definition~\ref{def:governance-equivalent-reading-paths}).

\section{Example Scheme II: Summarization with Compliance Constraints}


\paragraph{Setup.}
Let $D_{\mathrm{src}}$ be a long source document (e.g.\ a technical
report or internal policy document), and let the goal be to produce a
shorter summary $D_{\mathrm{sum}}$ that is faithful, informative, and
compliant with an external policy baseline. We associate:
\begin{itemize}
  \item a HACA law-object $L_{\mathrm{src}}$ for $D_{\mathrm{src}}$;
  \item a HACA law-object $L_{\mathrm{sum}}$ for the space of possible
        summaries, possibly with a different baseline geometry and
        normative subspace (for example, a stricter style guide);
  \item a family of certificate predicates
        $\mathcal{C}_{\mathrm{content}}$ for content faithfulness and
        $\mathcal{C}_{\mathrm{policy}}$ for policy compliance.
\end{itemize}
The task is realized as a PACER-driven mapping from reading paths in
$L_{\mathrm{src}}$ to operator paths in $L_{\mathrm{sum}}$, under
GRL-based compression and certificate constraints.

\paragraph{Operator Paths: Compression and Selection.}
On $L_{\mathrm{src}}$, PACER constructs reading paths
$\gamma_{\mathrm{src}}$ that identify segments relevant to the summary
objective (e.g.\ key results, main arguments). On $L_{\mathrm{sum}}$,
PACER operates on a lexical KAT monoid
$\mathcal{M}_{\mathrm{sum}}$ to construct a candidate summary path
$w_{\mathrm{sum}}$:
\begin{itemize}
  \item compression operators remove redundancy and low-salience
        packs;
  \item alignment operators arrange segments into a coherent
        narrative order;
  \item expansion operators introduce minimal connective sentences
        to maintain readability and logical flow.
\end{itemize}

\paragraph{Semantic and Policy Certificates.}
Semantic GRL path integrals are used at two levels:
\begin{itemize}
  \item for faithfulness, path integrals
        $\mathcal{Z}^{\mathrm{sem}}_\chi[\gamma_{\mathrm{src}},
        w_{\mathrm{sum}}]$ with $\chi\in\mathcal{C}_{\mathrm{content}}$
        evaluate whether the summary preserves key semantic content
        from the source reading path;
  \item for compliance, semantic predicates in
        $\mathcal{C}_{\mathrm{policy}}$ test whether the summary
        respects policy constraints (e.g.\ required disclaimers,
        forbidden claims, style restrictions).
\end{itemize}
The joint objective combines a compression reward (shorter summaries),
``faithfulness penalties'', and compliance penalties, together with MDQ
and safe-sector constraints.

\paragraph{Context Network and Governance Classes.}
The resulting context network involves nodes from both $L_{\mathrm{src}}$
and $L_{\mathrm{sum}}$, with edges representing:
\begin{itemize}
  \item reading-path segments in the source document;
  \item operator paths in the summary space;
  \item cross-law-object links encoding ``summarize-from'' relations.
\end{itemize}
Reading-path certificate objects record how each summary run navigated
this network and which certificates were satisfied. Governance
reading-path classes can then be defined over pairs
$(\gamma_{\mathrm{src}},w_{\mathrm{sum}})$, supporting tasks such as:
\begin{itemize}
  \item comparing different summary algorithms or system versions;
  \item verifying that all mandated policy constraints were enforced;
  \item tracing back, for a given problematic summary, which parts of
        the source document and which operators contributed to it.
\end{itemize}

\section{Example Scheme III: Cross-Document Consistency Checking}


\paragraph{Setup.}
Let $\{D_j\}_{j\in J}$ be a small collection of related documents
(e.g.\ a base standard, a set of amendments, and an implementation
guideline). Each $D_j$ is mapped to a HACA law-object $L_j$ with
law-sector $B_{L_j}\subseteq\GZLS$ and paragraph–monoid network.
Cross-references and logical links between documents are encoded as
cross-document morphisms in the sense of
Definition~\ref{def:cross-document-homotopy}.

The goal is to check whether, for a given family of tasks (e.g.\
compliance checks or technical requirements), the cross-document
reading paths implemented by HACA/PACER remain coherent: that is, they
do not lead to inconsistent or contradictory interpretations.

\paragraph{Operator Paths and Cross-Document Reading Paths.}
For a specific task type $T$, the PACER engine generates cross-document
reading paths of the form
\[
  \gamma_T^{(\ell)}:
  v_{i_0}^{(j_0)}\to\cdots\to v_{i_r}^{(j_r)}\to
  v_{k_0}^{(k_0)}\to\cdots\to v_{k_s}^{(k_s)},
\]
where superscripts index documents and arrows can represent both
intra-document and cross-document transitions. Corresponding operator
paths $w_T^{(\ell)}$ in the union of KAT monoids
$\bigsqcup_j\mathcal{M}^{(j)}$ are constructed, and structural
coordinates $(\mu^{(\ell)},\alpha^{(\ell)},\mathbf{v}^{(\ell)})$ are
recorded for each run~$\ell$.

\paragraph{Consistency Certificates.}
Semantic predicates in $\mathcal{C}_{\mathrm{consistency}}$ are
chosen to represent cross-document consistency requirements (e.g.\
``definition of term X is consistent across documents'' or ``no direct
contradiction between clauses about Y''). For each reading path
$\gamma_T^{(\ell)}$, the engine computes:
\begin{itemize}
  \item semantic GRL path integrals over the cross-document path;
  \item semantic predicates $\Xi_\chi(\gamma_T^{(\ell)})$ for
        $\chi\in\mathcal{C}_{\mathrm{consistency}}$;
  \item Jacobiator-gap certificate data for property-layer relations
        between the law-objects $L_j$ involved in the path.
\end{itemize}
A path that fails any of these checks indicates a potential source of
inconsistency or ambiguity, which can then be inspected by human
experts.

\paragraph{Homotopy Classes and Governance Interpretation.}
Cross-document reading paths for the same task type $T$ are grouped
into governance classes $[\gamma_T]_{\mathrm{gov}}$ as in
Definition~\ref{def:governance-equivalent-reading-paths}, now
interpreted over the collection-level law-space and certificate bundle
(Definition~\ref{def:lawspace-union-doc-collection}). Governance
analyses can then address questions such as:
\begin{itemize}
  \item whether all institutions or systems follow governance-equivalent
        paths for the same task;
  \item whether certain cross-document paths systematically fail
        consistency certificates;
  \item how revisions to one document change the homotopy classes of
        cross-document reading paths.
\end{itemize}
These abstract schemes are intended as templates: concrete deployments
can refine them with domain-specific choices of law-objects,
certificates, and task types while reusing the same underlying HACA/
PACER operator and context network structures.



\chapter{Proofs for KAT-Monoid Homotopy Networks and Certificate Conditions}


This appendix collects proofs and technical details for results used
in the main text concerning KAT-monoid homotopy networks and certificate
conditions. We focus on four themes:
\begin{itemize}
  \item basic properties of lexical KAT action monoids and power
        submonoid spectra;
  \item homotopy relations on paragraph-level reading paths and their
        algebraic properties;
  \item existence and structure of operator and certificate bundles
        under homotopy admissibility and homotopy completeness
        assumptions;
  \item stability and equivalence properties of governance-oriented
        reading-path classes.
\end{itemize}
Throughout, we use only the notation and definitions introduced in the
main part of the volume.

\section[Semantic Functor Algebras]{Semantic Functor Algebras as Degenerate PFB-GNLA \\Representations}
\label{app:sec-semantic-functor-proof}

\begin{proof}[Proof of Proposition~\ref{prop:semantic-functor-as-degenerate-rep}]
By Definition~\ref{def:semantic-degeneration-pfb-gnla}, the semantic
sector of the PFB-GNLA for a HACA law-object $L$ is encoded by a
principal bundle $P^{\mathrm{sem}}_L \to B_L$ equipped with a
generalized noncommutative Lie algebra
$\GAlgebra^{\mathrm{sem}}_L$ acting on horizontal lifts, together with
an associated Jacobiator-gap. The semantic degeneration identifies
those directions in the full PFB-GNLA that survive in the semantic
sector and quotients out the remaining directions by structural
constraints.

Definition~\ref{def:noncommutative-semantic-functor-algebra} constructs
$\mathfrak{F}_L$ as the algebra of noncommutative semantic functors
obtained by composing evaluation along the law-connection with elements
of $\GAlgebra^{\mathrm{sem}}_L$. In particular, the semantic bracket
$[\cdot,\cdot]_{\mathrm{sem}}$ and the Jacobiator-gap
$J(\cdot,\cdot,\cdot)$ in $\mathfrak{F}_L$ are defined by pushing
forward the generalized bracket and Jacobiator-gap in
$\GAlgebra^{\mathrm{sem}}_L$ along this evaluation map. This yields a
representation homomorphism
\[
  \rho_L \colon \GAlgebra^{\mathrm{sem}}_L
  \longrightarrow
  \mathrm{Der}(\mathfrak{F}_L),
\]
which intertwines brackets and Jacobiator-gaps as stated in
items~(1)--(3) of
Proposition~\ref{prop:semantic-functor-as-degenerate-rep}: the bracket
on $\mathfrak{F}_L$ and its Jacobiator-gap are the images of the
corresponding structures in $\GAlgebra^{\mathrm{sem}}_L$ under
$\rho_L$.

The representation is degenerate because the kernel of $\rho_L$
consists exactly of those directions in $\GAlgebra^{\mathrm{sem}}_L$
that act trivially on all semantic functors; these directions are
precisely those killed by the structural constraints defining the HACA/
PACER law-object $L$. Thus $\mathfrak{F}_L$ is a degenerate
structural-constraint realization of the underlying PFB-GNLA: it
retains the noncommutative and Jacobiator-gap features of the full
theory, but only along those directions singled out by the semantic
sector and subject to the constraints defining $L$. This is exactly the
packaged content of
Proposition~\ref{prop:semantic-functor-as-degenerate-rep}. The general
representation framework for PFB-GNLA is developed in detail in the
pure mathematics volume~\cite{GaoZheng2025GFramework}.
\end{proof}

\section[Curvature Control of Semantic MDQ]{Curvature Control of Semantic MDQ}
\label{app:sec-logical-pressure-mdq-proof}

\begin{proof}[Proof of Proposition~\ref{prop:logical-pressure-mdq-control}]
By the definition of the semantic MDQ density, for each admissible
operator lift $\omega_t$ there exist constants $A_1, A_2 \ge 0$ and a
finite family of semantic functors $\{F_\alpha\}$ such that
\[
  \bigl|q_{\mathrm{MDQ}}^{\mathrm{sem}}(t)\bigr|
  \;\le\;
  A_1 + A_2
  \sup_\alpha
  \bigl\|
    \nabla_{\dot{\gamma}(t)} F_\alpha
  \bigr\|
\]
for almost every $t\in[0,1]$. This is a direct consequence of the
construction of $q_{\mathrm{MDQ}}^{\mathrm{sem}}$ as a law-space
functional applied to the family $\{F_\alpha\}$.

Under the stated regularity and boundedness assumptions, the covariant
derivatives of the semantic functors along $\gamma$ satisfy a standard
curvature-control estimate for the law-space connection:
\[
  \bigl\|
    \nabla_{\dot{\gamma}(t)} F_\alpha
  \bigr\|
  \;\le\;
  a_1\,\bigl\|F_\alpha(\gamma(t))\bigr\|
  +
  a_2\,\bigl\|\mathcal{F}_L(\gamma(t))\bigr\|,
\]
where $\mathcal{F}_L$ denotes the semantic curvature two-form
associated with the law-space connection (see the pure mathematics
volume~\cite{GaoZheng2025GFramework} and Applications Volume~I for the
underlying MDQ estimates). The constants $a_1,a_2$ depend only on the
regularity bounds for the connection and the semantic functors.

By construction of the logical pressure field
(Definition~\ref{def:logical-pressure-field}) there exist constants
$b_1, b_2 \ge 0$ such that
\[
  \bigl\|\mathcal{F}_L(\gamma(t))\bigr\|
  \;\le\;
  b_1 + b_2\,\mathcal{F}_{\mathrm{log}}(\gamma(t))
\]
for almost every $t$. Combining the three inequalities and absorbing
all fixed constants into new constants $C_1, C_2 \ge 0$, we obtain
\[
  \bigl|q_{\mathrm{MDQ}}^{\mathrm{sem}}(t)\bigr|
  \;\le\;
  C_1 + C_2\,\mathcal{F}_{\mathrm{log}}(\gamma(t))
\]
for almost every $t \in [0,1]$, which is exactly the estimate claimed
in Proposition~\ref{prop:logical-pressure-mdq-control}.
\end{proof}

\section[Operator Bundles from Property Intersections]{Operator Bundles from Property Intersections and Homotopy Completeness}
\label{app:sec-operator-bundles}

\begin{lemma}[Trivializing the Jacobiator-gap obstruction]
\label{app:lem-jacobiator-gap-trivialization}
Let $U \subset B_L$ be a law-space region for a HACA law-object $L$,
equipped with a finite good open cover $\{U_i\}_{i\in I}$ by property
intersections in the sense of the main text. On each $U_i$ we have a
family of property-changing operators with structure group
$G_{\mathrm{op}}$, and on double overlaps $U_{ij} = U_i \cap U_j$ we
have transition maps
\[
  h_{ij} \colon U_{ij} \longrightarrow G_{\mathrm{op}}
\]
relating admissible operator lifts. On triple overlaps
$U_{ijk} = U_i \cap U_j \cap U_k$, the failure of the naive cocycle
condition $h_{ij}h_{jk} = h_{ik}$ is measured by a Jacobiator-gap
valued 2-cocycle
\[
  J_{ijk} \in Z\bigl(G_{\mathrm{op}}\bigr),
  \qquad
  h_{ij}h_{jk} = J_{ijk}\,h_{ik}.
\]
Assume that the higher homotopy harmonization conditions and the
homotopy completeness level LA/LB/LC hold on~$U$ for the
property-changing operators. Then the cohomology class of
$\{J_{ijk}\}$ in the Čech cohomology group
$\check{H}^2\bigl(\{U_i\},Z(G_{\mathrm{op}})\bigr)$ vanishes. In
particular, there exist $G_{\mathrm{op}}$-valued 1-cochains
$c_i \colon U_i \to G_{\mathrm{op}}$ such that the modified transition
maps
\[
  g_{ij} := c_i^{-1} h_{ij} c_j
\]
satisfy the strict cocycle condition
$g_{ij}g_{jk} = g_{ik}$ on every triple overlap $U_{ijk}$.
\end{lemma}

\begin{proof}
The higher homotopy harmonization conditions and the LA/LB/LC
homotopy completeness assumptions precisely state that the family of
heterogeneity morphisms and property-changing operators on~$U$
extends, up to controlled homotopy, to a coherent system on the nerve
of the cover $\{U_i\}_{i\in I}$. Concretely, the Jacobiator-gap data
$\{J_{ijk}\}$ encode a 2-cocycle in the Čech complex with coefficients
in the central subgroup $Z(G_{\mathrm{op}})$, and the $(n)$-th order
homotopy harmonization provides higher homotopies that contract this
2-cocycle in the appropriate homotopy category.

Since the cover is good, the nerve is homotopy equivalent to~$U$, and
the LA/LB/LC completeness ensures that the relevant homotopy groups
governing the obstruction class vanish in the dimensions where the
Jacobiator-gap lives. As a result, the cohomology class of
$\{J_{ijk}\}$ in $\check{H}^2(\{U_i\},Z(G_{\mathrm{op}}))$ is zero,
i.e., $\{J_{ijk}\}$ is a coboundary. Hence there exist 1-cochains
$c_i \colon U_i \to Z(G_{\mathrm{op}}) \subset G_{\mathrm{op}}$ such
that
\[
  J_{ijk}
  \;=\;
  c_i^{-1} c_j c_k^{-1}
  \quad\text{on }U_{ijk}.
\]
Defining $g_{ij} := c_i^{-1} h_{ij} c_j$ on each double overlap, a
direct computation shows that
\[
  g_{ij} g_{jk}
  \;=\;
  c_i^{-1} h_{ij} c_j c_j^{-1} h_{jk} c_k
  \;=\;
  c_i^{-1} h_{ij} h_{jk} c_k
  \;=\;
  c_i^{-1} J_{ijk} h_{ik} c_k
  \;=\;
  c_i^{-1} h_{ik} c_k
  \;=\;
  g_{ik},
\]
where in the penultimate equality we have used the explicit expression
of $J_{ijk}$ as a coboundary. This proves the existence of modified
transition maps $\{g_{ij}\}$ satisfying the strict cocycle condition.
\end{proof}

\begin{proposition}[Operator bundles from property intersections]
\label{app:prop-operator-bundle-from-property-intersections}
In the setting of Lemma~\ref{app:lem-jacobiator-gap-trivialization},
there exists a principal $G_{\mathrm{op}}$-bundle
\[
  \pi_{\mathrm{op}} \colon P_{\mathrm{op}} \longrightarrow U
\]
with local trivialisations
\[
  \Psi_i \colon U_i \times G_{\mathrm{op}} \longrightarrow
  \pi_{\mathrm{op}}^{-1}(U_i)
\]
whose transition functions on double overlaps $U_{ij}$ are exactly the
$g_{ij}$ constructed in Lemma~\ref{app:lem-jacobiator-gap-trivialization}.
The fibre over each point $x \in U$ can be identified with the set of
equivalence classes of admissible property-changing operators along
heterogeneity morphism chains passing through~$x$, modulo the natural
right action of $G_{\mathrm{op}}$. Moreover:
\begin{enumerate}
  \item On regions where the LA-level homotopy completeness holds, the
        bundle $P_{\mathrm{op}}$ admits local gauges adapted to the
        basic property intersections, corresponding to the existence of
        local trivializations with small homotopy diameter.
  \item On regions where LB-level completeness holds, these local
        gauges can be chosen coherently along prescribed path families,
        reflecting the ability to parameterize operator lifts along
        one-dimensional law-space trajectories.
  \item On regions where LC-level completeness holds, the gauges can be
        extended over higher-dimensional homotopy data, corresponding
        to the full bundle-level control discussed in the main text.
\end{enumerate}
\end{proposition}

\begin{proof}
Given the strict $G_{\mathrm{op}}$-valued 1-cocycle
$\{g_{ij}\}$ on the good cover $\{U_i\}$, the standard principal
bundle gluing construction (see, e.g., the principal bundle framework
developed in the pure mathematics volume~\cite{GaoZheng2025GFramework})
produces a principal $G_{\mathrm{op}}$-bundle 
$\pi_{\mathrm{op}} \colon P_{\mathrm{op}} \to U$ with the desired
local trivialisations $\Psi_i$ and transition functions $g_{ij}$. By
construction, two operator lifts on an overlap differ by right
multiplication by an element of $G_{\mathrm{op}}$, and the quotient
by this right action yields the fibre description stated above.

The three items on LA/LB/LC follow by unpacking the homotopy
completeness conditions in the concrete bundle setting. LA-level
completeness ensures that the nerve of the cover can be refined so
that each connected component admits a gauge in which the transition
functions are homotopically trivial, yielding local trivializations
with small homotopy diameter. LB-level completeness provides control
along 1-dimensional homotopies, which allows the local gauges to be
chosen coherently along path families, implementing a bundle-level
version of pathwise consistency. LC-level completeness, finally,
ensures that the obstruction to extending such gauges over higher
dimensional homotopies vanishes, so that the local gauges can be
globalized over the relevant homotopy data. This matches precisely the
interpretation of LA/LB/LC given in the main text and yields the
bundle-level structure used in
Proposition~\ref{prop:gaugeable-operator-bundle-from-property-intersections}.
\end{proof}

\section{Lexical KAT Action Monoids and Power Submonoid Spectra}

We begin by recording basic properties of lexical KAT action monoids
$\mathcal{M}$ and their power submonoid spectra
$\mathsf{Spec}_{\mathrm{KAT}}(\mathcal{M})$ as introduced in
Definition~\ref{def:lex-kat-action-monoid} and
Definition~\ref{def:power-submonoid-spectrum}.

\begin{proposition}[Submonoid closure and KAT structure]
\label{prop:submonoid-closure-kat}
Let $\mathcal{M}$ be a lexical KAT action monoid on the free monoid
$(\Sigma^*,\circ,\varepsilon)$ as in
Definition~\ref{def:lex-kat-action-monoid}, and let $N\subseteq\mathcal{M}$
be a subset closed under composition and containing the identity
$\mathrm{id}_{\Sigma^*}$. Then:
\begin{enumerate}
  \item $(N,\circ,\mathrm{id}_{\Sigma^*})$ is a submonoid of
        $(\mathcal{M},\circ,\mathrm{id}_{\Sigma^*})$.
  \item The KAT interpretation of $\mathcal{M}$ restricts to $N$ in
        the sense that the image of $N$ in the underlying KAT model
        $(K,T)$ is closed under the KAT operations and tests induced
        by $\mathcal{M}$.
\end{enumerate}
\end{proposition}

\begin{proof}
(1) Closure under composition and inclusion of the unit are exactly
the requirements for being a submonoid: for all $u,v\in N$, we have
$v\circ u\in N$, and $\mathrm{id}_{\Sigma^*}\in N$ by assumption.

(2) By Definition~\ref{def:lex-kat-action-monoid}, $\mathcal{M}$ comes
equipped with a map $u\mapsto [u]\in K$ into a KAT model $(K,T)$ such
that composition in $\mathcal{M}$ is compatible with multiplication in
$K$ and such that tests are mapped into $T\subseteq K$. The image of
$N$ in $K$ is then closed under multiplication because $N$ is closed
under composition and $[\cdot]$ is a monoid homomorphism. Tests
associated with elements of $N$ lie in $T$ by construction. Hence the
restrictions of the KAT operations to the image of $N$ again define a
KAT structure, now generated by $N$ rather than $\mathcal{M}$. This is
exactly the sense in which the KAT structure restricts to $N$.
\end{proof}

\begin{proposition}[Intersection of submonoids and spectra]
\label{prop:intersection-submonoids-spectra}
Let $\mathcal{M}$ be a monoid and let $\{N_a\}_{a\in A}$ be a finite
family of submonoids of $\mathcal{M}$. Then $\bigcap_{a\in A} N_a$ is
again a submonoid of~$\mathcal{M}$. In particular, if
$\mathsf{Spec}_{\mathrm{KAT}}(\mathcal{M})$ is a family of submonoids
closed under finite intersections, as required by
Definition~\ref{def:power-submonoid-spectrum}, then for any finite
subfamily $\{N_1,\dots,N_k\}\subseteq\mathsf{Spec}_{\mathrm{KAT}}(\mathcal{M})$
the intersection $N_1\cap\cdots\cap N_k$ is again in
$\mathsf{Spec}_{\mathrm{KAT}}(\mathcal{M})$.
\end{proposition}

\begin{proof}
Let $N:=\bigcap_{a\in A}N_a$. Since each $N_a$ is a submonoid, it
contains the identity $\mathrm{id}$, so $\mathrm{id}\in N$. If $u,v\in
N$, then $u,v\in N_a$ for all $a\in A$, so $v\circ u\in N_a$ for all
$a$; hence $v\circ u\in\bigcap_{a\in A}N_a=N$. Thus $N$ is closed
under composition and contains the identity, hence is a submonoid. The
statement about $\mathsf{Spec}_{\mathrm{KAT}}(\mathcal{M})$ is then
just the explicit closure requirement in
Definition~\ref{def:power-submonoid-spectrum}.
\end{proof}

\begin{remark}[Functional roles and stability]
\label{rem:functional-roles-stability}
In the main text, each member $N\in\mathsf{Spec}_{\mathrm{KAT}}(\mathcal{M})$
is assigned a functional role (historical topology, predictive
topology, halting closure, compliance pipeline, etc.). The above
propositions guarantee that intersections of such submonoids are again
submonoids and inherit well-defined functional roles (intersections of
roles). This justifies operations such as restricting attention to
``historical \emph{and} compliance'' submonoids when defining PACER
variants for specific governance tasks.
\end{remark}

\section{Homotopy Relations on Reading Paths}

We next show that the homotopy relation used for PACER reading paths
is an equivalence relation and interacts appropriately with
composition, providing a robust notion of path equivalence for
certificate-based AI.

\begin{proposition}[Homotopy of PACER reading paths is an equivalence relation]
\label{prop:pacer-homotopy-equivalence}
Let $\sim$ be the homotopy relation on PACER reading paths defined in
Definition~\ref{def:pacer-reading-path-homotopy}. Then $\sim$ is an
equivalence relation: it is reflexive, symmetric, and transitive.
\end{proposition}

\begin{proof}
By Definition~\ref{def:pacer-reading-path-homotopy}, two reading paths
are homotopy equivalent if they can be related by a finite sequence of
\emph{elementary homotopy moves}, each supported on a small subdiagram
(e.g.\ $v_i\to v_j\to v_k$ or $v_i\to v_j\leftarrow v_k$) and
satisfying certain commutativity and safety conditions.

Reflexivity: For any path $\gamma$, the empty sequence of homotopy
moves relates $\gamma$ to itself, so $\gamma\sim\gamma$.

Symmetry: If $\gamma_1$ is related to $\gamma_2$ by a finite sequence
of elementary moves, then we can invert each move (swapping the roles
of the two local paths in the corresponding 2-cell). Since the local
conditions (commutativity and safety) are symmetric under this swap,
the inverse sequence of moves relates $\gamma_2$ back to $\gamma_1$,
so $\gamma_2\sim\gamma_1$.

Transitivity: If $\gamma_1\sim\gamma_2$ and $\gamma_2\sim\gamma_3$,
then there exists a finite sequence of moves from $\gamma_1$ to
$\gamma_2$ and another from $\gamma_2$ to $\gamma_3$. Concatenating
these sequences, we obtain a finite sequence of moves relating
$\gamma_1$ to $\gamma_3$. The local conditions remain satisfied since
each individual move in the concatenated sequence satisfies them.
Hence $\gamma_1\sim\gamma_3$. This proves that $\sim$ is an
equivalence relation.
\end{proof}

\begin{proposition}[Homotopy and concatenation]
\label{prop:homotopy-concatenation}
Let $\gamma_1\sim\gamma_1'$ and $\gamma_2\sim\gamma_2'$ be homotopy-
equivalent reading paths such that the end node of $\gamma_1$ equals
the start node of $\gamma_2$, and similarly for $\gamma_1'$ and
$\gamma_2'$. Then the concatenated paths satisfy
\[
  \gamma_1\cdot\gamma_2 \;\sim\; \gamma_1'\cdot\gamma_2'.
\]
\end{proposition}

\begin{proof}
The path $\gamma_1$ can be transformed into $\gamma_1'$ by a finite
sequence of elementary homotopy moves, and similarly for $\gamma_2$
and $\gamma_2'$. Applying the moves that transform $\gamma_1$ into
$\gamma_1'$ to the initial segment of the concatenated path
$\gamma_1\cdot\gamma_2$ yields $\gamma_1'\cdot\gamma_2$; then applying
the moves that transform $\gamma_2$ into $\gamma_2'$ to the terminal
segment yields $\gamma_1'\cdot\gamma_2'$. Each move is local and
respects the boundary nodes, and the local commutativity and safety
conditions are satisfied by assumption. Hence the concatenated paths
are homotopy-equivalent.
\end{proof}

\begin{remark}[Homotopy-complete books and path algebra]
\label{rem:homotopy-complete-book-path-algebra}
Proposition~\ref{prop:pacer-homotopy-equivalence} and
Proposition~\ref{prop:homotopy-concatenation} endow the set of
reading-path homotopy classes with a natural concatenation operation:
concatenation is associative up to equality of classes and has
identity elements given by classes of trivial paths at a node. In a
homotopy-complete book (Definition~\ref{def:homotopy-complete-book}),
this yields a well-structured path algebra suitable for governance
analyses; the homotopy completeness conditions ensure that relevant
deviations between PACER runs can be captured as controlled
homotopies rather than pathological divergences.
\end{remark}

\section{Operator Bundles and Certificate Bundles}

We now give more detailed arguments behind the existence of operator
bundles and certificate bundles under homotopy admissibility
assumptions, expanding on
Proposition~\ref{prop:gaugeable-operator-bundle-from-property-intersections}
and Definition~\ref{def:certificate-bundle}. The results are standard
in spirit (gluing local trivializations into principal bundles), but
we state them here in the specific setting of HACA/PACER.

\begin{proposition}[Existence of operator bundles over property intersections]
\label{prop:operator-bundle-existence-proof}
Let $\mathcal{L}_{\mathrm{HACA}}$ be the category of HACA law-objects
with heterogeneous property-changing morphisms $(f_{ij},S_{ij})$ as in
Definition~\ref{def:property-changing-operators}. For a finite chain
\[
  L_0\to L_1\to\cdots\to L_n
\]
of law-objects, let $I(L_0,\dots,L_n)$ and $\mathcal{U}$ be as in
Proposition~\ref{prop:gaugeable-operator-bundle-from-property-intersections},
and assume that conditions~(1)--(3) there hold. Then there exists a
principal $G_{\mathrm{op}}$-bundle
\[
  \pi_{\mathrm{op}}:P_{\mathrm{op}}\to\mathcal{U}
\]
whose fibers coincide (locally) with the admissible operator fibers
$\mathcal{F}_p$.
\end{proposition}

\begin{proof}
We provide a detailed proof of Proposition~\ref{prop:operator-bundle-existence-proof}, making explicit the bundle gluing construction underlying the statement.

By condition~(1), the intersection sets $I(L_0,\dots,L_n)$ are non-
empty for homotopy-admissible chains and vary in a locally stable
manner under perturbations of the chains. This endows $\mathcal{U}$
with a topology such that each $I(L_0,\dots,L_n)$ is an open subset of
$\mathcal{U}$ (or, more generally, belongs to a basis of open sets).

For each $p\in\mathcal{U}$, let $\mathcal{F}_p$ be the set of
admissible operators at $p$ as in condition~(2). By assumption,
$G_{\mathrm{op}}$ acts freely and transitively on each non-empty
$\mathcal{F}_p$ (at least locally), so for $p$ in a small neighborhood
$U\subseteq\mathcal{U}$ we can identify $\mathcal{F}_p$ with
$G_{\mathrm{op}}$ via a choice of local gauge
\[
  \phi_U:\mathcal{F}_p\to G_{\mathrm{op}}.
\]
These local gauges provide local trivializations
\[
  \Phi_U:\pi_{\mathrm{op}}^{-1}(U)\to U\times G_{\mathrm{op}},
\]
where $\pi_{\mathrm{op}}^{-1}(U)$ is defined as
\[
  \bigsqcup_{p\in U}\mathcal{F}_p.
\]

On overlaps $U\cap V$, the two local trivializations $\Phi_U$ and
$\Phi_V$ differ by a transition function
\[
  g_{UV}:U\cap V\to G_{\mathrm{op}},
\]
defined by
\[
  \Phi_V\circ \Phi_U^{-1}(p,g) = (p,g_{UV}(p)\cdot g).
\]
Condition~(3) ensures that the higher homotopy data trivialize the
Jacobiator-gap cohomology classes over $\mathcal{U}$, which implies
that the transition functions satisfy the usual cocycle condition on
triple overlaps:
\[
  g_{UV}(p)\,g_{VW}(p)\,g_{WU}(p)=e
  \quad\text{for all }p\in U\cap V\cap W.
\]
This is exactly the condition required to glue the local trivial
bundles $U\times G_{\mathrm{op}}$ into a global principal
$G_{\mathrm{op}}$-bundle over $\mathcal{U}$. Thus
$\pi_{\mathrm{op}}:P_{\mathrm{op}}\to\mathcal{U}$ exists with fibers
locally identified with $\mathcal{F}_p$, as claimed.
\end{proof}

\begin{proposition}[Basic properties of certificate bundles]
\label{prop:certificate-bundle-properties}
Let $L$ be a HACA law-object, and let
$\pi_{\mathrm{cert}}:P_{\mathrm{cert}}\to B_{\mathrm{cert}}(L)$ be the
certificate bundle defined in
Definition~\ref{def:certificate-bundle}. Then:
\begin{enumerate}
  \item For each $b\in B_{\mathrm{cert}}(L)$, the fiber
        $\mathcal{H}_b$ is non-empty if and only if the corresponding
        chain of law-objects and properties is homotopy admissible in
        the sense of Definition~\ref{def:homotopy-admissible-jgap}.
  \item Restrictions of $P_{\mathrm{cert}}$ to subsets of
        $B_{\mathrm{cert}}(L)$ corresponding to Level~A/B/C completeness
        regimes admit local or global trivializations consistent with
        the LA/LB/LC interpretation in the main text.
\end{enumerate}
\end{proposition}

\begin{proof}
(1) If the chain is homotopy admissible, then by definition there exist
higher-homotopy and smoothing data $\{H_i\}$ filling the relevant
Jacobiator-gaps and stabilizing property intersections and operator
fibers; hence $\mathcal{H}_b$ is non-empty. Conversely, if
$\mathcal{H}_b$ is non-empty, then any $h\in\mathcal{H}_b$ furnishes a
choice of homotopy-filling data that trivializes the Jacobiator-gap
classes and satisfies the admissibility conditions, so the chain is
homotopy admissible.

(2) For Level~A, we restrict $P_{\mathrm{cert}}$ to the image of a
single homotopy-complete path. This image is contractible, so any
principal bundle (or bundle with fibers $\mathcal{H}_b$) over it is
trivializable: we can pick a continuous section along the path by
following the homotopy data $H_i$. For Level~B, we restrict to a
neighborhood of a homotopy common core
$\mathcal{C}\subseteq B_{\mathrm{cert}}(L)$: the existence of a common
core ensures that the relevant fibers can be trivialized over
$\mathcal{C}$ and that this trivialization can be extended to a
neighborhood using the homotopy data. For Level~C, we restrict to a
normative subspace $\mathcal{N}$ with the property that all admissible
paths can be projected into $\mathcal{N}$; by design, the certificate
data over $\mathcal{N}$ are chosen so that
$P_{\mathrm{cert}}|_{\mathcal{N}}$ is globally trivializable up to a
chosen baseline. This matches the interpretations given in the main
text.
\end{proof}

\section{Governance-Equivalence of Reading Paths}

Finally, we establish basic properties of the governance-equivalence
relation $\approx_{\mathrm{gov}}$ on reading paths introduced in
Definition~\ref{def:governance-equivalent-reading-paths}.

\begin{proposition}[Governance-equivalence is an equivalence relation]
\label{prop:gov-equivalence-equivalence}
Let $\approx_{\mathrm{gov}}$ be the relation on PACER reading paths
defined in Definition~\ref{def:governance-equivalent-reading-paths}.
Then $\approx_{\mathrm{gov}}$ is an equivalence relation.
\end{proposition}

\begin{proof}
By Definition~\ref{def:governance-equivalent-reading-paths},
$\gamma^{(1)}\approx_{\mathrm{gov}}\gamma^{(2)}$ holds if:
\begin{itemize}
  \item $\gamma^{(1)}\sim\gamma^{(2)}$ in the homotopy sense of
        Definition~\ref{def:pacer-reading-path-homotopy}; and
  \item their certificate records agree up to specified tolerances
        in a manner fixed by governance choices.
\end{itemize}

Reflexivity: For any path $\gamma$, we have $\gamma\sim\gamma$ by
Proposition~\ref{prop:pacer-homotopy-equivalence}, and its certificate
record $\mathsf{Cert}(\gamma)$ agrees with itself. Thus
$\gamma\approx_{\mathrm{gov}}\gamma$.

Symmetry: If $\gamma^{(1)}\approx_{\mathrm{gov}}\gamma^{(2)}$, then
$\gamma^{(1)}\sim\gamma^{(2)}$, and the certificate records for
$\gamma^{(1)}$ and $\gamma^{(2)}$ agree up to tolerances. Since
$\sim$ is symmetric, $\gamma^{(2)}\sim\gamma^{(1)}$, and the tolerance
condition is symmetric by design (differences bounded by specified
norms or metrics are symmetric). Hence
$\gamma^{(2)}\approx_{\mathrm{gov}}\gamma^{(1)}$.

Transitivity: If $\gamma^{(1)}\approx_{\mathrm{gov}}\gamma^{(2)}$ and
$\gamma^{(2)}\approx_{\mathrm{gov}}\gamma^{(3)}$, then
$\gamma^{(1)}\sim\gamma^{(2)}$ and $\gamma^{(2)}\sim\gamma^{(3)}$, so
by transitivity of $\sim$ we have $\gamma^{(1)}\sim\gamma^{(3)}$. The
certificate record of $\gamma^{(1)}$ is within tolerances of that for
$\gamma^{(2)}$, and the latter is within tolerances of that for
$\gamma^{(3)}$. Assuming that the tolerance relation is defined via a
metric or seminorm satisfying the triangle inequality (which is a
natural governance requirement), it follows that the certificate
record of $\gamma^{(1)}$ is within tolerances of that for
$\gamma^{(3)}$. Thus
$\gamma^{(1)}\approx_{\mathrm{gov}}\gamma^{(3)}$.
\end{proof}

\begin{proposition}[Stability of governance classes under small updates]
\label{prop:gov-class-stability}
Let $\gamma^{(1)}$ be a PACER reading path and suppose
$\gamma^{(2)}$ is obtained from $\gamma^{(1)}$ by a finite number of
local policy updates that:
\begin{enumerate}
  \item change only finitely many edges in the paragraph–monoid
        network, and only by elementary homotopy moves;
  \item modify certificate-relevant quantities (e.g.\ $\mu_i^\pi$,
        $\alpha_i^\pi$, $\mathbf{v}$, $\mathbf{q}^{\mathrm{sem}}(\pi)$)
        by amounts smaller than pre-specified tolerance margins.
\end{enumerate}
Then $\gamma^{(1)}\approx_{\mathrm{gov}}\gamma^{(2)}$.
\end{proposition}

\begin{proof}
Condition~(1) implies $\gamma^{(1)}\sim\gamma^{(2)}$ by the definition
of the homotopy relation $\sim$ as generated by elementary homotopy
moves. Condition~(2) ensures that the certificate records for
$\gamma^{(1)}$ and $\gamma^{(2)}$ differ by less than the tolerance
thresholds used to define $\approx_{\mathrm{gov}}$. Hence the two
paths are governance-equivalent.
\end{proof}

\begin{remark}[Interpretation for policy updates and system evolution]
\label{rem:policy-update-interpretation}
Proposition~\ref{prop:gov-class-stability} formalizes the intuitive
idea that small, well-controlled updates to a PACER policy should not
change the governance classification of reading paths: as long as
changes can be implemented as local homotopies in the paragraph–monoid
network and do not significantly alter certificate-relevant statistics,
the governance reading-path class $[\gamma]_{\mathrm{gov}}$ remains
unchanged. This provides a mathematically precise notion of stability
for governance-oriented behavior under incremental system evolution.
\end{remark}



\chapter{Notation and Cross-Volume Index}


This chapter summarizes the main notation used throughout this volume
and provides a cross-volume index linking key symbols to their roles
in the pure mathematics volume and in Applications Volumes~I and~II.
The list is not exhaustive, but it covers the principal symbols
introduced or heavily used in the HACA/PACER context, with an
emphasis on law-space geometry, KAT action monoids, GRL path integrals,
and certificate structures.

\section*{General Conventions}

Unless otherwise stated, the following general conventions are in force:

\begin{description}
  \item[$X\subseteq Y$] $X$ is a subset of $Y$.
  \item[$\mathcal{P}(X)$] The power set of $X$.
  \item[$\bigsqcup$] Disjoint union of sets.
  \item[$\sqcup$] Disjoint union when the context is clear.
  \item[$\mathcal{C}_L$] A distinguished family of semantic certificate
        predicates for a law-object $L$.
  \item[$\mathbb{N},\mathbb{Z},\mathbb{R}$] The sets of natural numbers,
        integers, and real numbers, respectively.
  \item[$\mathcal{O}(\cdot)$] Big-$\mathcal{O}$ notation for asymptotic
        bounds when needed.
\end{description}

\section*{Law-Space and PFB-GNLA Notation}

The core law-space and generalized Lie-algebraic structures are
inherited from the pure mathematics volume and Applications Volume~I.

\begin{description}
  \item[$\GZLS$] GaoZheng law-space (global law-space manifold).
  \item[$\GZLC$] Law-space connection on a principal bundle over
        $\GZLS$.
  \item[$\GZLCurv$] Law-space curvature associated with $\GZLC$.
  \item[$\GZTLC$] GaoZheng triple-layer connection (three coupled
        connections on law-space bundles, specialized in this volume to
        symbolic/semantic/decision layers of HACA).
  \item[$\GAlgebra$] GaoZheng generalized noncommutative Lie algebra
        (PFB-GNLA) bundle over $\GZLS$.
  \item[$\GZGMS$] GaoZheng generalized mathematical structure (GZ-GMS)
        as introduced in the pure mathematics volume.
  \item[$B_L\subseteq\GZLS$] Law-space sector supporting a given
        law-object $L$.
  \item[\textbf{PFB-GNLA}] \emph{Principal-Bundle–Based Generalized
        Noncommutative Lie Algebra}, the global algebraic-geometric
        framework underlying all three volumes.
\end{description}

\section*{HACA Law-Objects and Semantic Structures}

\begin{description}
  \item[$\mathcal{L}_{\mathrm{HACA}}$] Category of HACA law-objects
        (Definition~\ref{def:haca-law-objects-category}).
  \item[$L$] A HACA law-object.
  \item[$L^{(\mathrm{sym})}$] Symbolic layer of $L$ (lexical/KAT level).
  \item[$L^{(\mathrm{sem})}$] Semantic layer of $L$ (semantic
        structure space level).
  \item[$L^{(\mathrm{dec})}$] Decision layer of $L$ (policy/decision
        layer).
  \item[$\Omega^{(0)}(L)$] Character- or token-level operator
        configurations.
  \item[$\Omega^{(1)}(L)$] Word- or base-pack–level operators.
  \item[$\Omega^{(2)}(L)$] Sentence-level operator configurations.
  \item[$\Omega^{(3)}(L)$] Paragraph- or reading-path operators (KAT
        path integrals).
  \item[$\Omega(L)$] Multi-scale operator universe of $L$,
        \(
          \Omega(L)=\Omega^{(0)}(L)\sqcup\Omega^{(1)}(L)
          \sqcup\Omega^{(2)}(L)\sqcup\Omega^{(3)}(L)
        \).
  \item[$\mathcal{S}(L)$] Semantic structure space of $L$.
  \item[$\operatorname{Sem}_L$] Semantic map
        $\operatorname{Sem}_L:\Omega(L)\to\mathcal{S}(L)$.
  \item[$\operatorname{Law}_L$] Law-space projection
        $\operatorname{Law}_L:\mathcal{S}(L)\to B_L\subseteq\GZLS$.
  \item[$\Phi_L$] Composite law-space footprint map
        $\Phi_L:=\operatorname{Law}_L\circ\operatorname{Sem}_L:
         \Omega(L)\to B_L$.
\end{description}

\section*{Property Layers, Semantic Functors, and Jacobiator-Gaps}

\begin{description}
  \item[$\mathrm{Prop}(L)$] Property layer (property set) on a HACA
        law-object $L$; $\mathrm{Prop}:\mathcal{L}_{\mathrm{HACA}}\to
        \mathbf{Set}$ is the property-layer functor
        (Definition~\ref{def:haca-property-layer}).
  \item[$p\in\mathrm{Prop}(L)$] A property instance of $L$.
  \item[$S_{ij}$] Property-changing operator $S_{ij}:\mathrm{Prop}(L_i)
        \to\mathrm{Prop}(L_j)$ between HACA law-objects
        (Definition~\ref{def:property-changing-operators}).
  \item[$(f_{ij},S_{ij})$] Heterogeneous property-changing morphism
        from $L_i$ to $L_j$ combining a law-space morphism $f_{ij}$ and
        a property-changing operator $S_{ij}$
        (Definition~\ref{def:property-changing-operators}).
  \item[$J_{\mathrm{prop}}$] Property-layer Jacobiator-gap
        \(
          J_{\mathrm{prop}}:=S_{23}\circ S_{12}-S_{13}
        \)
        for a triple $(L_1,L_2,L_3)$
        (Definition~\ref{def:haca-jacobiator-gap-property-layer}).
  \item[$\mathsf{Obs}(L)$] Family of observable functionals on
        $\Omega(L)$ (Definition~\ref{def:haca-observable-properties}).
  \item[$\chi$] Truth-valued semantic predicate
        $\chi:\mathcal{S}(L)\to\{\mathrm{true},\mathrm{false}\}$
        (Definition~\ref{def:haca-semantic-predicates}).
  \item[$P_\chi$] Semantic property class associated with $\chi$:
        $P_\chi=\{\omega\in\Omega(L)\mid
        \chi(\operatorname{Sem}_L(\omega))=\mathrm{true}\}$.
  \item[$\mathcal{L}_{\mathrm{obs}}(L)$] Collection of observable
        subsets of $\Omega(L)$ whose law-space pullbacks are
        geometrically realizable and homotopy-stable
        (Definition~\ref{def:haca-observable-properties}).
  \item[$\mathfrak{F}_L$] Semantic functor algebra on $L$; a space of
        semantic functors $F:\mathcal{S}(L)\to V$ equipped with a
        noncommutative bracket and Jacobiator-gap data
        (Definition~\ref{def:noncommutative-semantic-functor-algebra}).
\end{description}

\section*{Lexical KAT Action Monoids and PACER Notation}

\begin{description}
  \item[$\Sigma$] Lexical alphabet (characters, subword tokens, etc.).
  \item[$\Sigma^*$] Free monoid of lexical sequences
        $(\Sigma^*,\circ,\varepsilon)$
        (Definition~\ref{def:free-monoid-lexical}).
  \item[$\mathrm{End}(\Sigma^*)$] Endomorphism monoid of $\Sigma^*$,
        i.e.\ all functions $f:\Sigma^*\to\Sigma^*$.
  \item[$\mathcal{M}$] Lexical KAT action monoid on $\Sigma^*$
        (Definition~\ref{def:lex-kat-action-monoid}).
  \item[$K,T$] Underlying KAT model $(K,T)$: $K$ is a Kleene algebra
        and $T\subseteq K$ is the set of tests.
  \item[$\mathcal{B}^{(0)}(L)$] Token-level base packs for $L$
        (Definition~\ref{def:base-packs}).
  \item[$\mathcal{B}^{(1)}(L)$] Phrase-level base packs for $L$.
  \item[$\mathcal{B}^{(2)}(L)$] Sentence-level base packs for $L$.
  \item[$\mathcal{B}(L)$] Union of all base packs for $L$.
  \item[$\mathsf{Spec}_{\mathrm{KAT}}(\mathcal{M})$] Power submonoid
        spectrum of $\mathcal{M}$, i.e.\ a family of submonoids
        $N\subseteq\mathcal{M}$ with assigned functional roles (memory,
        prediction, halting, compliance, etc.)
        (Definition~\ref{def:power-submonoid-spectrum}).
  \item[$\mathcal{G}_t$] Local generator set at context step $t$
        (Definition~\ref{def:local-context-submonoid}).
  \item[$M_t=\langle\mathcal{G}_t\rangle$] Local context submonoid at
        step~$t$, generated by $\mathcal{G}_t$
        (Definition~\ref{def:local-context-submonoid}).
  \item[$P_i$] Paragraph~$i$ of a document.
  \item[$M(P_i)$] Paragraph-level submonoid associated with paragraph
        $P_i$ (Definition~\ref{def:paragraph-submonoids}).
  \item[$v_i=(P_i,M(P_i))$] Paragraph node in the paragraph–monoid
        network (Definition~\ref{def:paragraph-submonoids}).
  \item[$F_i$] Paragraph-level KAT monoid morphism
        $F_i:M(P_i)\to M(P_{i+1})$
        (Definition~\ref{def:paragraph-monoid-morphism}).
  \item[$\gamma^{(\ell)}$] PACER reading path, i.e.\ a sequence of
        paragraph nodes and morphisms indexed by a run or policy
        $\ell$ (Definition~\ref{def:pacer-reading-paths}).
\end{description}

\section*{GRL, MDQ, Occupancy Measures, and Value Geometry}

\begin{description}
  \item[\textbf{GRL}] \emph{GaoZheng Reinforcement Learning (GRL) in
        law-space}, the path-integral–based reinforcement learning
        framework introduced in Applications Volume~I.
  \item[\textbf{MDQ}] \emph{Differential dynamical quanta}, local
        quanta of law-aware dynamics along GRL path integrals, used in
        Applications Volume~I and extended in this volume to semantic
        MDQ.
  \item[$w\in\mathsf{Paths}(\mathcal{M})$] Lexical operator path on the
        KAT action monoid $\mathcal{M}$
        (Definition~\ref{def:operator-paths-on-M}).
  \item[$S_{\mathrm{lex}}$] Lexical GRL action functional on operator
        paths $w$ (Definition~\ref{def:lexical-grl-action-on-M}).
  \item[$\mathcal{Z}_{\mathrm{lex}}$] Lexical GRL path integral on
        operator paths $w$ (Definition~\ref{def:lexical-grl-action-on-M}).
  \item[$\mathcal{Z}^{\mathrm{sem}}_\chi$] Semantic GRL path integral
        associated with predicate $\chi$ and paragraph path
        $\gamma\in\Omega^{(3)}(L)$
        (Definition~\ref{def:semantic-grl-path-integral}).
  \item[$\Xi_\chi(\gamma)$] Truth-valued semantic GRL path-integral
        predicate induced by $\mathcal{Z}^{\mathrm{sem}}_\chi[\gamma]$
        (Definition~\ref{def:semantic-grl-path-integral}).
  \item[$\Omega^{(3)}_{\mathrm{adm}}(L)$] Set of semantically
        admissible paragraph paths for law-object $L$, i.e.\ those for
        which $\Xi_\chi(\gamma)=\mathrm{true}$ for all
        $\chi\in\mathcal{C}_L$
        (Definition~\ref{def:semantically-admissible-paragraph-operators}).
  \item[$\mu_i^\pi$] Occupancy measure of segment $i$ under policy
        $\pi$ (Definition~\ref{def:pack-occupancy-measures}).
  \item[$\mu^\pi=(\mu_i^\pi)_{i\in I_L}$] Occupancy measure vector for
        policy $\pi$.
  \item[$\mathbf{v}=(v_i)_{i\in I_L}$] Gradient-based value baseline
        vector (value geometry vector), with
        $v_i=\partial\widehat{\mathcal{J}}/\partial\mu_i$
        (Definitions~\ref{def:pack-value-geometry},
         \ref{def:gradient-based-baseline}).
  \item[$\alpha_i^\pi$] Strategy weight for segment $i$ under policy
        $\pi$ (Definition~\ref{def:mu-alpha-structural-coordinates}).
  \item[$\alpha^\pi=(\alpha_i^\pi)_{i\in I_L}$] Strategy weight vector
        for policy $\pi$.
  \item[$\mathbf{q}^{\mathrm{sem}}(\pi)$] Semantic MDQ profile for
        policy $\pi$, with components
        $q_i^{\mathrm{sem}}(\pi)$ describing semantic MDQ at segment
        $i$ (Definition~\ref{def:semantic-mdq-profile}).
  \item[$\mathcal{J}(\pi)$] SAC-PACER-HACA joint objective functional
        for policy $\pi$, combining lexical and semantic GRL actions,
        MDQ costs, value geometry, and entropy
        (SAC-PACER-HACA chapter).
\end{description}

\section*{Jacobiator-Gap, Homotopy Levels, and Certificate Bundles}

\begin{description}
  \item[$J_{\mathrm{prop}}$] Property-layer Jacobiator-gap
        $S_{23}\circ S_{12}-S_{13}$
        (Definition~\ref{def:haca-jacobiator-gap-property-layer}).
  \item[$H$] Homotopy data on the property layer used to fill
        Jacobiator-gaps (Definition~\ref{def:haca-jacobiator-gap-property-layer}).
  \item[$I(L_0,\dots,L_n)$] Property-layer intersection set for a
        chain $L_0\to\cdots\to L_n$ (Proposition~\ref{prop:gaugeable-operator-bundle-from-property-intersections}).
  \item[$\mathcal{U}$] Union of property intersections over
        homotopy-admissible chains (Proposition~\ref{prop:gaugeable-operator-bundle-from-property-intersections}).
  \item[$\pi_{\mathrm{op}}:P_{\mathrm{op}}\to\mathcal{U}$] Principal
        operator bundle whose fibers consist of admissible operators at
        each property configuration
        (Proposition~\ref{prop:gaugeable-operator-bundle-from-property-intersections}).
  \item[LA/LB/LC] Homotopy completeness levels (Level~A: local
        completeness along paths; Level~B: common-core completeness;
        Level~C: normative completeness over a normative subspace)
        (Definition~\ref{def:homotopy-completeness-levels}).
  \item[$B_{\mathrm{cert}}(L)$] Certificate base space for law-object
        $L$, parametrizing chains of law-objects, paths, and
        property-layer data
        (Definition~\ref{def:certificate-base-space}).
  \item[$P_{\mathrm{cert}}$] Certificate bundle over
        $B_{\mathrm{cert}}(L)$ whose fiber $\mathcal{H}_b$ at $b$
        consists of admissible Jacobiator-gap homotopies and gauge
        data (Definition~\ref{def:certificate-bundle}).
  \item[$\mathcal{H}_b$] Fiber of $P_{\mathrm{cert}}$ at base point
        $b\in B_{\mathrm{cert}}(L)$, i.e.\ the set of certificates at
        $b$ (Definition~\ref{def:certificate-bundle}).
  \item[$\mathsf{Cert}(\gamma)$] Law-space certificate record for a
        reading path $\gamma$, constructed by the law-space certificate
        extraction pipeline
        (Definition~\ref{def:lawspace-certificate-extraction}).
  \item[$\gamma$-class] Homotopy class of reading path $\gamma$ under
        the relation $\sim$
        (Definitions~\ref{def:pacer-reading-path-homotopy},
         \ref{prop:pacer-homotopy-equivalence}).
  \item[$\gamma_{\mathrm{gov}}$-class] Governance reading-path class for
        $\gamma$ under the governance-equivalence relation
        $\approx_{\mathrm{gov}}$
        (Definition~\ref{def:governance-equivalent-reading-paths}).
\end{description}

\section*{Abbreviations and Acronyms}

For convenience, we collect here the main abbreviations and their
English expansions used across the three volumes:

\begin{description}
  \item[\textbf{HACA}] \emph{Hierarchical Algebraic Cognitive
        Architectures} (in this volume often realized as hierarchical
        algebraic cognitive architectures with explicit KAT and
        semantic layers).
  \item[\textbf{PACER}] \emph{Pack-Aligned Compressive–Expansion
        Reasoner}, the engine implementing \\compressive–alignment–
        expansion–recompression cycles on lexical KAT action monoids
        within HACA law-objects.
  \item[\textbf{GRL}] \emph{GaoZheng Reinforcement Learning in
        law-space}, the path-integral–based reinforcement learning
        framework introduced in Applications Volume~I.
  \item[\textbf{MDQ}] \emph{Differential Dynamical Quanta}, local
        quanta of law-aware dynamics along GRL path integrals (for
        both physical/lexical and semantic variants).
  \item[\textbf{PFB-GNLA}] \emph{Principal-Bundle–Based Generalized
        Noncommutative Lie Algebra}, the algebraic-geometric framework
        for law-space and generalized Lie brackets.
  \item[\textbf{LBOPB}] \emph{Life Bio-Operator Principal Bundle},
        life-science total-operator principal bundle used in
        Applications Volume~II.
  \item[\textbf{FGOPB}] \emph{Financial Generalized Operator Principal
        Bundle} (when used), for financial law-space modeling (typically
        in FGOPB-related project documents).
  \item[\textbf{PGOM}] \emph{Pharmacogenomics Monoids}, monoid
        structures encoding pharmacogenomics \\law-objects in
        Applications Volume~II.
  \item[\textbf{KAT}] \emph{Kleene Algebra with Tests}, an algebraic
        framework for reasoning about programs and control flow,
        applied here to lexical operators and paragraph paths.
  \item[\textbf{SAC}] \emph{Soft Actor–Critic}, a reinforcement
        learning algorithm whose objective is adapted in this volume
        to the SAC-PACER-HACA setting.
  \item[\textbf{HACA/PACER}] Shorthand for the combined HACA law-space
        architecture and PACER engine framework, with SAC-PACER-HACA
        referring to GRL/SAC-style learning over this architecture.
\end{description}

\section*{Cross-Volume Index}

Finally, we summarize where key concepts are introduced across the
three main volumes of the GaoZheng G-Framework series:

\begin{description}
  \item[Pure mathematics volume] Introduces PL-PI meta-mathematical
        theory, GaoZheng law-space $\GZLS$, PFB-GNLA structures
        $\GAlgebra$, triple-layer connection $\GZTLC$, CH-layering,
        and the general theory of Jacobiator-gaps and higher homotopy
        in law-space~\cite{GaoZheng2025GFramework}.
  \item[Applications Volume~I (GRL and Physics)] Introduces GaoZheng
        Reinforcement Learning (GRL) in law-space, differential
        dynamical quanta (MDQ), law-space measures and path integrals
        for physical systems, and the first instances of law-space
        minimal-action principles~\cite{GaoZhengAppVol1}.
  \item[Applications Volume~II (LBOPB)] Develops LBOPB (Life Bio-
        Operator Principal Bundle) and related structures (PK/PD,
        PGOM, life-science law-objects) and applies GRL and MDQ to
        life-science pipelines (e.g.\ pathology-to-reverse-design
        trajectories)~\cite{GaoZhengAppVol2}.
  \item[Applications Volume~III (HACA/PACER, this volume)] Introduces
        HACA \\law-objects, lexical KAT action monoids, PACER engines,
        semantic functor algebras, semantic GRL path integrals,
        certificate bundles, homotopy networks of packs, and the
        certificate-based AI and governance framework built on top of
        the previous volumes.
\end{description}

This cross-volume index is intended to help the reader navigate
between pure law-space geometry, GRL/MDQ-based applications, and the
HACA/PACER law-space semantics and governance structures developed in
the present volume.

\clearpage


\begin{thebibliography}{99}

\phantomsection
\addcontentsline{toc}{chapter}{\bibname}

\bibitem{GaoZheng2025GFramework}
Gao, Z. (2025).
\newblock Meta-Mathematical Theory based on Pan-Logic Analysis and
Pan-Iterative Analysis (GaoZheng G-Framework) and the
Principal-Bundle-Based Generalized Noncommutative Lie Algebra
(GaoZheng G-Algebra).
\newblock In \emph{Meta-Mathematical Theory based on Pan-Logic Analysis and
Pan-Iterative Analysis (GaoZheng G-Framework) and the
Principal-Bundle-Based Generalized Noncommutative Lie Algebra
(GaoZheng G-Algebra): An Integrated Construction (v1.0)}.
\newblock Zenodo.
\newblock \url{https://doi.org/10.5281/zenodo.17651584}.

\bibitem{GaoZhengAppVol1}
Gao, Z. (2025).
\newblock GaoZheng G-Framework and GaoZheng G-Algebra: Applied Mathematics Volume I — Law-Space Geometry, GRL, Quantum Computing, and Superconductivity.
\newblock In \emph{GaoZheng G-Framework and GaoZheng G-Algebra: Applied Mathematics Volume I — Law-Space Geometry, GRL, Quantum Computing, and Superconductivity (v1.2)}.
\newblock Zenodo.
\newblock \url{https://doi.org/10.5281/zenodo.17686133}.

\bibitem{GaoZhengAppVol2}
Gao, Z. (2025).
\newblock GaoZheng G-Framework and GaoZheng G-Algebra: Applied Mathematics Volume II — LBOPB, Life-Science Monoids, and Generative Precision Medicine.
\newblock In \emph{GaoZheng G-Framework and GaoZheng G-Algebra: Applied Mathematics Volume II — LBOPB, Life-Science Monoids, and Generative Precision Medicine (v1.1)}.
\newblock Zenodo.
\newblock \url{https://doi.org/10.5281/zenodo.17744672}.



\bibitem{KozenSmithKAT1996}
Kozen, D., \& Smith, F. (1996).
\newblock Kleene algebra with tests: Completeness and decidability.
\newblock In D.~van Dalen \& M.~Bezem (Eds.), \emph{Computer Science Logic (CSL 1994)}, Lecture Notes in Computer Science, vol.~1414, 244--259.
\newblock Springer.


\bibitem{Haarnoja2018SAC}
Haarnoja, T., Zhou, A., Abbeel, P., \& Levine, S. (2018).
\newblock Soft Actor-Critic: Off-Policy Maximum Entropy Deep Reinforcement Learning with a Stochastic Actor.
\newblock In \emph{Proceedings of the 35th International Conference on Machine Learning (ICML 2018)}, Proceedings of Machine Learning Research, vol.~80, 1861--1870.


\bibitem{Vaswani2017Attention}
Vaswani, A., Shazeer, N., Parmar, N., Uszkoreit, J., Jones, L., Gomez, A.~N., Kaiser, {\L}., \& Polosukhin, I. (2017).
\newblock Attention Is All You Need.
\newblock In \emph{Advances in Neural Information Processing Systems 30 (NeurIPS 2017)}, 5998--6008.


\end{thebibliography}

\bigskip
\noindent
\textit{License notice.} This arXiv version is licensed under
CC BY 4.0 (Creative Commons Attribution 4.0 International).

\end{document}
